\documentclass[a4paper,11pt,reqno]{amsart}
% This is the authoritative version. main.tex is an earlier version retained for reference.
\usepackage[utf8]{inputenc}
\usepackage{enumerate}
\usepackage{amsmath, amssymb,amsthm}
\usepackage{mathtools}
\usepackage{leftidx}
\usepackage{setspace}
\numberwithin{equation}{section}
%\usepackage{mathbbol}
\setcounter{tocdepth}{1} %2=show subsections, 1= show only sections
\usepackage[left=3cm, right=3cm]{geometry}
\usepackage{hyperref}
\usepackage{xcolor}
%\definecolor{name}{model}{color-spec}
\definecolor{my-black}{rgb}{0,0,0}
\definecolor{my-blue}{rgb}{0,0,0.8}
\definecolor{my-red}{rgb}{0.8,0,0} 
\definecolor{my-green}{rgb}{0,0.5,0}
\hypersetup{colorlinks,
    linkcolor	= {my-blue},
    citecolor	= {my-red},
    urlcolor		= {my-black}
}
\usepackage{graphicx}
\usepackage{tikz-cd}
\onehalfspacing
%Javier's suggestion to avoid overfull boxes
%\overfullrule=10cm better not to use because ArXiv might mess up things
%\usepackage{microtype}

\theoremstyle{plain} %italic

\newtheorem{lemma}{Lemma}[section]
\newtheorem{theorem}{Theorem}
\newtheorem{corollary}{Corollary}[section]
\newtheorem{proposition}{Proposition}[section]
\newtheorem{assumption}{Assumption}[section]
\newtheorem{problem}{Problem}[section]
\newtheorem{consequence}{Consequence}[section]
\newtheorem*{theorem*}{Theorem} %to avoid numbering


\theoremstyle{definition} %non-italic
\newtheorem{remark}{Remark}[section]

\newtheorem{definition}{Definition}[section]
\newtheorem{notation}{Notation}[section]
\newtheorem{conjecture}{Conjecture}[section]
\newtheorem{example}{Example}[section]
\newtheorem{claim}{Claim}[section]

\theoremstyle{remark}
\newtheorem*{remark-non}{Remark}

\DeclareMathOperator{\id}{id}
\DeclareMathOperator{\real}{Re}
\DeclareMathOperator{\imag}{Im}
\DeclareMathOperator{\image}{im}
\DeclareMathOperator{\supp}{supp}
\DeclareMathOperator{\vol}{vol}
\DeclareMathOperator{\inv}{inv}
\DeclareMathOperator{\Tr}{Tr}
\DeclareMathOperator{\ord}{ord}
\DeclareMathOperator{\Hom}{Hom}
\DeclareMathOperator{\Aut}{Aut}
\DeclareMathOperator{\End}{End}
\DeclareMathOperator{\Hol}{Hol}
\DeclareMathOperator{\GL}{GL}
\DeclareMathOperator{\SL}{SL}
\DeclareMathOperator{\PGL}{PGL}
\DeclareMathOperator{\PSL}{PSL}
\DeclareMathOperator{\Spur}{Sp}
\DeclareMathOperator{\PSO}{PSO}
\DeclareMathOperator{\SO}{SO}
\DeclareMathOperator{\Or}{O}
\DeclareMathOperator{\Res}{Res}
\DeclareMathOperator{\diam}{diam}
\DeclareMathOperator{\sgn}{sgn}
\DeclareMathOperator{\sinc}{sinc}
\DeclareMathOperator{\disc}{disc}


\newcommand{\R}{\mathbb{R}}
\renewcommand{\P}{\mathbb{P}}
\newcommand{\C}{\mathbb{C}}
\newcommand{\N}{\mathbb{N}}
\newcommand{\Q}{\mathbb{Q}}
\newcommand{\Z}{\mathbb{Z}}
\newcommand{\F}{\mathbb{F}}
\newcommand{\g}{\mathfrak{g}}
\renewcommand{\o}{\mathfrak{o}}
\renewcommand{\d}{\mathfrak{d}}
\renewcommand{\a}{\mathfrak{a}}
\renewcommand{\b}{\mathfrak{b}}
\newcommand{\p}{\mathfrak{p}}
\newcommand{\I}{\mathbf{I}}

\renewcommand{\sl}{\mathfrak{sl}}
\newcommand{\h}{\mathfrak{h}}
\renewcommand{\H}{\mathbb{H}}
\textwidth 16cm
\oddsidemargin  0cm \evensidemargin 0cm \parindent0.7em
\textheight 23cm
\allowdisplaybreaks


\title[Time-frequency localization]{Eigenfunctions, Free Boundaries, and Time-Frequency Localization}
\author{Jo{\~a}o P.G. Ramos}



\begin{document}
\maketitle 
\begin{abstract}
We study Gaussian time--frequency localization operators through their Bargmann--Fock representation, developing an inverse and free-boundary approach to their spectral theory. Our perturbative inverse theory shows that functions sufficiently close to the ground-state Hermite function can be realized as eigenfunctions of localization operators; as a byproduct we recover the Abreu--D\"orfler disk characterization in the simply connected setting, and we prove that the stability estimate of G\'omez--Guerra, Ramos and Tilli is sharp.

We then apply these inverse results to Faber--Krahn-type problems for Gaussian time--frequency concentration: local maximizers of the first localization eigenvalue under a measure constraint are necessarily disks, and global maximizers exist by a concentration--compactness argument in the Bargmann--Fock space, recovering the Nicola--Tilli theorem from this perspective.
\end{abstract}

\noindent\textbf{Keywords:} Gaussian STFT, localization operators, Bargmann--Fock space, inverse spectral problem, free boundary, quadrature domains, Faber--Krahn inequality, time-frequency concentration.

\medskip

\noindent\textbf{MSC2020:} 42B10, 42C15, 30H20, 35R35, 35R30, 47B35, 49Q10.


% ==============================================================
% Section 1: Introduction
% ==============================================================
\section{Introduction}

\subsection{Time--frequency localization and the Gaussian model}

The study of time--frequency concentration lies at the heart of modern harmonic analysis, with deep connections to signal processing, mathematical physics, and phase--space methods; see, for instance, \cite{Grochenig,Mallat,Folland}. A central object in this theory is the \emph{short-time Fourier transform} (STFT), defined for a function $f \in L^2(\mathbb R)$ and a window $\varphi \in L^2(\mathbb R)$ by
\[
V_\varphi f(x,\omega) = \int_{\mathbb R} f(t)\,\overline{\varphi(t-x)}\,e^{-2\pi i t \omega}\,dt.
\]
This transform encodes the joint time--frequency content of $f$, and its squared modulus $|V_\varphi f|^2$ provides a phase--space energy density. Understanding how this energy can be concentrated in regions of the time--frequency plane is a fundamental problem, with both theoretical and applied significance.

In applications ranging from signal recovery to uncertainty principles, one is often interested in maximizing the quantity
\[
\int_\Omega |V_\varphi f(x,\omega)|^2\,dx\,d\omega
\]
over functions $f$ of unit norm and over sets $\Omega$ of prescribed measure. This leads naturally to the study of \emph{time--frequency localization operators}, introduced by Daubechies in \cite{Daubechies} and further developed in \cite{Daubechies-Paul,Ramanathan-Topiwala-Weyl,Ramanathan-Topiwala-Spectrogram,Cordero-Grochenig,Boggiatto-Cordero-Grochenig,Bayer-Grochenig,Kisil-CrossToeplitz,Nicola-Tilli-2,Ramos-time-frequency}. For general windows, sharp concentration bounds are notoriously elusive, and Lieb's uncertainty principle \cite{Lieb} remains a fundamental benchmark; see also \cite{abreu2021donoho} for related large-sieve phenomena in modulation and polyanalytic Fock spaces. In the Gaussian setting, by contrast, these operators acquire a particularly rigid structure: via the Bargmann transform \cite{Bargmann}, they can be realized as Toeplitz-type integral operators on the Bargmann--Fock space, and the problem becomes one of complex analysis in weighted spaces; compare \cite{Berezin,Berger-Coburn-Symbol,Berger-Coburn,Berger-Coburn-Heat,Zhu,Grochenig,Folland}. In this picture, the STFT is intertwined with an entire function whose weighted modulus encodes the concentration properties of the signal, and geometric questions translate into analytic ones.

\subsection{Spectral and geometric interplay}

A central theme in this area is the interplay between the \emph{spectral properties} of localization operators and the \emph{geometry} of the underlying domains. In the Gaussian case, Nicola and Tilli proved the sharp Faber--Krahn inequality for the STFT, showing that disks maximize concentration among all measurable sets of a fixed measure \cite{Nicola-Tilli}; see also \cite{Nicola-Tilli-2,Bastianoni-Cordero-Nicola,Bastianoni-Teofanov} for complementary spectral and regularity results on localization operators. The stability theory for this variational problem was subsequently developed in \cite{Gomez-Guerra-Ramos-Tilli}, where quantitative control of both the first eigenfunction and the Fraenkel asymmetry is obtained from the spectral deficit. Related concentration questions for other transforms, including wavelet transforms, have been investigated in \cite{Ramos-Tilli}.

From a complementary perspective, one may pose inverse questions: to what extent does spectral information determine the geometry of the localization domain? Conversely, given a function, can one construct a domain for which it is an eigenfunction of the associated localization operator? The first systematic result in this direction is due to Abreu and D\"orfler \cite{Abreu-Doerfler}, who showed that if a Hermite function is an eigenfunction of a localization operator on a simply connected domain, then the domain must be a disk. This inverse viewpoint sits naturally within the broader circle of ideas surrounding weighted quadrature identities, free-boundary problems, and inverse potential theory: classical results of Aharonov, Schiffer, and Zalcman \cite{Aharonov-Schiffer-Zalcman} characterize disks via mean-value identities for harmonic functions; Karp \cite{Karp} and Cherednichenko \cite{Cherednichenko} address inverse problems for Newtonian potentials; Gustafsson and collaborators \cite{Gustafsson,Gustafsson-Shahgholian,Shahgholian} develop the theory of quadrature domains and obstacle problems; and Isakov's monograph \cite{Isakov} provides a perturbative framework for related inverse source problems. The inverse viewpoint is also tied to the special role of Gaussian decay and Hermite structures in the Bargmann model; see \cite{Kulikov-Oliveira-Ramos,Goncalves-Oliveira-Steinerberger,Thangavelu}.

\subsection{Eigenfunctions as geometric data}

The guiding principle of the present paper is that, for Gaussian localization operators, eigenfunctions should be regarded as \emph{geometric data}. More precisely, the eigenfunction should determine the domain in rigid situations, and the possibility of prescribing the eigenfunction should in turn have strong consequences for nearby variational problems. The three main threads of the paper are organized around this idea:
\begin{enumerate}
\item rigid eigenfunctions force rigid domains;
\item near the Gaussian, one can prescribe the eigenfunction and construct the corresponding localization domain;
\item once such an eigenfunction-to-domain mechanism is available, it can be used to analyze sharpness and extremizers in concentration inequalities.
\end{enumerate}

The present work contributes to this line of research by developing a unified inverse and free-boundary framework for Gaussian time--frequency localization operators. Our approach builds on, and extends, the rigidity result of Abreu and D\"orfler \cite{Abreu-Doerfler}, the variational theory of Nicola and Tilli \cite{Nicola-Tilli,Nicola-Tilli-2}, and the quantitative stability theorems of G\'omez, Guerra, Ramos, and Tilli \cite{Gomez-Guerra-Ramos-Tilli}. At the technical level, our arguments also draw on ideas from inverse potential problems and quadrature-domain theory, in the spirit of \cite{Isakov,Cherednichenko,Gustafsson,Gustafsson-Shahgholian,Shahgholian}. In this way, the paper places localization operators squarely within the broader free-boundary and inverse-potential framework suggested by the Gaussian Bargmann model.

A few words on the relation with classical quadrature-domain theory are in order. In a (classical) quadrature domain $D\subset \R^2$, integration of harmonic functions against Lebesgue measure on $D$ reduces to a finite linear combination of point evaluations and their derivatives; see Gustafsson--Shahgholian \cite{Gustafsson-Shahgholian}, the Gustafsson lectures \cite{Gustafsson}, and Shahgholian \cite{Shahgholian} for a survey of the obstacle-problem characterization, the moving-boundary dynamics, and the connection with inverse potential theory. The moment system for our problem, derived in Section~3, is structurally different in two ways. First, the moments are weighted by the Gaussian density $e^{-\pi|z|^2}$ rather than by Lebesgue measure, so the underlying functional identity is not a quadrature identity for harmonic functions but a weighted-mean identity for entire functions in the Bargmann--Fock space. Second, the eigenvalue $\lambda$ enters the moment system as an explicit multiplicative parameter, whereas classical quadrature domains have no analogous spectral parameter. The free boundary obtained in Section~3 is therefore a \emph{$\lambda$-dependent weighted quadrature domain} for an entire-function class, and it is the dependence on $\lambda$ together with the Gaussian weight that makes the inverse problem nontrivial even at the linearized level around the disk.

The inverse-spectral viewpoint we adopt here also fits naturally into the broader study of Toeplitz operators on Bergman and Fock spaces. Symbol-recovery and inverse-spectral problems for Toeplitz operators have been pursued, in slightly different directions, by Berger and Coburn \cite{Berger-Coburn-Symbol,Berger-Coburn,Berger-Coburn-Heat}, by Berezin \cite{Berezin}, by Zhu \cite{Zhu}, and more recently by Bastianoni, Cordero, Nicola, and Teofanov \cite{Bastianoni-Cordero-Nicola,Bastianoni-Teofanov}, Fulsche and Hagger \cite{Fulsche-Hagger}, and Kisil \cite{Kisil-CrossToeplitz}; see also \cite{abreu2021donoho} for related polyanalytic and large-sieve phenomena. In all of these the symbol typically belongs to a function class on phase space, and the inverse problem asks to read off properties of the symbol from the spectrum. By contrast, the inverse problem we solve here is geometric: the symbol is the indicator function $\chi_U$ of an unknown domain $U$, and the spectral datum is a single eigenfunction. The free-boundary moment system we derive can therefore be regarded as the appropriate Toeplitz-symbol inverse problem on the Bargmann--Fock space when the symbol class is restricted to indicators of domains, and the present paper appears to be the first to study it in this form.

Despite these advances, two fundamental questions have remained open. First, no general construction of localization domains associated with prescribed near-Gaussian functions has been available, beyond a handful of highly symmetric examples. Second, the role of eigenfunction rigidity in variational concentration problems has not been clarified, leaving the sharpness of recent stability inequalities and the structure of local maximizers out of reach. The results of the present paper address these questions in a unified manner.

\subsection{An inverse construction near the Gaussian}

Our first main result is the constructive half of this picture. It shows that, near the Gaussian, one can indeed prescribe the eigenfunction and realize it through a suitable localization domain.

\begin{theorem}\label{thm:perturb}
Let $\lambda\in (0,1)$ and $N\in\mathbb N$. There exists $\varepsilon_0(\lambda,N)>0$ such that if
\[
f_0 = h_0 + \sum_{k=1}^N a_k h_k
\]
satisfies $\max_k |a_k|<\varepsilon_0$, then there exists a domain $U_\lambda$ such that
\[
\mathcal L_{U_\lambda} f_0 = \lambda f_0 .
\]
\end{theorem}

This theorem provides, to the best of our knowledge, the first systematic construction of localization domains associated with nontrivial linear combinations of Hermite functions. In particular, it supplies the inverse counterpart to the Abreu--D\"orfler rigidity phenomenon: rather than deducing the domain from a rigid eigenfunction, we construct domains starting from prescribed near-Gaussian eigenfunctions. The dependence of $\varepsilon_0$ on $\lambda$ and $N$ is explicit in our argument, and reflects the size of the spectral gap of the disk localization operator on the relevant Hermite subspace.

The proof is a direct perturbation argument at the Daubechies disk. After
applying the coefficientwise conjugate differential operator, we reformulate
the problem as an infinite system of weighted moment equations for an analytic
radial graph over the disk. At the disk the linearized moment map is diagonal
in Fourier modes, with an explicit inverse. An implicit-function theorem in
analytic boundary spaces produces a nearby domain satisfying the symmetrized
identity exactly, and a final elementary ODE argument recovers the original
eigenvalue equation.

A noteworthy aspect of this construction is that it reaches well beyond monomial-type eigenfunctions, the only case previously accessible by direct symmetry arguments. Our theorem shows that a much broader class of near-Gaussian functions can be realized spectrally, opening the door to a richer inverse theory and revealing new connections with nonlinear potential theory and weighted quadrature domains in the spirit of \cite{Gustafsson-Shahgholian,Shahgholian}.

\subsection{A new proof of the Abreu--D\"orfler rigidity theorem}

The same free-boundary framework also yields the rigid half of the correspondence, namely a short recovery of the classical theorem of Abreu and D\"orfler.

More precisely, we provide two new proofs of the Abreu--D\"orfler rigidity result based on our methods: one through the local radiality mechanism for the logarithmic potential, and a second through a holomorphic auxiliary function attached to the associated Poisson problem.

\begin{theorem}\label{thm:eig-loc-U}
Let $U$ be a simply connected bounded domain. If a Hermite function is an eigenfunction of the localization operator associated with $U$, then $U$ is a disk.
\end{theorem}

The proof transports the eigenvalue equation to the Bargmann--Fock space, where Hermite functions correspond to monomials and the localization operator becomes an integral operator with reproducing-kernel structure. In this setting, the eigenvalue equation translates into a functional identity involving exponential generating functions. Differentiating appropriately and invoking a Paley--Wiener type argument, we reduce the problem to the vanishing of a Fourier transform along complex lines.

This allows us to apply a divisibility lemma of Ehrenpreis--Malgrange type, which produces a logarithmic potential solving a Poisson equation with source supported on the domain. The decisive insight is that the spectral assumption imposes strong constraints on this potential: it must exhibit a form of radial behavior on the boundary. In the simply connected setting, this constraint forces the domain to be a disk, recovering the Abreu--D\"orfler rigidity theorem from the same inverse/free-boundary perspective.

\subsection{Sharpness of the GGRT stability inequality}

We now turn to consequences of this inverse point of view. Our first application concerns the quantitative stability theorem of G\'omez, Guerra, Ramos, and Tilli for the Faber--Krahn inequality for the Gaussian STFT. In the formulation most relevant here, if $U\subset \R^2$ is a measurable set of finite positive measure, $\lambda_1(U)$ denotes the first eigenvalue of the associated localization operator, and $f_U$ is a corresponding $L^2$-normalized first eigenfunction, then Corollary~1.4 of \cite{Gomez-Guerra-Ramos-Tilli} gives the bounds
\[
\min_{z_0\in\C,\ |c|=1}\|f_U-c\phi_{z_0}\|_{L^2(\R)}
\leq
C e^{|U|/2}
\left(
1-\frac{\lambda_1(U)}{1-e^{-|U|}}
\right)^{1/2},
\]
and
\[
\mathcal A(U)
\leq
K(|U|)
\left(
1-\frac{\lambda_1(U)}{1-e^{-|U|}}
\right)^{1/2},
\]
where $\phi_{z_0}$ ranges over the Gaussian optimizers and $\mathcal A(U)$ denotes the Fraenkel asymmetry of $U$.

Our perturbative inverse theorem provides the missing input needed to test this stability estimate on genuine spectral data. Starting from a prescribed polynomial perturbation of the Gaussian, Theorem~\ref{thm:perturb} produces an exact eigenpair $(U_\varepsilon,f_\varepsilon)$ for the localization operator, with $U_\varepsilon$ a smooth perturbation of the disk and $f_\varepsilon$ a corresponding first eigenfunction. The constructive direction of the paper thus yields explicit near-extremizers, against which one can compare the size of the spectral deficit with the actual geometric displacement of the domain from the class of disks. The resulting one-parameter family shows that the square-root dependence in Corollary~1.4 of \cite{Gomez-Guerra-Ramos-Tilli} is the best possible.

\begin{theorem} \label{thm:GGRT-sharp}
The exponent $1/2$ in Corollary~1.4 of \cite{Gomez-Guerra-Ramos-Tilli} is optimal. More precisely, the powers of the deficit
\[
1-\frac{\lambda_1(U)}{1-e^{-|U|}}
\]
appearing in both the eigenfunction estimate and the Fraenkel-asymmetry estimate there cannot be replaced by $\beta$ for any $\beta>1/2$.
\end{theorem}

The argument rests on constructing a carefully chosen family of perturbations of the disk via the inverse theorem above. Starting from a polynomial perturbation of the ground-state Hermite function, we obtain a family of domains whose geometry deviates from that of a disk at first order. A detailed asymptotic analysis of the eigenvalue equation, including a first-variation computation by a transport formula, then identifies the leading Fourier modes of the boundary deformation.

This analysis shows that the deformation is governed by second-order harmonics, which cannot be eliminated by translations of the disk. Comparing the perturbed domains with arbitrary disks then yields lower bounds on their symmetric difference that match the upper bounds in the stability estimate of \cite{Gomez-Guerra-Ramos-Tilli}. The proof combines spectral perturbation theory, geometric measure estimates, and harmonic analysis on the circle, illustrating how inverse constructions can be used to probe the sharpness of variational inequalities.

\subsection{Local maximizers are disks}

We next connect this eigenfunction-to-domain philosophy to variational problems. Our fourth main result concerns the structure of local maximizers for Gaussian time--frequency concentration.

\begin{theorem}\label{thm:local-maximizers}
Any local maximizer of $\lambda_{1,\varphi}$ under a measure constraint is, up to translation, a disk.
\end{theorem}

The proof relies on combining variational arguments with the inverse spectral perspective developed earlier. We first show that any local maximizer must coincide with a superlevel set of the density associated with its first eigenfunction in the Bargmann--Fock space. This reduces the geometric problem to a functional one.

We then analyze the spectral properties of the localization operator at a maximizer, showing that the first eigenvalue is simple and that derivatives of the eigenfunction remain eigenfunctions. This strong rigidity property forces the maximizing eigenfunction into a Gaussian form. Once this is known, the corresponding domain is determined as well, and one concludes that the only admissible configuration is a disk. The argument thus shows that extremizers are governed by the same rigidity mechanism that underlies the inverse theory.

\subsection{A new proof of the Nicola--Tilli theorem}

Finally, we address the global variational problem and recover the Nicola--Tilli theorem from the same perspective.

\begin{theorem}[New proof of Nicola--Tilli]
Disks maximize Gaussian time--frequency concentration among all sets of fixed measure.
\end{theorem}

Our proof is based on a concentration--compactness argument in the Bargmann--Fock space. We consider maximizing sequences and analyze their possible limiting behavior, ruling out vanishing and dichotomy scenarios using the structure of the Fock space and the invariance properties of the STFT. This yields the existence of a maximizer. Related concentration--compactness methods have recently been used in nearby time--frequency concentration problems, notably by Nicola, Romero, and Trapasso for ambiguity-function type concentration functionals \cite{Nicola-Romero-Trapasso}; see also \cite{Marceca-Romero-Speckbacher,Samuelsen} for recent work on Fourier concentration operators, and \cite{Fulsche-Hagger} for closely related developments in the polyanalytic Fock-space setting.

We then combine this existence result with the rigidity of local maximizers established earlier. Since any global maximizer is, in particular, a local maximizer, its eigenfunction is rigid and therefore determines the maximizing domain. This provides a new conceptual route to the Nicola--Tilli theorem, avoiding rearrangement inequalities and instead relying on inverse and free-boundary methods. The argument highlights the power of treating eigenfunctions as geometric data.

\subsection{New results versus new proofs}

We pause to clarify the status of each main result, since the paper combines genuinely new theorems with new proofs of theorems that are already in the literature. Theorem~\ref{thm:perturb} (the perturbative inverse construction near the Gaussian), Theorem~\ref{thm:GGRT-sharp} (sharpness of the GGRT stability inequality), and Theorem~\ref{thm:local-maximizers} (local maximizers are disks) are \emph{new results}, with no precursor known to the author. By contrast, Theorem~\ref{thm:eig-loc-U} (the Abreu--D\"orfler rigidity theorem in the simply connected setting) and the global Faber--Krahn theorem of Section~7 (the Nicola--Tilli theorem) are \emph{new proofs of known results}: they are reproved here from the inverse and free-boundary perspective developed in this paper, and we make no claim of priority on the statements themselves. The interest of these new proofs is methodological: they place the two classical theorems within the same unified inverse/free-boundary framework that produces the new results above.

\subsection{Organization of the paper}

The paper is organized as follows. In Section~2 we review the Bargmann--Fock representation and establish several auxiliary results, including the key divisibility lemma. Section~3 proves Theorem~\ref{thm:perturb}; the self-contained disk inverse theorem is proved there as part of the construction. Section~4 explains how the same machinery recovers the classical theorem of Abreu and D\"orfler in the simply connected case (Theorem~\ref{thm:eig-loc-U}). Section~5 contains the sharpness result, Theorem~\ref{thm:GGRT-sharp}, while the remaining sections are devoted to variational applications, culminating in the new proof of the Nicola--Tilli theorem.

\medskip

In summary, the results of this paper establish a bridge between time--frequency analysis, inverse spectral theory, and free-boundary problems in potential theory. This perspective not only resolves several open problems, but also supplies a flexible and robust framework for studying a wide range of questions in phase--space analysis.

% ==============================================================
% Section 2: Preliminaries
% ==============================================================
\section{Preliminaries}

\subsection{Main technical inputs}\label{subsec:main-technical-inputs}

For the convenience of the reader, and to make the boundaries between
quoted material, adapted material, and original arguments transparent, we
collect here the principal external theorems used in this paper, together
with an indication of the form in which each is used.

\begin{itemize}
\item \emph{Isakov's perturbative free-boundary theorem.} In Section~3 we
use the implicit-function-theorem framework for inverse free-boundary and
inverse-source problems developed in Isakov \cite{Isakov}. The general
theorem is \emph{adapted with proof}: the underlying perturbative
implicit-function strategy is Isakov's, but the function spaces, the
linearized moment map at the disk, and its invertibility on the relevant
subspace are reformulated and reproved here in a form tailored to the
Gaussian Bargmann--Fock setting.

\item \emph{Cherednichenko's inverse Newtonian-potential result.} The
recovery of a domain from its exterior logarithmic potential, used as a
conceptual guide in Sections~3 and~4, is taken from Cherednichenko
\cite{Cherednichenko}; see also \cite{Karp,Aharonov-Schiffer-Zalcman} for
the closely related classical statements. We \emph{quote} these
results in the form needed and do not reprove them.

\item \emph{Bathtub principle.} The Lieb--Loss bathtub principle is used
in Sections~6 and~7 to identify maximizers of $\int_\Omega u_F\,dA$ under
a measure constraint with superlevel sets of $u_F$. We quote it
\emph{verbatim} from \cite{Lieb}; see also Lemma~\ref{lem:basic-properties-global}(i)
above.

\item \emph{Concentration--compactness profile decomposition.} In
Section~7 we use a profile-decomposition statement in $\mathcal F^2(\C)$
in the spirit of P.-L.~Lions's concentration--compactness method and of
its time--frequency adaptations by Nicola and Tilli \cite{Nicola-Tilli,Nicola-Tilli-2}
and by Nicola, Romero, and Trapasso \cite{Nicola-Romero-Trapasso}. The
form we use is \emph{adapted with proof}: the abstract concentration
mechanism is classical, while the precise profile decomposition for
maximizing sequences of $I_F(s)$ in the Fock space, including the
treatment of the lack of translation compactness, is reproved here.

\item \emph{Ehrenpreis--Malgrange divisibility.} The divisibility of
entire functions on $\C^2$ vanishing on a complex hypersurface, which
underlies Lemma~\ref{lemma:fourier-to-fboundary} above, is the classical
Ehrenpreis--Malgrange divisibility theorem. The version we need is
\emph{proved as a self-contained lemma} (Lemma~\ref{lemma:fourier-to-fboundary}),
using only the Paley--Wiener--Schwartz theorem and elementary
two-variable polynomial divisibility, in order to keep track of explicit
support and regularity information.

\item \emph{Paley--Wiener and Paley--Wiener--Schwartz theorems.} The
classical Paley--Wiener theorem (and its Schwartz extension to compactly
supported distributions) is used several times to translate vanishing of
a Fourier transform along complex lines into vanishing of a function on
the support side. We \emph{quote} these results from
\cite{Hormander,Folland} in the standard form.
\end{itemize}

The remainder of Section~2 is devoted to the analytic preliminaries on
the Bargmann--Fock space and the divisibility lemma proper.

\subsection{Properties of the Bargmann-Fock space} We briefly recall the connection between the Gaussian short-time Fourier transform
(STFT) and the Bargmann--Fock space, which will be used systematically in what
follows. Excellent references for the material in this subsection are
\cite{Bargmann,Folland,Grochenig,Zhu}; see also
\cite{Berezin,Berger-Coburn,Daubechies}.

\medskip

Let $\varphi(x) = 2^{1/4} e^{-\pi x^2}$ be the normalized Gaussian. For $f \in L^2(\mathbb{R})$, its Gaussian STFT is
\begin{equation}\label{eq:STFT}
V_\varphi f(x,\omega)
= \int_{\mathbb{R}} f(t)\, \varphi(x-t)\, e^{-2\pi i t \omega}\, dt.
\end{equation}
The Bargmann transform
$\mathcal{B} : L^2(\mathbb{R}) \to \mathcal{F}^2(\mathbb{C})$ is defined by
\begin{equation}\label{eq:Bargmann}
(\mathcal{B}f)(z)
= 2^{1/4} \int_{\mathbb{R}} f(t)\, e^{2\pi t z - \pi t^2 - \frac{\pi}{2} z^2}\, dt,
\qquad z \in \mathbb{C}.
\end{equation}
It is well known that $\mathcal{B}$ is a unitary isomorphism onto the Bargmann--Fock
space $\mathcal{F}^2(\mathbb{C})$, endowed with the norm
\[
\|F\|_{\mathcal{F}^2}^2
= \int_{\mathbb{C}} |F(z)|^2 e^{-\pi |z|^2} \, dA(z).
\]
A direct computation shows that, for $z = x + i\omega$,
\begin{equation}\label{eq:STFT-B}
V_\varphi f(x,-\omega)
= e^{\pi i x \omega}\, (\mathcal{B}f)(z)\, e^{-\frac{\pi}{2}|z|^2}.
\end{equation}
In particular, setting $F = \mathcal{B}f$, one obtains
\begin{equation}\label{eq:energy}
\int_{\mathbb{R}^2} |V_\varphi f(x,\omega)|^2 \, dx\, d\omega
= \int_{\mathbb{C}} |F(z)|^2 e^{-\pi |z|^2} \, dA(z)
= \|F\|_{\mathcal{F}^2}^2.
\end{equation}

\medskip

\noindent
\subsubsection*{Weyl operators and invariance.}
For $a \in \mathbb{C}$, define the operator
\begin{equation}\label{eq:Ta}
(T_a F)(z) := e^{\pi a z - \frac{\pi}{2}|a|^2} F(z - a).
\end{equation}
Then $T_a$ is unitary on $\mathcal{F}^2(\mathbb{C})$. Indeed,
\begin{align*}
\|T_aF\|_{\mathcal F^2}^2
&=
\int_{\C} e^{2\pi \Re(az)-\pi |a|^2}\,|F(z-a)|^2 e^{-\pi |z|^2}\,dA(z) \\
&=
\int_{\C} |F(z-a)|^2 e^{-\pi |z-a|^2}\,dA(z)
=
\|F\|_{\mathcal F^2}^2,
\end{align*}
because
\[
2\Re(az)-|a|^2-|z|^2=-|z-a|^2.
\]
Moreover,
\begin{equation}\label{eq:uF-translate}
u_{T_a F}(z) = u_F(z - a).
\end{equation}
Indeed, since
\(2\Re(az)-|a|^2-|z|^2 = -|z-a|^2\),
\begin{equation}\label{eq:u-translation-global}
u_{T_aF}(z)
=
|T_aF(z)|^2e^{-\pi |z|^2}
=
e^{2\pi\Re(az)-\pi|a|^2}|F(z-a)|^2e^{-\pi |z|^2}
=
|F(z-a)|^2e^{-\pi |z-a|^2}=u_F(z-a).
\end{equation}
In particular, the functionals $I_F$ defined by,
for \(s>0\),
\[
I_F(s):=\sup_{|\Omega|=s}\int_\Omega u_F\,dA,
\]
satisfy, by the translation-invariance of the Lebesgue measure on \(\C\),
\begin{equation}\label{eq:I-translation-global}
I_{T_a F}(s) = I_F(s), \qquad a \in \mathbb{C},
\end{equation}
reflecting the invariance of the Gaussian STFT under time--frequency shifts;
see \cite{Daubechies,Grochenig,Kulikov-Oliveira-Ramos,Goncalves-Oliveira-Steinerberger}.

\medskip

\noindent
\subsubsection*{Connection with localization.}
If $F = \mathcal{B}f$, then for every measurable $\Omega \subset \mathbb{C}$,
\[
\int_\Omega u_F \, dA
= \int_{\widetilde{\Omega}} |V_\varphi f(x,\omega)|^2 \, dx\, d\omega,
\]
where $\widetilde{\Omega} = \{(x,\omega) : (x,-\omega) \in \Omega\}$.
Thus, concentration properties of the STFT are naturally encoded by the function
\begin{equation}\label{eq:uF}
u_F(z) := |F(z)|^2 e^{-\pi |z|^2}, \qquad z \in \mathbb{C},
\end{equation}
which is the main object of interest in the Bargmann--Fock setting, in terms of the
quantities $I_F(s)$.

\medskip

We next collect a few elementary facts that will be used repeatedly.

\begin{lemma}\label{lem:Fock-basic}
Let $F,G \in \mathcal{F}^2(\mathbb{C})$. Then:
\begin{enumerate}
\item[(i)] The space $\mathcal{F}^2(\mathbb{C})$ is a reproducing kernel Hilbert space with kernel
\[
K_w(z) = e^{\pi z \overline{w}}, \qquad z,w \in \mathbb{C}.
\]

\item[(ii)] (Sharp pointwise bound) For every $z \in \mathbb{C}$,
\[
u_F(z) \le \|F\|_{\mathcal{F}^2}^2,
\]
and equality holds at some point if and only if $F$ is proportional to $F_w$
for some $w \in \mathbb{C}$.

\item[(iii)] $u_F(z) \to 0$ as $|z| \to \infty$.

\item[(iv)] One has the estimate
\[
\|u_F - u_G\|_{L^1(\mathbb{C})}
\le \big( \|F\|_{\mathcal{F}^2} + \|G\|_{\mathcal{F}^2} \big)
\|F - G\|_{\mathcal{F}^2}.
\]
\end{enumerate}
\end{lemma}

\begin{proof}
(i) For \(w\in \C\), the function \(K_w(z)=e^{\pi z\overline w}\) is entire and
satisfies, after completing the square via
\(2\Re(z\overline w)-|z|^2=-|z-w|^2+|w|^2\),
\[
\|K_w\|_{\mathcal F^2}^2
=
\int_{\C} e^{2\pi \Re(z\overline w)-\pi |z|^2}\,dA(z)
=
e^{\pi |w|^2}\int_{\C}e^{-\pi |z-w|^2}\,dA(z)
=
e^{\pi |w|^2},
\]
since the Gaussian \(e^{-\pi |\cdot|^2}\) integrates to \(1\) over \(\C\) (with
the Lebesgue measure \(dA\) on \(\C\simeq \mathbb R^2\)). In particular
\(K_w\in \mathcal F^2(\C)\). Now, expanding \(F\in \mathcal F^2(\C)\) along the
orthonormal basis \(\{e_k\}_{k\geq 0}\) of Lemma \ref{lem:monomials} and using
the Taylor expansion
\(K_w(z)=\sum_{k\geq 0}\frac{\pi^k\overline w^k}{k!}z^k=\sum_{k\geq 0}\overline{e_k(w)}\,e_k(z)\),
one obtains the reproducing identity
\[
\langle F,K_w\rangle_{\mathcal F^2}
=
\sum_{k\geq 0} a_k\,e_k(w)
=
F(w),
\qquad F(z)=\sum_{k\geq 0}a_k e_k(z).
\]
Point evaluation is therefore continuous, and \(F\) is entire as a uniform
limit on compact sets of polynomials. This is the standard description of the
Bargmann--Fock space, see \cite{Bargmann,Folland,Grochenig,Zhu}.

(ii) Applying Cauchy--Schwarz to the reproducing identity from (i),
\begin{equation}\label{eq:pointwise}
|F(z)|
=
|\langle F,K_z\rangle_{\mathcal F^2}|
\le
\|F\|_{\mathcal F^2}\,\|K_z\|_{\mathcal F^2}
=
\|F\|_{\mathcal F^2}\,e^{\frac{\pi}{2}|z|^2},
\end{equation}
which immediately gives \(u_F(z)=|F(z)|^2 e^{-\pi |z|^2}\le \|F\|_{\mathcal F^2}^2\).
The case of equality at a point \(z=w\) corresponds to equality in
Cauchy--Schwarz, that is, \(F=c\,K_w\) for some \(c\in \C\); writing
\(F_w:=K_w\), this proves the rigidity statement. The sharpness of this
pointwise bound goes back to \cite{Lieb}.

(iii) Fix \(\varepsilon>0\), and choose a polynomial \(P\) such that
\[
\|F-P\|_{\mathcal F^2}<\varepsilon
\]
using the density of polynomials in \(\mathcal F^2(\C)\). By the pointwise bound from (ii),
\[
|F(z)-P(z)|e^{-\frac{\pi}{2}|z|^2}\le \varepsilon
\qquad\text{for every }z\in \C.
\]
Since \(P\) is a polynomial,
\[
P(z)e^{-\frac{\pi}{2}|z|^2}\to 0
\qquad\text{as }|z|\to\infty.
\]
Hence
\[
\limsup_{|z|\to\infty}|F(z)|e^{-\frac{\pi}{2}|z|^2}\le \varepsilon.
\]
As \(\varepsilon>0\) is arbitrary, we conclude that
\[
|F(z)|e^{-\frac{\pi}{2}|z|^2}\to 0,
\]
that is, \(u_F(z)\to 0\) as \(|z|\to\infty\).

(iv) Writing
\[
|u_F - u_G|
= \big||F|^2 - |G|^2\big| e^{-\pi|z|^2}
\le |F-G|(|F|+|G|) e^{-\pi|z|^2},
\]
the claim follows by Cauchy--Schwarz.
\end{proof}

\medskip

We now note the following crucial fact about the Bargmann transform: it
transforms the basis of normalized Hermite functions \(\{h_k\}_{k\geq 0}\) on
\(L^2(\R)\) (see \cite{Thangavelu}) onto normalized monomials, yielding a
natural orthonormal basis on $\mathcal{F}^2(\C).$ More precisely, one has
\(\mathcal B h_k=e_k\) for every \(k\ge 0\), see \cite{Bargmann,Folland,Grochenig}.

\begin{lemma}\label{lem:monomials}
The functions
\[
e_k(z) := b_k z^k, \qquad k \ge 0,
\]
form an orthonormal basis of $\mathcal{F}^2(\mathbb{C})$, where
\[
b_k = \left( \frac{\pi^k}{k!} \right)^{1/2}.
\]
In particular,
\[
\langle e_k, e_m \rangle_{\mathcal{F}^2} = \delta_{k,m}.
\]
\end{lemma}

\begin{proof}
Writing \(z=re^{i\theta}\) and using \(dA(z)=r\,dr\,d\theta\) on \(\C\),
a direct computation in polar coordinates gives
\[
\int_{\mathbb{C}} |z|^{2k} e^{-\pi |z|^2}\,dA(z)
= \int_0^{2\pi}\!\!\int_0^\infty r^{2k+1} e^{-\pi r^2}\,dr\,d\theta
= 2\pi \int_0^\infty r^{2k+1} e^{-\pi r^2}\,dr
= \frac{k!}{\pi^k},
\]
where the last identity follows from the substitution \(s=\pi r^2\), giving
\(2\pi\int_0^\infty r^{2k+1}e^{-\pi r^2}\,dr=\pi^{-k}\int_0^\infty s^k e^{-s}\,ds
=\pi^{-k}\Gamma(k+1)\). Thus
\(\|z^k\|_{\mathcal{F}^2}^2 = \frac{k!}{\pi^k}\), which yields the normalization.
For orthogonality, again writing \(z=re^{i\theta}\),
\[
\langle z^k,z^m\rangle_{\mathcal F^2}
=
\int_0^\infty r^{k+m+1}e^{-\pi r^2}\,dr\int_0^{2\pi}e^{i(k-m)\theta}\,d\theta
=
0
\qquad\text{for }k\neq m.
\]
For completeness, observe that if \(F\in \mathcal F^2(\C)\) satisfies
\(\langle F,e_k\rangle_{\mathcal F^2}=0\) for every \(k\ge 0\), then in
particular all Taylor coefficients of \(F\) at the origin vanish: since
\(F=\mathcal Bf\) for a unique \(f\in L^2(\R)\), and \(\mathcal Bh_k=e_k\) (so
that \(\{e_k\}\) is the image of the orthonormal basis of Hermite functions
under the unitary \(\mathcal B\)), the family \(\{e_k\}_{k\ge 0}\) is total in
\(\mathcal F^2(\C)\). Equivalently, one may invoke directly the fact that any
entire function of order at most \(2\) lying in \(\mathcal F^2(\C)\) is the
limit of its Taylor polynomials in the Fock norm; see
\cite{Bargmann,Zhu}.
\end{proof}

\medskip

\begin{lemma}\label{lem:growth}
Let $F \in \mathcal{F}^2(\mathbb{C})$. Then:
\begin{enumerate}
\item[(i)] $F$ admits the expansion
\[
F(z) = \sum_{k=0}^\infty a_k e_k(z),
\qquad \sum_{k=0}^\infty |a_k|^2 = \|F\|_{\mathcal{F}^2}^2.
\]

\item[(ii)] (Growth bound) For every $z \in \mathbb{C}$,
\[
|F(z)| \le \|F\|_{\mathcal{F}^2} e^{\frac{\pi}{2}|z|^2}.
\]

\item[(iii)] (Derivative bounds) For every $z \in \mathbb{C}$,
\[
|F'(z)| \le C(1+|z|)\, \|F\|_{\mathcal{F}^2} e^{\frac{\pi}{2}|z|^2}.
\]
\end{enumerate}
\end{lemma}

\begin{proof}
(i) follows from the orthonormal basis in Lemma \ref{lem:monomials}.

(ii) follows from the reproducing kernel estimate:
\[
|F(z)| = |\langle F, K_z \rangle|
\le \|F\|_{\mathcal{F}^2} \|K_z\|_{\mathcal{F}^2}
= \|F\|_{\mathcal{F}^2} e^{\frac{\pi}{2}|z|^2}.
\]

(iii) Differentiating the reproducing identity
\(F(z)=\langle F,K_z\rangle_{\mathcal F^2}\) in the parameter \(z\), and using
\(\partial_z K_w(z) =\pi \overline{w}\,K_w(z)\) (which is to be read in the
antiholomorphic variable \(w\), namely \(\overline{\partial_{\overline w}K_z(w)}=\pi w K_z(w)\)
once one views \(z\mapsto K_z\) as an antiholomorphic family), we obtain
\[
F'(z)=\pi \langle F,\overline{w}K_z(w)\rangle_{\mathcal F^2}.
\]
Cauchy--Schwarz then gives
\[
|F'(z)|
\le
\pi \|F\|_{\mathcal F^2}\,\|\overline{w}K_z\|_{\mathcal F^2}.
\]
Now, by completing the square as in (i),
\[
\|\overline{w}K_z\|_{\mathcal F^2}^2
=
\int_{\C} |w|^2 e^{2\pi\Re(w\overline z)-\pi |w|^2}\,dA(w)
=
e^{\pi |z|^2}\int_{\C} |w+z|^2 e^{-\pi |w|^2}\,dA(w),
\]
after the change of variables \(w\mapsto w+z\). Expanding
\(|w+z|^2=|w|^2+2\Re(w\overline z)+|z|^2\), and using that
\(\int_\C e^{-\pi |w|^2}\,dA(w)=1\),
\(\int_\C |w|^2 e^{-\pi |w|^2}\,dA(w)=\frac{1}{\pi}\), and
\(\int_\C w\,e^{-\pi |w|^2}\,dA(w)=0\) (by rotational invariance), we conclude
\[
\int_{\C} |w+z|^2 e^{-\pi |w|^2}\,dA(w)
=
\frac{1}{\pi}+|z|^2
\le
C(1+|z|^2),
\]
for some absolute constant \(C>0\). Therefore
\[
|F'(z)| \le C(1+|z|)\,\|F\|_{\mathcal F^2}e^{\frac{\pi}{2}|z|^2}.
\]
\end{proof}

\medskip

\begin{lemma}\label{lem:max}
Let $F \in \mathcal{F}^2(\mathbb{C})$ and set
\[
T := \sup_{z \in \mathbb{C}} u_F(z).
\]
Then $0 \le T \le \|F\|_{\mathcal{F}^2}^2$, and $T = \|F\|_{\mathcal{F}^2}^2$
if and only if $F$ is proportional to $F_w$ for some $w \in \mathbb{C}$.
\end{lemma}

\begin{proof}
The bound follows from Lemma \ref{lem:Fock-basic}(ii). The characterization of
equality is the same as in the reproducing kernel inequality.
    \end{proof}

\medskip

    We proceed with a few elementary facts.

\begin{lemma}\label{lem:basic-properties-global}
Let \(F,G\in \mathcal F^2(\C)\). Then:
\begin{enumerate}
    \item[(i)] \(I_F(s)\) is attained by a superlevel set of \(u_F\), and
    \[
    0\le I_F(s)\le \|F\|_{\mathcal F^2}^2.
    \]
    \item[(ii)] For every \(c\in \C\),
    \[
    I_{cF}(s)=|c|^2 I_F(s).
    \]
    \item[(iii)] One has the Lipschitz bound
    \[
    |I_F(s)-I_G(s)|
    \le
    \|u_F-u_G\|_{L^1(\C)}
    \le
    \big(\|F\|_{\mathcal F^2}+\|G\|_{\mathcal F^2}\big)\,
    \|F-G\|_{\mathcal F^2}.
    \]
\end{enumerate}
\end{lemma}

\begin{proof}
The identity
\[
\int_\C u_F\,dA=\int_\C |F(z)|^2e^{-\pi |z|^2}\,dA(z)=\|F\|_{\mathcal F^2}^2
\]
is immediate from the definition of the Fock norm. Since \(F\in\mathcal F^2\),
the standard reproducing-kernel estimate implies that \(u_F\) is bounded and
continuous, and in fact \(u_F(z)\to 0\) as \(|z|\to \infty\); thus
\(u_F\in C_0(\C)\cap L^1(\C)\).

The existence of a maximizing set for \(I_F(s)\), given by a superlevel set of
\(u_F\), is the classical bathtub principle. The bound
\[
0\le I_F(s)\le \int_\C u_F\,dA=\|F\|_{\mathcal F^2}^2
\]
is immediate. The scaling identity in (ii) is obvious.

For (iii), for every measurable \(\Omega\subset \C\) with \(|\Omega|=s\),
\[
\left|\int_\Omega u_F\,dA-\int_\Omega u_G\,dA\right|
\le \int_\C |u_F-u_G|\,dA.
\]
Taking suprema over \(|\Omega|=s\) gives
\[
|I_F(s)-I_G(s)|\le \|u_F-u_G\|_{L^1(\C)}.
\]
Finally,
\begin{align*}
\|u_F-u_G\|_{L^1(\C)}
&=
\int_\C \big||F|^2-|G|^2\big|\,e^{-\pi |z|^2}\,dA \\
&\le
\int_\C |F-G|(|F|+|G|)e^{-\pi |z|^2}\,dA \\
&\le
\|F-G\|_{\mathcal F^2}\big(\|F\|_{\mathcal F^2}+\|G\|_{\mathcal F^2}\big)
\end{align*}
by Cauchy--Schwarz.
\end{proof}

\subsection{From time-frequency analysis to free boundary problems}\label{sec:fourier-to-fboundary}

Having translated, in the previous subsection, the time--frequency
localization problem into a question about the vanishing of a Fourier
transform along the two complex lines $\eta=\pm i\xi$, we now bridge the gap
between this Fourier-analytic statement and the free boundary formulation
that will dominate the remainder of the paper. The bridge is provided by a
divisibility lemma of Ehrenpreis--Malgrange flavour: any compactly supported
distribution whose Fourier transform vanishes on the characteristic variety
$\{\eta=\pm i\xi\}$ of the Laplacian must be the Laplacian of a compactly
supported $L^2$ function which, in our setting, is moreover continuous off a
finite set of points.

In order to state the lemma cleanly, we make precise the class of
distributions with which we work. We say that a compactly supported
distribution $\omega$ is a \emph{singular-continuous free measure} if it
defines a (finite, signed) Radon measure on $\R^2$ whose Lebesgue
decomposition with respect to planar Lebesgue measure $\mathrm{d}x\,
\mathrm{d}y$ takes the form
\[
\omega = \omega_{\mathrm{ac}} + \omega_{\mathrm{pp}}
       = f\,\mathrm{d}x\,\mathrm{d}y + \sum_{j=1}^N c_j\,\delta_{p_j},
\]
with \(f\in L^\infty(\R^2)\) compactly supported, \(c_j\in \R\setminus\{0\}\), and
$p_1,\ldots,p_N\in \R^2$ pairwise distinct. In other words, the continuous
singular part $\omega_{\mathrm{cs}}$ in the Lebesgue--Radon--Nikodym
decomposition $\omega=\omega_{\mathrm{ac}}+\omega_{\mathrm{cs}}+
\omega_{\mathrm{pp}}$ is required to vanish identically, and the pure
point part is finite. The terminology \emph{free measure} is suggested by
the role of $\omega$ as the right-hand side of a free boundary problem
arising in the next subsection.

\begin{lemma}\label{lemma:fourier-to-fboundary}
Let $\omega\in \mathcal{S}'(\R^2)$ be a compactly supported singular-continuous
free distribution such that
\[
\widehat{\omega}(w,\pm iw)=0,\qquad \forall w\in\C.
\]
Then there exist a finite set $P\subset \R^2$ and a compactly supported
function $u\in L^2(\R^2)$ such that
\[
\Delta u=\omega
\]
in the sense of distributions, and
\[
u\in C(\R^2\setminus P).
\]
\end{lemma}

\begin{proof}
By the Paley--Wiener--Schwartz theorem (see, e.g.,
\cite[Theorem~7.3.1]{Hormander} or \cite[\S 9.3]{Folland}), the Fourier
transform $\widehat\omega(\xi,\eta)$ extends to an entire function on $\C^2$
of exponential type, satisfying the bound
\[
|\widehat\omega(\xi,\eta)|\le C(1+|\xi|+|\eta|)^{N}\,e^{R(|\Im\xi|+|\Im\eta|)}
\]
for some $C,R>0$ and some integer $N\ge 0$ depending on the order of
$\omega$, where $R$ is the radius of any disk containing
$\operatorname{supp}\omega$. In particular, $\widehat\omega$ admits an
everywhere convergent power series expansion
\[
\widehat{\omega}(\xi,\eta)=\sum_{k,l\ge 0} a_{k,l}\,\xi^k\eta^l,
\qquad (\xi,\eta)\in\C^2.
\]
Using the vanishing hypothesis on the complex lines $\eta=\pm i\xi$, we may
group terms by total homogeneity: setting $\eta=\pm iw$, $\xi=w$, and
collecting powers of $w$,
\[
0=\widehat{\omega}(w,\pm iw)
   =\sum_{n\ge 0}\Bigl(\sum_{k=0}^n a_{k,n-k}(\pm i)^{\,n-k}\Bigr)w^n,
\qquad \forall w\in\C.
\]
Since this holds for all $w\in\C$, comparing coefficients yields
\[
\sum_{k=0}^n a_{k,n-k}(\pm i)^{\,n-k}=0,\qquad \forall n\ge 0.
\]
For each $n\ge 0$, define the auxiliary one-variable polynomial
\[
p_n(t):=\sum_{k=0}^n a_{k,n-k}\, t^{\,n-k}.
\]
The two scalar equations above precisely say that $p_n(\pm i)=0$, so
$(t-i)(t+i)=t^2+1$ divides $p_n(t)$. Hence we may write
\[
p_n(t)=(t^2+1)\,q_n(t),
\]
with $q_n$ a polynomial of degree at most $n-2$ (and $q_n\equiv 0$ when
$n<2$).

We now reassemble the homogeneous pieces. Let
\[
H_n(\xi,\eta):=\sum_{k=0}^n a_{k,n-k}\,\xi^k\eta^{\,n-k},
\qquad n\ge 0,
\]
so that $\widehat\omega=\sum_{n\ge 0}H_n$ is the decomposition of
$\widehat\omega$ into homogeneous components (of degree $n$). The argument
above shows that each $H_n$ vanishes on the two complex lines
$\eta=\pm i\xi$, hence is divisible, as a homogeneous polynomial in two
variables, by $\xi^2+\eta^2$. Concretely,
\[
H_n(\xi,\eta)=(\xi^2+\eta^2)\,G_{n-2}(\xi,\eta)
\]
for a (uniquely determined) homogeneous polynomial $G_{n-2}$ of degree
$n-2$, with $G_{n-2}\equiv 0$ for $n<2$. Summing the homogeneous pieces, the
formal series
\[
F(\xi,\eta):=\sum_{n\ge 2} G_{n-2}(\xi,\eta)
\]
converges absolutely on every polydisc on which $\widehat\omega$ does, and
defines an entire function on $\C^2$ satisfying
\[
\widehat{\omega}(\xi,\eta)=(\xi^2+\eta^2)\,F(\xi,\eta).
\]
Since $\widehat\omega$ has a zero of order at least two at the origin
(coming from the vanishing of $H_0$ and $H_1$, which is forced by
$p_0(\pm i)=p_1(\pm i)=0$), this factorization is consistent at $0$ and $F$
is entire there as well.

We next quantify the size of $F$ on $\R^2$. Since $\omega$ is a finite
(signed) Radon measure with compact support, its Fourier transform is in
fact bounded on $\R^2$:
\[
|\widehat{\omega}(\xi,\eta)|\le \|\omega\|_{\mathrm{TV}},
\qquad (\xi,\eta)\in\R^2,
\]
where $\|\omega\|_{\mathrm{TV}}$ is the total variation; equivalently,
the Paley--Wiener bound holds with $N=0$. On $\R^2$, $\xi^2+\eta^2=
|\xi|^2+|\eta|^2$, so for $|\xi|^2+|\eta|^2\ge 1$ we have $\xi^2+\eta^2
\ge \tfrac{1}{2}(1+|\xi|^2+|\eta|^2)$, whence
\[
|F(\xi,\eta)|=\frac{|\widehat\omega(\xi,\eta)|}{\xi^2+\eta^2}
\le \frac{2\|\omega\|_{\mathrm{TV}}}{1+|\xi|^2+|\eta|^2},
\qquad |\xi|^2+|\eta|^2\ge 1.
\]
Near the origin, $F$ is entire (as established above) and hence bounded;
absorbing this region into the constants we obtain
\[
|F(\xi,\eta)|\le \frac{C''}{1+|\xi|^2+|\eta|^2},
\qquad (\xi,\eta)\in\R^2,
\]
for some $C''>0$. Therefore $F|_{\R^2}\in L^2(\R^2)\cap L^1(\R^2)$. On
$\C^2$, the analogous Paley--Wiener-type bound
\[
|F(\xi,\eta)|\le C'\,
\frac{e^{R(|\Im \xi|+|\Im \eta|)}}{1+|\xi|^2+|\eta|^2}
\]
follows from the same argument off a neighbourhood of $\{\xi^2+\eta^2=0\}$
and from boundedness of $F$ on a fixed compact neighbourhood of $0$ inside
$\C^2$, after enlarging $C'$ if necessary.

By the Paley--Wiener--Schwartz theorem applied in the reverse direction
(see again \cite[Theorem~7.3.1]{Hormander}), the entire function $F$
satisfies an exponential-type bound on $\C^2$ with the same exponential
rate $R$ as $\widehat\omega$, and integrable decay along $\R^2$, so it is
the Fourier--Laplace transform of a distribution of order $\le N$
compactly supported in the closed disk $\overline{B(0,R)}$. The $L^2$
bound on the real locus shows that this distribution is, in fact, an
$L^2(\R^2)$ function. Let $u\in L^2(\R^2)$ denote this compactly supported
representative, characterised by $\widehat u = F$ on $\R^2$ and
$\operatorname{supp} u\subset \overline{B(0,R)}$.

Computing distributional Laplacians on the Fourier side, with the
convention $\widehat g(\xi,\eta)=\int g(x,y)e^{-2\pi i(x\xi+y\eta)}\,
\mathrm{d}x\,\mathrm{d}y$, we have
\[
\widehat{\Delta u}(\xi,\eta) = -4\pi^2(\xi^2+\eta^2)\,\widehat{u}(\xi,\eta)
= -4\pi^2(\xi^2+\eta^2)\,F(\xi,\eta) = -4\pi^2\,\widehat\omega(\xi,\eta);
\]
with any other standard normalisation, the same identity holds up to a
nonzero real multiplicative factor. In all cases, $\Delta u$ is a constant
nonzero real multiple of $\omega$. Replacing $u$ by $\lambda u$ for the
appropriate (nonzero, real) $\lambda$ -- this is the harmless rescaling
referred to above -- we may, and do, assume that
\[
\Delta u = \omega
\]
in the sense of distributions. The rescaling preserves all the previously
established properties of $u$: compact support, membership in $L^2(\R^2)$,
and the Fourier identity $\widehat u = \lambda F$.

To pin down $u$ pointwise, observe that since $\omega$ is a compactly
supported Radon measure on $\R^2$, the convolution
\[
v(x,y) := \frac{1}{2\pi}(\log|\cdot|*\omega)(x,y)
\]
is a well-defined locally integrable function on $\R^2$ which solves
$\Delta v = \omega$ in the sense of distributions; here $G(z)=\log|z|$ is
(up to the standard normalisation $\tfrac{1}{2\pi}$) the fundamental solution
of the Laplacian in $\R^2$, see, e.g., \cite[Ch.~II]{Garnett} or
\cite[\S 6.2]{Folland} for the relevant logarithmic-potential theory, and
\cite{Aharonov-Schiffer-Zalcman} for related representations in the planar
quadrature-domain context. The difference $u-v$ is then a tempered
distribution which is harmonic on $\R^2$ (since
$\Delta(u-v)=\omega-\omega=0$), and hence, being harmonic, real-analytic.
Outside the union of the supports of $u$ and $\omega$, the function $v$ is
in fact real-analytic and grows at most logarithmically at infinity, with
$v(x)=\frac{M}{2\pi}\log|x|+O(|x|^{-1})$ as $|x|\to\infty$ where
$M=\omega(\R^2)$, while $u$ vanishes there because of its compact support.
Hence $u-v$ is harmonic on $\R^2$ and grows at most logarithmically at
infinity; a standard refinement of the harmonic Liouville theorem (any
harmonic function on $\R^2$ growing slower than a polynomial of degree $1$
is constant) then forces $u-v$ to be a constant, and this constant must
vanish to be consistent with $u\in L^2(\R^2)$ and the local integrability
of $v$. Equivalently, and more transparently: by the compact support of
$\omega$ together with $\widehat u = F$ on $\R^2$, the function $u$ is
uniquely determined among compactly supported $L^2$ solutions of
$\Delta u = \omega$, since any two such would differ by an entire $L^2$
function which is harmonic, hence vanishes. Thus
\[
u(x,y)=\frac{1}{2\pi}\int_{\R^2}\log|(x,y)-(x',y')|\,\mathrm{d}\omega(x',y'),
\]
the logarithmic potential of $\omega$.

We finally analyse the regularity of $u$. Decomposing
$\omega = f\,\mathrm{d}x\,\mathrm{d}y+\sum_{j=1}^N c_j\,\delta_{p_j}$ as in
the definition of singular-continuous free measure, we obtain
correspondingly
\[
u(x,y)=\frac{1}{2\pi}\int_{\R^2}\log|(x,y)-(x',y')|\,f(x',y')\,
\mathrm{d}x'\mathrm{d}y'
+\frac{1}{2\pi}\sum_{j=1}^N c_j\,\log|(x,y)-p_j|.
\]
The first summand is the logarithmic potential of the bounded, compactly
supported density \(f\) and is therefore continuous on all of \(\R^2\) -- this
is a classical fact, see e.g.\ \cite[Ch.~II, Theorem~1.3]{Garnett} or
\cite[\S 6.2]{Folland}. It is here that the absence of a continuous
singular part in $\omega$ is essential: a generic continuous singular
measure (e.g.\ the Hausdorff measure on a Cantor-type set of positive
$1$-capacity) gives rise to a logarithmic potential which need not be
continuous, only quasi-continuous, and the conclusion of the lemma would
fail in that generality. Each term $\log|(x,y)-p_j|$ is continuous on
$\R^2\setminus\{p_j\}$, with a logarithmic singularity at $p_j$. Setting
$P:=\{p_1,\ldots,p_N\}$, we conclude that $u\in C(\R^2\setminus P)$, as
claimed. The fact that $u$ is compactly supported has already been
established in the construction above and is consistent with the
representation as a logarithmic potential, which differs from a compactly
supported function only by a globally harmonic correction that, as
explained, must vanish identically.
\end{proof}

We postpone the proof of Theorem~\ref{thm:eig-loc-U} until after the
perturbative inverse theorem. The same inverse/free-boundary mechanism used
to construct domains for prescribed eigenfunctions also recovers the
classical Abreu--D\"orfler disk theorem in the simply connected setting; see
Section~\ref{sec:classical-rigidity} below.


% ==============================================================
% Section 3
% ==============================================================
\section{Proof of Theorem \ref{thm:perturb}}

\subsection{The disk inverse step}\label{subsec:disk-inverse-step}

The proof rests on a local inverse theorem at the Daubechies disk. The result
below is self-contained, but its point of view is inspired by perturbative
inverse-source and quadrature-domain methods: Isakov's local free-boundary
scheme for inverse source problems, Cherednichenko's local inverse logarithmic
potential theorem, and the quadrature-domain framework developed by
Gustafsson, Shahgholian, Karp, and Aharonov--Schiffer--Zalcman. In the present
disk-based setting the required inverse step can be proved directly by
diagonalizing the full moment map in Fourier modes.

Fix \(0<\lambda<1\), and let \(r_\lambda>0\) be determined by
\[
1-e^{-\pi r_\lambda^2}=\lambda .
\]
For a real-valued \(2\pi\)-periodic function \(\varphi\) with
\(\|\varphi\|_{C^0}<r_\lambda/2\), set
\[
U_\varphi
:=
\{re^{i\theta}:0\le r<r_\lambda+\varphi(\theta)\}.
\]
For \(\tau>0\), let \(\mathcal A_\tau^\R\) be the Banach space of real-valued
periodic functions
\begin{equation}\label{eq:Atau-def}
\varphi(\theta)=\sum_{k\in\Z}\varphi_k e^{ik\theta},
\qquad
\varphi_{-k}=\overline{\varphi_k},
\end{equation}
with norm
\begin{equation}\label{eq:Atau-norm}
\|\varphi\|_{\mathcal A_\tau}
:=
\sum_{k\in\Z}|\varphi_k|e^{\tau |k|}<\infty .
\end{equation}
This is the real subspace of the Wiener-type Banach algebra of functions on
\(\T:=\R/2\pi\Z\) admitting a holomorphic extension to the strip
\(|\Im\theta|<\tau\) with absolutely summable Fourier expansion. It is a Banach
algebra, continuously embedded into \(C^m(\T)\) for every \(m\ge0\), and in
particular into the H\"older space \(C^{1,\alpha}(\T)\) for every
\(\alpha\in(0,1)\). Small elements of \(\mathcal A_\tau^\R\) therefore produce
real-analytic \(C^2\) Jordan domains, given as radial graphs over the disk
\(D_{r_\lambda}\).

For sequences indexed by \(j\ge0\), the natural target Banach space dual to
\(\mathcal A_\tau^\R\) under the moment map of the proposition below is
\begin{equation}\label{eq:Ytau-def}
\mathcal Y_\tau:=
\Big\{y=(y_j)_{j\ge0}\colon y_0\in\R,\ \|y\|_{\mathcal Y_\tau}
:=\tfrac{|y_0|}{a_0}+2\sum_{j=1}^\infty e^{\tau j}\tfrac{|y_j|}{a_j}<\infty\Big\},
\end{equation}
with the weight sequence \(\{a_j\}_{j\ge0}\) given by \eqref{eq:aj-weight}
below. The factor of two in \eqref{eq:Ytau-def} accounts for the conjugate
pairing \(\varphi_{-j}=\overline{\varphi_j}\) used in \eqref{eq:Atau-def}.

\begin{proposition}[Disk symmetrized inverse theorem]\label{prop:disk-inverse-tailored}
Let \(N\ge1\), \(0<\lambda<1\), and \(r_\lambda\) be as above. There are
\(\tau>0\), \(\varepsilon>0\), and a \(C^1\) map
\[
b=(b_1,\ldots,b_N)\longmapsto \varphi_b\in\mathcal A_\tau^\R,
\qquad
\varphi_0=0,
\]
defined for \(\max_j|b_j|<\varepsilon\), such that, for
\[
P_b(z):=1+\sum_{j=1}^N b_jz^j,
\qquad
P_b^\#(z):=\overline{P_b(\overline z)},
\qquad
Q_b(w):=P_b^\#(\partial_w/\pi)P_b(w),
\]
and
\[
U_b:=U_{\varphi_b},
\]
one has
\begin{equation}\label{eq:disk-inverse-identity}
\int_{U_b}|P_b(z)|^2e^{-\pi |z|^2}e^{\pi\overline z w}\,dA(z)
=
\lambda Q_b(w),
\qquad w\in\C .
\end{equation}
Moreover,
\[
\|\varphi_b\|_{C^2(\mathbb S^1)}
\le C_{\lambda,N}\max_{1\le j\le N}|b_j|,
\]
and the branch is locally unique among sufficiently small analytic radial
graphs in \(\mathcal A_\tau^\R\).
\end{proposition}

\begin{proof}
We solve the coefficient equations obtained by expanding the entire identity
\eqref{eq:disk-inverse-identity}. Write
\[
Q_b(w)=\sum_{j=0}^N q_j(b)w^j,
\]
and \(q_j(b)=0\) for \(j>N\). The constant coefficient is real:
\[
q_0(b)=1+\sum_{k=1}^N \frac{k!}{\pi^k}|b_k|^2.
\]
Since
\[
e^{\pi\overline z w}
=
\sum_{j=0}^\infty \frac{\pi^j}{j!}\overline z^{\,j}w^j,
\]
the desired identity is equivalent to the infinite system
\begin{equation}\label{eq:disk-moment-system}
\frac{\pi^j}{j!}\int_{U_\varphi}
|P_b(z)|^2\overline z^{\,j}e^{-\pi |z|^2}\,dA(z)
=
\lambda q_j(b),
\qquad j\ge0 .
\end{equation}

We introduce the weighted sequence space adapted to the linearization at the
circle. Put
\begin{equation}\label{eq:aj-weight}
a_j:=2\pi e^{-\pi r_\lambda^2}r_\lambda^{j+1}\frac{\pi^j}{j!},
\qquad j\ge0,
\end{equation}
so that \(\{a_j\}\) decays super-exponentially in \(j\), and \(\mathcal Y_\tau\)
is the Banach space defined in \eqref{eq:Ytau-def}.
Define
\[
\mathcal F_j(b,\varphi)
:=
\frac{\pi^j}{j!}\int_{U_\varphi}
|P_b(z)|^2\overline z^{\,j}e^{-\pi |z|^2}\,dA(z)
-\lambda q_j(b),
\qquad j\ge0 .
\]
Thus \(\mathcal F(b,\varphi)=0\) is precisely
\eqref{eq:disk-moment-system}. At \(b=0\), \(\varphi=0\), the domain is the
disk \(D_{r_\lambda}\), and radial symmetry gives
\[
\frac{\pi^j}{j!}\int_{D_{r_\lambda}}
\overline z^{\,j}e^{-\pi |z|^2}\,dA(z)
=
\begin{cases}
\lambda, & j=0,\\
0, & j\ge1.
\end{cases}
\]
Hence \(\mathcal F(0,0)=0\).

We next compute the Fr\'echet derivative of the moment map with respect to
the boundary graph. Let
\(\psi(\theta)=\sum_{k\in\Z}\psi_k e^{ik\theta}\in\mathcal A_\tau^\R\), and
recall the polar-coordinate formula
\[
\int_{U_\varphi}H(z)\,dA(z)
=
\int_0^{2\pi}\!\!\int_0^{r_\lambda+\varphi(\theta)}
H(re^{i\theta})\,r\,dr\,d\theta,
\]
valid for any continuous \(H\colon\C\to\C\). Differentiating the right-hand
side at \(\varphi=0\) in the direction of \(\psi\) by the fundamental
theorem of calculus gives
\begin{equation}\label{eq:disk-frechet-domain}
D_\varphi\!\Big(\int_{U_\varphi}H(z)\,dA(z)\Big)\bigg|_{\varphi=0}\![\psi]
=
r_\lambda\!\int_0^{2\pi}\!\psi(\theta)\,H(r_\lambda e^{i\theta})\,d\theta .
\end{equation}
Applied to \(H(z)=\frac{\pi^j}{j!}\overline z^{\,j}e^{-\pi|z|^2}\) and combined
with the orthogonality relation
\(\int_0^{2\pi}e^{i(k-j)\theta}\,d\theta=2\pi\delta_{kj}\), formula
\eqref{eq:disk-frechet-domain} yields the closed-form Fourier-multiplier
expression
\begin{equation}\label{eq:disk-frechet-multiplier}
D_\varphi\mathcal F_j(0,0)[\psi]
=
\frac{\pi^j}{j!}\,r_\lambda^{j+1}e^{-\pi r_\lambda^2}
\int_0^{2\pi}\!\psi(\theta)\,e^{-ij\theta}\,d\theta
=
a_j\,\psi_j,
\qquad j\ge0,
\end{equation}
where \(a_j\) is the weight in \eqref{eq:aj-weight}. Consequently the
Fr\'echet derivative \(D_\varphi\mathcal F(0,0)\) is the Fourier multiplier
\begin{equation}\label{eq:disk-frechet-iso}
D_\varphi\mathcal F(0,0)\colon
\mathcal A_\tau^\R\longrightarrow\mathcal Y_\tau,
\qquad
\psi=\sum_{k\in\Z}\psi_k e^{ik\theta}\longmapsto(a_j\psi_j)_{j\ge0},
\end{equation}
acting diagonally on the non-negative Fourier modes of \(\psi\).

The reality constraint \(\psi_{-j}=\overline{\psi_j}\) shows that the kernel
of the unrestricted multiplier on the complexified space consists exactly of
the negative-frequency modes; on \(\mathcal A_\tau^\R\) the reality
constraint determines those negative modes from the non-negative ones, so
\(D_\varphi\mathcal F(0,0)\) has trivial kernel as a real-linear map. We
verify directly that \eqref{eq:disk-frechet-iso} is bounded and invertible.

\emph{Boundedness.} Using \eqref{eq:Atau-norm} and \eqref{eq:Ytau-def},
\[
\|D_\varphi\mathcal F(0,0)\psi\|_{\mathcal Y_\tau}
=
\frac{|a_0\psi_0|}{a_0}+2\sum_{j=1}^\infty e^{\tau j}\frac{|a_j\psi_j|}{a_j}
=
|\psi_0|+2\sum_{j=1}^\infty e^{\tau j}|\psi_j|
=
\|\psi\|_{\mathcal A_\tau}.
\]

\emph{Invertibility.} For \(y\in\mathcal Y_\tau\) define
\begin{equation}\label{eq:disk-frechet-inverse}
\psi_0:=\frac{y_0}{a_0}\in\R,
\qquad
\psi_j:=\frac{y_j}{a_j}\ (j\ge1),
\qquad
\psi_{-j}:=\overline{\psi_j}\ (j\ge1),
\end{equation}
and \(\psi(\theta):=\sum_{k\in\Z}\psi_k e^{ik\theta}\). The above
computation gives \(\|\psi\|_{\mathcal A_\tau}=\|y\|_{\mathcal Y_\tau}\), so
\(\psi\in\mathcal A_\tau^\R\), and by construction
\(D_\varphi\mathcal F(0,0)\psi=y\). Hence \eqref{eq:disk-frechet-iso} is an
isometric isomorphism of \(\mathcal A_\tau^\R\) onto \(\mathcal Y_\tau\); in
particular
\(\|D_\varphi\mathcal F(0,0)^{-1}\|_{\mathcal Y_\tau\to\mathcal A_\tau^\R}=1\),
which provides the explicit quantitative invertibility needed for the
Banach-space implicit function theorem.

It remains to justify the use of the Banach-space implicit function theorem.
For \(\|\varphi\|_{\mathcal A_\tau}\) small, the function
\(\theta\mapsto r_\lambda+\varphi(\theta)\) extends holomorphically to the
strip \(|\Im\theta|<\tau\) and stays in a fixed annulus
\(r_\lambda/2<|r|<3r_\lambda/2\). In polar coordinates,
\[
\mathcal F_j(b,\varphi)+\lambda q_j(b)
=
\frac{\pi^j}{j!}\int_0^{2\pi}e^{-ij\theta}
\int_0^{r_\lambda+\varphi(\theta)}
P_b(re^{i\theta})P_b^\#(re^{-i\theta})r^{j+1}e^{-\pi r^2}\,dr\,d\theta .
\]
Expanding the inner integral at \(r_\lambda\) gives an absolutely convergent
power series in \(\varphi\), whose coefficients are analytic functions of
\(\theta\) and polynomial functions of \(b\). Since \(\mathcal A_\tau\) is a
Banach algebra and \(P_b(re^{i\theta})P_b^\#(re^{-i\theta})\) is a finite
trigonometric polynomial of degree at most \(N\), the resulting series and
its first derivatives in \((b,\varphi)\) converge in \(\mathcal Y_\tau\) after
possibly decreasing \(\tau\) and restricting to
\(\|b\|+\|\varphi\|_{\mathcal A_\tau}<\eta\). More explicitly, Cauchy's
estimate on a slightly wider strip gives, for \(0<\tau'<\tau\),
\[
\sum_{j\ge0}e^{\tau' j}\frac1{a_j}
\left|
\frac{\pi^j}{j!}\widehat{R_j}(j)
\right|
\le
C_{\lambda,N,\tau,\tau'}\|R\|_{\mathcal A_\tau},
\]
where \(\widehat{R_j}(j)\) denotes the \(j\)-th Fourier coefficient of the
analytic coefficient appearing in the \(j\)-th moment equation. Applying this
estimate term by term to the Taylor expansion in \(\varphi\) yields
\[
\|\mathcal F(b,\varphi)-\mathcal F(b,\psi)
-D_\varphi\mathcal F(b,\psi)[\varphi-\psi]\|_{\mathcal Y_{\tau'}}
\le
C(\|b\|+\|\varphi\|_{\mathcal A_\tau}
+\|\psi\|_{\mathcal A_\tau})
\|\varphi-\psi\|_{\mathcal A_\tau}^2,
\]
and analogous estimates for differentiation in \(b\). Shrinking the strip
once and then relabelling \(\tau'\) as \(\tau\), this proves that
\[
\mathcal F:
\{(b,\varphi):\|b\|+\|\varphi\|_{\mathcal A_\tau}<\eta\}
\longrightarrow \mathcal Y_\tau
\]
is \(C^1\).

The Banach implicit function theorem now applies at \((0,0)\), because
\(D_\varphi\mathcal F(0,0)\) is an isomorphism. Thus, for \(\|b\|\) small,
there is a unique \(\varphi_b\in\mathcal A_\tau^\R\), small in
\(\mathcal A_\tau\), such that \(\mathcal F(b,\varphi_b)=0\), and
\(b\mapsto\varphi_b\) is \(C^1\). The embedding
\(\mathcal A_\tau\hookrightarrow C^2\) gives the displayed \(C^2\) bound.

It remains to upgrade the coefficient identities \eqref{eq:disk-moment-system}
to the global identity \eqref{eq:disk-inverse-identity}. Define the entire
functions
\begin{equation}\label{eq:LHS-RHS-entire}
\mathcal L(w):=\int_{U_b}|P_b(z)|^2 e^{-\pi|z|^2}e^{\pi\overline z w}\,dA(z),
\qquad
\mathcal R(w):=\lambda Q_b(w),
\qquad w\in\C.
\end{equation}
The function \(\mathcal R\) is a polynomial of degree \(\le N\), and
\(\mathcal L\) is entire of exponential type \(\pi M\), with
\(M=\sup_{z\in U_b}|z|\), since
\(|e^{\pi\overline z w}|\le e^{\pi M|w|}\) on \(U_b\). Both \(\mathcal L\) and
\(\mathcal R\) admit absolutely convergent Taylor expansions at the origin,
\[
\mathcal L(w)=\sum_{j\ge0}\frac{\pi^j}{j!}
\Big(\!\int_{U_b}\!|P_b(z)|^2\overline z^{\,j}e^{-\pi|z|^2}dA(z)\!\Big)w^j,
\qquad
\mathcal R(w)=\sum_{j=0}^N\lambda q_j(b)\,w^j,
\]
the first by interchange of summation and integration (justified by
\(\sum_j\frac{(\pi M|w|)^j}{j!}<\infty\) and the boundedness of \(U_b\)).
The system \eqref{eq:disk-moment-system}, namely
\(\mathcal F_j(b,\varphi_b)=0\) for every \(j\ge0\), is exactly the
identification of the Taylor coefficients of \(\mathcal L\) and
\(\mathcal R\) at the origin, including the moments of degree \(j>N\) for
which \(q_j(b)=0\) and the integral on the left-hand side accordingly
vanishes. Hence \(\mathcal L=\mathcal R\) as formal power series, and by the
identity theorem for entire functions \(\mathcal L\equiv\mathcal R\) on
\(\C\), which is \eqref{eq:disk-inverse-identity}.

Equivalently, the same identity may be read off from the exterior
Cauchy-transform picture: restricting \eqref{eq:disk-inverse-identity} to
real \(w=t>0\), multiplying by \(e^{-\pi\zeta t}\) and integrating in
\(t\in(0,\infty)\) for real \(\zeta>M\) gives, by Fubini's theorem (applicable
since \(\Re(\overline z-\zeta)\le|z|-\zeta<0\) on \(U_b\)),
\begin{equation}\label{eq:cauchy-from-moments}
\frac1\pi\int_{U_b}\frac{|P_b(z)|^2 e^{-\pi|z|^2}}{\zeta-\overline z}\,dA(z)
=
\lambda\sum_{j=0}^N q_j(b)\,\frac{j!}{\pi^{j+1}\zeta^{j+1}},
\qquad \zeta\in(M,\infty).
\end{equation}
Both sides of \eqref{eq:cauchy-from-moments} are holomorphic in \(\zeta\) on
the connected open set \(\C\setminus\overline{U_b^*}\) (where
\(U_b^*=\{\overline z:z\in U_b\}\)), and they agree on the real ray
\((M,\infty)\), which has an accumulation point in this set. By the
identity theorem for holomorphic functions on a connected domain,
\eqref{eq:cauchy-from-moments} extends to all
\(\zeta\in\C\setminus\overline{U_b^*}\). Expanding the geometric series
\(1/(\zeta-\overline z)=\sum_{j\ge0}\overline z^{\,j}/\zeta^{\,j+1}\) at
\(\zeta=\infty\) and matching the resulting Laurent coefficients gives the
moment identities \(\frac{\pi^j}{j!}\int_{U_b}|P_b|^2\overline z^{\,j}e^{-\pi|z|^2}dA=\lambda q_j(b)\)
of \eqref{eq:disk-moment-system}, including the vanishing of the higher
moments \(\int_{U_b}|P_b|^2\overline z^{\,j}e^{-\pi|z|^2}dA=0\) for every
\(j>N\) (equivalently \(q_j(b)=0\) for \(j>N\)). The two formulations are
therefore equivalent, and either yields \eqref{eq:disk-inverse-identity}.
This proves the proposition.
\end{proof}

\begin{remark}
The proof above is the only place where analyticity of the boundary
parametrization is used. The constructed domains are therefore \(C^\infty\),
hence \(C^2\), which is the regularity needed later. The argument does not
assert a general local quadrature-domain theorem near an arbitrary reference
domain; it proves exactly the disk-based inverse statement used below.
\end{remark}

We isolate the ODE-theoretic ingredient needed for the unsymmetrization step
of the proof of Theorem~\ref{thm:perturb}, namely the claim that the only
entire function annihilated by \(P^\#(\partial_w/\pi)\) and growing no faster
than the Cauchy-transform of a bounded weighted measure on \(U\) is the zero
function, provided the zeros of the symbol \(P^\#\) escape to infinity.

\begin{lemma}[Step 3: rigidity of exponential-polynomial solutions]\label{lem:step3-rigidity}
Let $N\ge1$ and let $P(z)=1+\sum_{k=1}^N b_k z^k$ be a polynomial with
$\max_k|b_k|\le 1$. Let \(P^\#(z)=\overline{P(\overline z)}\) be its
coefficientwise conjugate, and consider the constant-coefficient differential
operator \(P^\#(\partial_w/\pi)\) defined as in \eqref{eq:Phash-operator}.
Suppose $H\colon\C\to\C$ is an entire function satisfying
\begin{equation}\label{eq:lemma-step3-eq}
P^\#(\partial_w/\pi)H(w)=0,
\qquad w\in\C,
\end{equation}
together with the growth bound
\begin{equation}\label{eq:lemma-step3-growth}
|H(w)|\le C(1+|w|)^N e^{\pi M|w|},
\qquad w\in\C,
\end{equation}
for some \(C>0\) and some \(M>0\). If, in addition, the zeros
$\zeta_1,\dots,\zeta_\ell$ of the characteristic polynomial
$\widetilde P^\#(\zeta):=P^\#(\zeta/\pi)$ on $\C$ all satisfy
\begin{equation}\label{eq:lemma-step3-zeros}
|\zeta_j|>\pi M,\qquad 1\le j\le\ell,
\end{equation}
then \(H\equiv 0\).
\end{lemma}

\begin{proof}
Let $m_1,\dots,m_\ell$ be the multiplicities of the zeros
$\zeta_1,\dots,\zeta_\ell$ of $\widetilde P^\#$, and write
\(\widetilde P^\#(\zeta)=\overline{b_N}\,\pi^{-N}\prod_{j=1}^\ell(\zeta-\zeta_j)^{m_j}\)
when $b_N\neq0$ (otherwise the leading coefficient is the highest non-zero
$\overline{b_k}$ divided by $\pi^k$, and the same argument applies). The
elementary structure theorem for entire solutions of constant-coefficient
linear ODEs of finite exponential type \cite[Ch.~XII]{Hormander} states that
every entire solution of \eqref{eq:lemma-step3-eq} with finite exponential
type takes the form
\begin{equation}\label{eq:lemma-step3-decomposition}
H(w)=\sum_{j=1}^\ell R_j(w)\,e^{\zeta_j w},
\end{equation}
where \(R_j\in\C[w]\) is a polynomial of degree at most \(m_j-1\). The
hypothesis \eqref{eq:lemma-step3-growth} ensures that \(H\) has exponential
type at most \(\pi M\), so \eqref{eq:lemma-step3-decomposition} indeed
applies.

Suppose, for a contradiction, that some $R_j$ does not vanish identically.
Re-indexing if necessary, choose $j_0$ with $R_{j_0}\not\equiv0$ and
$|\zeta_{j_0}|=\max\{|\zeta_j|\colon R_j\not\equiv 0\}$. After a small
generic rotation of the coordinate $w\mapsto e^{i\theta}w$ (which leaves both
\eqref{eq:lemma-step3-eq} and \eqref{eq:lemma-step3-growth} invariant up to
relabeling the $\zeta_j$), we may assume that for every \(j\) with
$j\ne j_0$ and $R_j\not\equiv 0$,
\begin{equation}\label{eq:lemma-step3-strict}
\Re(\zeta_j\,\overline{\zeta_{j_0}})<|\zeta_{j_0}|^2.
\end{equation}
This is automatic when $|\zeta_j|<|\zeta_{j_0}|$ from the Cauchy--Schwarz
inequality $\Re(\zeta_j\,\overline{\zeta_{j_0}})\le|\zeta_j||\zeta_{j_0}|$,
and ties at maximal modulus are broken by the rotation.

Along the half-line \(w_t:=t\,\overline{\zeta_{j_0}}\), \(t\ge 0\), the
exponential factor of the \(j_0\)th term has modulus
$|e^{\zeta_{j_0}w_t}|=e^{t|\zeta_{j_0}|^2}$, while every other term
satisfies
\(|R_j(w_t)e^{\zeta_j w_t}|\le C_j(1+t)^{m_j-1}e^{t\,\Re(\zeta_j\overline{\zeta_{j_0}})}\),
which by \eqref{eq:lemma-step3-strict} grows strictly slower than
$e^{t|\zeta_{j_0}|^2}$. Combining these and using $R_{j_0}\not\equiv 0$,
\begin{equation}\label{eq:lemma-step3-lower}
|H(w_t)|\ge c_0\,e^{t|\zeta_{j_0}|^2}
\qquad
\text{for all sufficiently large $t$,}
\end{equation}
with $c_0>0$ depending only on $R_{j_0}$ and the gap in
\eqref{eq:lemma-step3-strict}.

On the other hand, \eqref{eq:lemma-step3-growth} evaluated at $w_t$ gives
\begin{equation}\label{eq:lemma-step3-upper}
|H(w_t)|
\le
C(1+t|\zeta_{j_0}|)^N e^{\pi M\,t|\zeta_{j_0}|}.
\end{equation}
Comparing \eqref{eq:lemma-step3-lower} and \eqref{eq:lemma-step3-upper},
and absorbing the polynomial factor into an arbitrarily small exponential,
yields \(|\zeta_{j_0}|^2\le\pi M|\zeta_{j_0}|\), i.e.
\(|\zeta_{j_0}|\le\pi M\), contradicting \eqref{eq:lemma-step3-zeros}. Hence
all $R_j$ are identically zero, and \(H\equiv0\).
\end{proof}

\begin{remark}\label{rmk:step3-no-hidden-obstruction}
Lemma~\ref{lem:step3-rigidity} is the precise reason why no
finite-dimensional obstruction is left behind by the symmetrization
\eqref{eq:eigenfunction-general-new}~\(\Rightarrow\)~\eqref{eq:almost-reduced-new}:
the kernel of \(P^\#(\partial_w/\pi)\) on the full space of entire functions
is the finite-dimensional exponential-polynomial space spanned by
\(\{w^ie^{\zeta_jw}:0\le i<m_j,\ 1\le j\le\ell\}\), but the growth
constraint \eqref{eq:lemma-step3-growth} together with the non-vanishing of
\(P^\#\) on the closure of \(U\) (which forces \eqref{eq:lemma-step3-zeros})
collapses this kernel to $\{0\}$. The role of $P^\#$ being non-vanishing on
$U$ is therefore not cosmetic: it is what guarantees that the zeros of
$\widetilde P^\#$ are far from the origin, hence beyond the exponential-type
radius $\pi M$ that any candidate solution can have.
\end{remark}

\begin{proof}[Proof of Theorem \ref{thm:perturb}]
We again use the Bargmann transform. Let
\[
P(z)=1+\sum_{k=1}^N b_k z^k,
\]
where the coefficients $b_k$ are chosen so that the Bargmann transform of
\[
f_0=h_0+\sum_{k=1}^N a_k h_k
\]
is exactly $P$. We also write
\begin{equation}\label{eq:Phash-def}
P^\#(z):=\overline{P(\overline z)}
=1+\sum_{k=1}^N \overline{b_k}\,z^k
\end{equation}
for the \emph{coefficientwise conjugate polynomial}: the operation
\(P\mapsto P^\#\) leaves the monomials \(z^k\) intact and conjugates only the
scalar coefficients, so \(P^\#\) is again a polynomial of the same degree as
\(P\), holomorphic in \(z\), satisfying
\(P^\#(\overline z)=\overline{P(z)}\) for every \(z\in\C\).
Associated with \(P^\#\) is the holomorphic constant-coefficient differential
operator on the variable \(w\in\C\) defined by
\begin{equation}\label{eq:Phash-operator}
P^\#(\partial_w/\pi)
:=
\sum_{k=0}^N\overline{b_k}\,\Big(\frac{\partial_w}{\pi}\Big)^k,
\qquad b_0:=1,
\end{equation}
i.e. by substituting \(z=\partial_w/\pi\) into \eqref{eq:Phash-def}. Since
\(\partial_w=\frac12(\partial_x-i\partial_y)\) is the standard Wirtinger
derivative on entire functions of \(w=u+iv\), the operator
\(P^\#(\partial_w/\pi)\) acts on every entire function as the finite linear
combination of the iterated Wirtinger derivatives \((\partial_w/\pi)^k\),
\(0\le k\le N\); in particular it is a finite-order, holomorphic differential
operator. The key calculation is that
\(P^\#(\partial_w/\pi)\) acts as the multiplier \(P^\#(\overline z)\) on
each plane wave \(w\mapsto e^{\pi\overline z w}\):
\begin{equation}\label{eq:Phash-on-planewaves}
P^\#(\partial_w/\pi)\big[e^{\pi\overline z w}\big]
=
P^\#(\overline z)\,e^{\pi\overline z w}
=
\overline{P(z)}\,e^{\pi\overline z w},
\qquad z,w\in\C,
\end{equation}
because \((\partial_w/\pi)^k e^{\pi\overline z w}=\overline z^{\,k}e^{\pi\overline z w}\).

Term-by-term differentiation of an integrand of the form
\(P(z)e^{-\pi|z|^2}e^{\pi\overline z w}\) against \(z\in U\) is justified by
dominated convergence: \(U\) is bounded, \(P(z)e^{-\pi|z|^2}\) is bounded on
\(U\), and for each \(k\le N\) and each compact \(K\subset\C\) one has the
locally uniform estimate
\[
\Big|\Big(\frac{\partial_w}{\pi}\Big)^k\!\big[e^{\pi\overline z w}\big]\Big|
=
|\overline z|^k\,e^{\pi\Re(\overline z w)}
\le
M^k e^{\pi M|w|},
\qquad z\in U,\ w\in K,
\]
with \(M:=\sup_{z\in U}|z|\). Thus differentiation and integration may be
exchanged finitely many times, which is all that is required to apply
\(P^\#(\partial_w/\pi)\) term-by-term to the entire function defined by an
integral against \(z\in U\). Then the condition that $f_0$ be an eigenfunction
of $\mathcal{L}_U$ with eigenvalue $\lambda$ is equivalent to
\begin{equation}\label{eq:eigenfunction-general-new}
\int_U P(z)e^{-\pi|z|^2}e^{\pi \overline z w}\,dA(z)=\lambda P(w),
\qquad w\in \C.
\end{equation}
Set
\begin{equation}\label{eq:Q-def}
Q(w):=P^\#(\partial_w/\pi)P(w).
\end{equation}
By Proposition~\ref{prop:disk-inverse-tailored}, after decreasing
\(\varepsilon_0(\lambda,N)\) if necessary, there exists a real-analytic radial
graph \(U=U_b\), \(C^2\)-close to the Daubechies disk of radius
\(r_\lambda\), for which
\begin{equation}\label{eq:almost-reduced-new}
\int_U |P(z)|^2e^{-\pi|z|^2}e^{\pi\overline z w}\,dA(z)
=
\lambda Q(w),
\qquad w\in\C .
\end{equation}
We now recover the unsymmetrized identity
\eqref{eq:eigenfunction-general-new}. Define
\[
G(w):=\int_U P(z)e^{-\pi|z|^2}e^{\pi\overline z w}\,dA(z).
\]
The function \(G\) is entire, and differentiating under the integral is
justified by boundedness of \(U\). Since
\[
\Big(\frac{\partial_w}{\pi}\Big)^j e^{\pi\overline z w}
=\overline z^{\,j}e^{\pi\overline z w},
\]
we have
\[
P^\#(\partial_w/\pi)G(w)
=
\int_U |P(z)|^2e^{-\pi|z|^2}e^{\pi\overline z w}\,dA(z)
=
\lambda Q(w)
=
P^\#(\partial_w/\pi)(\lambda P)(w).
\]
Thus \(H:=G-\lambda P\) is an entire solution of
\[
P^\#(\partial_w/\pi)H=0.
\]

We now verify the hypotheses of Lemma~\ref{lem:step3-rigidity}. Setting
\(M:=\sup_{z\in U}|z|\), the integral defining \(G\) yields the standard
exponential-type bound
\begin{equation}\label{eq:G-growth-bound}
|G(w)|\le\Big(\sup_{z\in U}|P(z)|e^{-\pi|z|^2}\Big)\,|U|\,e^{\pi M|w|}
\le C(1+|w|)^N e^{\pi M|w|},
\qquad w\in\C,
\end{equation}
because \(|e^{\pi\overline z w}|\le e^{\pi|z||w|}\) for \(z\in U\). Since
\(\lambda P\) grows only polynomially, \(H=G-\lambda P\) inherits the same
exponential-type bound, with possibly larger constant: \(H\) satisfies
\eqref{eq:lemma-step3-growth}.

The location of the zeros of the characteristic polynomial \(\widetilde P^\#\)
is controlled by \(b\). Since \(U\subset V_2\) (with \(V_2\) defined in
Proposition~\ref{prop:disk-inverse-tailored}), and since
\(P^\#(z)=1+\sum_k\overline{b_k}z^k\) reduces to the constant \(1\) at
\(b=0\), the zeros of \(P^\#\) on \(\C\) escape to infinity uniformly as
\(b\to0\); a fortiori the zeros of \(\widetilde P^\#(\zeta)=P^\#(\zeta/\pi)\)
escape to infinity. After shrinking \(\varepsilon_0(\lambda,N)\) if
necessary, all \(\ell\) zeros \(\zeta_j\) of \(\widetilde P^\#\) satisfy
\[
|\zeta_j|>10\,M\ge\pi M,
\qquad 1\le j\le\ell,
\]
which is hypothesis \eqref{eq:lemma-step3-zeros}. Lemma~\ref{lem:step3-rigidity}
therefore implies $H\equiv 0$, which is precisely
\eqref{eq:eigenfunction-general-new}.

The Bargmann transform is unitary and intertwines the localization operator
with the integral operator in \eqref{eq:eigenfunction-general-new}. Therefore
the inverse Bargmann transform \(f_0\) of \(P\) satisfies
\[
\mathcal L_U f_0=\lambda f_0 .
\]
This proves the theorem.
\end{proof}

\iffalse
The proof proceeds in four steps. In Step~0 we apply the coefficientwise conjugate differential operator $P^\#(\partial_w/\pi)$ to \eqref{eq:eigenfunction-general-new} in order to symmetrize the integrand and reduce matters to a moment identity for $|P|^2 e^{-\pi|z|^2}$. In Step~1 we construct an auxiliary radially symmetric domain $U_0$ for which the corresponding eigenvalue is the constant function $\lambda$; this is achieved by perturbing a Daubechies disk via the Isakov free-boundary scheme of Proposition~\ref{prop:isakov-tailored}. In Step~2 we use a Laplace transform in the variable $w$ to convert the polynomial right-hand side $\lambda Q(w)$ into the rational right-hand side required by Proposition~\ref{prop:cherednichenko-tailored} and produce a domain $U$ satisfying the symmetrized moment identity. Finally, in Step~3 we invert the differentiation of Step~0 by analyzing the entire solutions of the constant-coefficient ODE $P^\#(\partial_w/\pi)H=0$ and ruling out nontrivial exponential solutions through a growth estimate.

\medskip
\noindent\textbf{Step 0: a first reduction.}
For a polynomial
\[
R(z)=\alpha_0+\alpha_1 z+\cdots+\alpha_n z^n,
\]
define the differential operator
\[
R(\partial_w/\pi):=\sum_{k=0}^n \alpha_k\Big(\frac{\partial_w}{\pi}\Big)^k.
\]
We apply $P^\#(\partial_w/\pi)$ to both sides of \eqref{eq:eigenfunction-general-new}. The right-hand side is the polynomial $\lambda P(w)$, on which the differential operator acts in the obvious way. For the left-hand side, write
\[
F(w):=\int_U P(z)e^{-\pi|z|^2}e^{\pi \overline z w}\,dA(z).
\]
Since $U$ is bounded, the integrand is jointly continuous on $\overline U\times \C$, and for every compact set $K\subset\C$ and every nonnegative integer $j$ one has the uniform bound
\[
\big|\partial_w^{\,j}\big(P(z)e^{-\pi|z|^2}e^{\pi\overline z w}\big)\big|
\le
\big(\pi\,\sup_{z\in U}|z|\big)^{j}\,\sup_{z\in U}\!\big(|P(z)|e^{-\pi|z|^2}\big)\,e^{\pi(\sup_{z\in U}|z|)\sup_{w\in K}|w|},
\]
valid for $(z,w)\in U\times K$. The right-hand side is integrable on the bounded set $U$, so dominated convergence justifies differentiating under the integral sign arbitrarily often, and $F$ is entire. Combining this with the identity
\[
\Big(\frac{\partial_w}{\pi}\Big)^j e^{\pi \overline z w}=\overline z^{\,j} e^{\pi \overline z w},
\]
we obtain
\[
P^\#(\partial_w/\pi)F(w)
=
\int_U P(z)\,\overline{P(z)}\,e^{-\pi|z|^2}e^{\pi\overline z w}\,dA(z)
=
\int_U |P(z)|^2 e^{-\pi|z|^2}e^{\pi\overline z w}\,dA(z).
\]
Hence \eqref{eq:eigenfunction-general-new} is equivalent to
\begin{equation}\label{eq:almost-reduced-new}
\int_U |P(z)|^2 e^{-\pi|z|^2} e^{\pi \overline z w}\,dA(z)=\lambda Q(w),
\qquad w\in\C,
\end{equation}
where
\[
Q(w):=P^\#(\partial_w/\pi)\big(P(w)\big).
\]
Strictly speaking, applying $P^\#(\partial_w/\pi)$ only gives the implication \eqref{eq:eigenfunction-general-new}~$\Rightarrow$~\eqref{eq:almost-reduced-new}; the converse implication will be recovered in Step~3 by an ODE-theoretic argument. Our first task is therefore to construct a domain $U$ satisfying \eqref{eq:almost-reduced-new}.

\medskip
\noindent\textbf{Step 1: construction of an auxiliary domain.}
We begin by constructing a $C^2$ domain $U_0$, close to a fixed disk, for which
\begin{equation}\label{eq:first-domain-new}
\int_{U_0} |P(z)|^2 e^{-\pi|z|^2} e^{\pi \overline z w}\,dA(z)=\lambda,
\qquad w\in\C.
\end{equation}
Note that \eqref{eq:first-domain-new} is the special case of \eqref{eq:almost-reduced-new} in which the right-hand side is the constant $\lambda$; it will serve as the unperturbed reference for Step~2.

Let $\psi\in C_c^\infty(\C)$ be a nonnegative radial function such that
\[
\supp \psi \subset D_2(0)\setminus D_1(0),
\qquad
\int_\C \psi\,dA=1,
\]
and set
\[
\psi_\delta(z):=\delta^{-2}\psi(z/\delta),\qquad 0<\delta<1.
\]
Then $\psi_\delta$ has total mass one and $\supp\psi_\delta\subset D_{2\delta}(0)\setminus D_\delta(0)$. Since $\psi_\delta$ is radial, expanding $e^{\pi\overline z w}=\sum_{j\ge 0}(\pi\overline z w)^j/j!$ and integrating in polar coordinates kills every angular harmonic with $j\ge 1$, so that
\[
\int_{\C} \psi_\delta(z)e^{-\pi|z|^2}e^{\pi \overline z w}\,dA(z)=c_\delta,
\qquad w\in\C,
\]
where
\[
c_\delta:=\int_{\C}\psi_\delta(z)e^{-\pi|z|^2}\,dA(z)>0.
\]
Note that $c_\delta\to 1$ as $\delta\to 0$. Hence \eqref{eq:first-domain-new} is equivalent to
\begin{equation}\label{eq:second-domain-new}
\int_{U_0} |P(z)|^2 e^{-\pi|z|^2} e^{\pi \overline z w}\,dA(z)
=
\frac{\lambda}{c_\delta}
\int_{\C}\psi_\delta(z)e^{-\pi|z|^2}e^{\pi \overline z w}\,dA(z).
\end{equation}

We shall construct a compactly supported distributional solution of the Poisson equation
\begin{equation}\label{eq:inverse-first-new}
\Delta u_0
=
|P(z)|^2 e^{-\pi|z|^2}\mathbf{1}_{U_0}
-\frac{\lambda}{c_\delta}\psi_\delta(z)e^{-\pi|z|^2}.
\end{equation}
Once such a solution has been constructed, the moment identity \eqref{eq:first-domain-new} follows by applying the same Fourier-analytic divisibility argument used in the proof of Theorem~\ref{thm:eig-loc-U} to the corresponding compactly supported distribution
\[
\omega_0:=
|P(z)|^2e^{-\pi|z|^2}\mathbf 1_{U_0}
-\frac{\lambda}{c_\delta}\psi_\delta(z)e^{-\pi|z|^2}.
\]
Indeed, \eqref{eq:inverse-first-new} says exactly that $\Delta u_0=\omega_0$, and since $u_0$ is compactly supported, the vanishing of $\widehat{\omega_0}$ on the two characteristic lines of the Bargmann transform is equivalent to \eqref{eq:first-domain-new}.

\smallskip
\emph{Unperturbed reference.} We first treat the unperturbed case $P\equiv 1$. By the theorem of Daubechies \cite{Daubechies}, there exists a disk $D_0$ centered at the origin such that
\[
\int_{D_0} e^{-\pi|z|^2}e^{\pi \overline z w}\,dA(z)=\lambda,
\qquad w\in\C.
\]
Equivalently, $\Delta u_*=e^{-\pi|z|^2}\mathbf 1_{D_0}-\lambda\,\delta_0$ admits a compactly supported radial solution $u_*$. This solution is unique among compactly supported distributional solutions: by radial symmetry, the Poisson equation reduces to the ODE
\[
\frac{1}{r}\big(r u_*'(r)\big)'=e^{-\pi r^2}\mathbf 1_{(0,r_0)}(r),\qquad r>0,
\]
together with the asymptotic condition $u_*'(r)\sim -\lambda/(2\pi r)$ as $r\to 0^+$ enforced by the distributional source \(-\lambda\delta_0\). Integrating twice with the natural boundary condition that $u_*$ vanish for $r\ge r_0$ gives a unique radial profile. Explicitly,
\[
u_*(r)=\frac{\lambda}{2\pi}\log\frac{r_0}{r}-\int_r^{r_0}\frac{1}{s}\int_0^{s}\rho\,e^{-\pi\rho^2}\,d\rho\,ds,
\qquad 0<r\le r_0,
\]
and $u_*\equiv 0$ for $r\ge r_0$, where $r_0$ is the radius of $D_0$. From this explicit formula one verifies directly:
\begin{itemize}
\item $\Delta u_*=e^{-\pi|z|^2}>0$ on $D_0\setminus\{0\}$;
\item $u_*=0$ on $\partial D_0$;
\item the logarithmic singularity created by the mass \(-\lambda\delta_0\) has positive coefficient $\lambda/(2\pi)$ in \(\log(r_0/r)\), so $u_*(r)\to+\infty$ as $r\to 0^+$.
\end{itemize}
Equivalently, if \(A(s):=\int_0^s \rho e^{-\pi\rho^2}\,d\rho\), then the
Daubechies identity gives \(A(r_0)=\lambda/(2\pi)\), and the displayed formula
can be rewritten as
\[
u_*(r)=\int_r^{r_0}\frac{A(r_0)-A(s)}{s}\,ds>0,
\qquad 0<r<r_0.
\]
In particular \(u_*>0\) on every compact subset of \(D_0\setminus\{0\}\).

For $\delta$ small enough that $\supp\psi_\delta\subset D_0$, the same argument yields a compactly supported radial solution $\widetilde u_0$ of
\begin{equation}\label{eq:unperturbed-free-boundary-new}
\Delta \widetilde u_0
=
e^{-\pi|z|^2}\mathbf{1}_{D_0}
-\frac{\lambda}{c_\delta}\psi_\delta(z)e^{-\pi|z|^2},
\end{equation}
and $\widetilde u_0$ is uniquely determined among compactly supported radial distributional solutions by the same ODE reduction (the smooth radial source on the right-hand side replaces the Dirac mass, and the requirements of compact support and continuity through the origin pin down the integration constants).

We claim that, for $\delta$ sufficiently small,
\[
\widetilde u_0>0 \qquad \text{in } D_0.
\]
To see this, we first verify that
\begin{equation}\label{eq:psi-delta-to-mass}
\frac{\lambda}{c_\delta}\,\psi_\delta(z)\,e^{-\pi|z|^2}\xrightarrow[\delta\to 0]{}\lambda\,\delta_0
\qquad \text{in }\mathcal D'(\C).
\end{equation}
Indeed, for any $\varphi\in C_c^\infty(\C)$,
\[
\Big\langle \frac{\lambda}{c_\delta}\psi_\delta\, e^{-\pi|\cdot|^2},\varphi\Big\rangle
=
\frac{\lambda}{c_\delta}\int_\C \psi_\delta(z)\,e^{-\pi|z|^2}\varphi(z)\,dA(z).
\]
Decompose $\varphi(z)e^{-\pi|z|^2}=\varphi(0)+r(z)$ with $r$ continuous and $r(0)=0$; then
\[
\frac{1}{c_\delta}\int_\C \psi_\delta(z)\,e^{-\pi|z|^2}\varphi(z)\,dA(z)
=
\varphi(0)\cdot\frac{1}{c_\delta}\int_\C\psi_\delta(z)e^{-\pi|z|^2}\,dA(z)
+
\frac{1}{c_\delta}\int_\C\psi_\delta(z)\,r(z)\,dA(z).
\]
The first summand equals $\varphi(0)$. The second is bounded in absolute value by $c_\delta^{-1}\,\|\psi_\delta\|_{L^1}\,\sup_{|z|\le 2\delta}|r(z)|=c_\delta^{-1}\sup_{|z|\le 2\delta}|r(z)|\to 0$ as $\delta\to 0$, since $r$ is continuous, $r(0)=0$, and $c_\delta\to 1$. Multiplying by $\lambda$ proves \eqref{eq:psi-delta-to-mass}.

In particular, $\Delta\widetilde u_0\to\Delta u_*$ in $\mathcal D'(\C)$. To upgrade this to local uniform convergence, we use the Newtonian-potential representations of $\widetilde u_0$ and $u_*$ together with the harmonic correction enforcing compact support: writing
\[
\widetilde u_0(z)
=
\frac{1}{2\pi}\int_\C \log\frac{1}{|z-\zeta|}\Big(-e^{-\pi|\zeta|^2}\mathbf 1_{D_0}(\zeta)+\frac{\lambda}{c_\delta}\psi_\delta(\zeta)e^{-\pi|\zeta|^2}\Big)dA(\zeta)+h_\delta(z),
\]
\[
u_*(z)
=
\frac{1}{2\pi}\int_\C \log\frac{1}{|z-\zeta|}\Big(-e^{-\pi|\zeta|^2}\mathbf 1_{D_0}(\zeta)\Big)dA(\zeta)+\frac{\lambda}{2\pi}\log\frac{1}{|z|}+h_*(z),
\]
where $h_\delta,h_*$ are harmonic functions chosen so as to enforce compact support, the contribution from $-e^{-\pi|\zeta|^2}\mathbf 1_{D_0}$ cancels in the difference and one is left with
\[
\widetilde u_0(z)-u_*(z)
=
\frac{\lambda}{2\pi}\Big(\int_\C \log\frac{1}{|z-\zeta|}\frac{1}{c_\delta}\psi_\delta(\zeta)e^{-\pi|\zeta|^2}\,dA(\zeta)-\log\frac{1}{|z|}\Big)+\big(h_\delta(z)-h_*(z)\big).
\]
By dominated convergence applied to the logarithmic potential with kernel uniformly bounded on compact subsets of $\overline D_0\setminus\{0\}$, together with \eqref{eq:psi-delta-to-mass}, the displayed difference tends to zero locally uniformly on $\overline D_0\setminus\{0\}$. Combined with interior elliptic regularity for $\Delta(\widetilde u_0-u_*)$, which is harmonic on any compact subset of $D_0\setminus\{0\}$ once $\delta$ is so small that $\supp\psi_\delta$ is contained in any prescribed neighborhood of the origin, this gives
\[
\widetilde u_0\xrightarrow[\delta\to 0]{}u_* \qquad \text{locally uniformly on $\overline D_0\setminus\{0\}$}.
\]
Since $u_*$ is strictly positive on every compact subset of $D_0\setminus\{0\}$ and continuous up to $\partial D_0$ where it vanishes, this implies $\widetilde u_0>0$ on every fixed compact subset of $D_0\setminus\{0\}$ for all $\delta$ small enough; the strict inequality $u_*>0$ at interior points provides the room needed to absorb the local-uniform error. On the small annulus $\delta\le|z|\le 2\delta$, where $\psi_\delta$ is supported, the explicit profile near the origin shows that $\widetilde u_0$ is comparable to a truncated logarithm $-\log(|z|/r_0)$ shifted by an $O(1)$ smooth correction, and is therefore strictly positive. We conclude that $\widetilde u_0>0$ throughout $D_0$ for $\delta>0$ small.

\smallskip
\emph{Perturbed problem via Isakov.}
We now invoke the perturbative free-boundary theorem of Isakov \cite{Isakov}, given by Proposition~\ref{prop:isakov-tailored}, verifying its hypotheses carefully. Set
\[
f_0(z):=e^{-\pi|z|^2}\Big(1-\frac{\lambda}{c_\delta}\psi_\delta(z)\Big).
\]
Since $\psi_\delta\in C_c^\infty(\C)$, we have $f_0\in C^\infty(\R^2)$, hence in particular $f_0\in C^5(\R^2)$. If $\delta$ is chosen so small that $\supp\psi_\delta\subset D_0$, then on $\partial D_0$,
\[
f_0(z)=e^{-\pi|z|^2}>0.
\]
Then \eqref{eq:unperturbed-free-boundary-new} becomes
\[
\Delta \widetilde u_0=f_0\,\mathbf{1}_{D_0}
\qquad \text{in }\C,
\]
with \(\widetilde u_0=0\) on \(\C\setminus D_0\) and \(\widetilde u_0>0\) in \(D_0\), by the positivity established above. Since \(\widetilde u_0\) was constructed as the compactly supported radial solution of \eqref{eq:unperturbed-free-boundary-new}, it is \(C^1\) across \(\partial D_0\); equivalently \(\widetilde u_0=\partial_\nu \widetilde u_0=0\) on \(\partial D_0\). Thus all hypotheses of Proposition~\ref{prop:isakov-tailored} are satisfied with \(\Omega_0=D_0\).

For the perturbed source, define
\[
f(z):=e^{-\pi|z|^2}\Big(|P(z)|^2-\frac{\lambda}{c_\delta}\psi_\delta(z)\Big).
\]
Since $P$ is a polynomial, $f\in C^\infty(\R^2)$ as well. Moreover, for any fixed disk $\Omega'\supset D_0$, the $C^5(\Omega')$-norm of $f-f_0$ is controlled by the coefficients of $P-1$: writing
\[
|P(z)|^2-1=\sum_{k=1}^N b_k z^k+\sum_{k=1}^N \overline{b_k}\,\overline z^{\,k}+\sum_{k,\ell=1}^N b_k\overline{b_\ell}\,z^k\overline z^{\,\ell},
\]
which is a polynomial in $z,\overline z$ that vanishes when all $b_k=0$ and depends polynomially on the coefficients, a standard estimate yields
\[
\|f-f_0\|_{C^5(\Omega')}\le C(\Omega',N,\lambda)\max_{1\le k\le N}|b_k|.
\]
Hence, if the coefficients $b_k$ are sufficiently small, we may arrange that
\[
\|f-f_0\|_{C^5(\Omega')}<\varepsilon_1,
\]
where $\varepsilon_1$ is the smallness threshold furnished by Proposition~\ref{prop:isakov-tailored}. Applying that proposition, we obtain a \(C^2\) domain \(U_0\), represented as a small normal deformation of \(D_0\), and a compactly supported function \(u_0\) such that
\[
\Delta u_0=f\,\mathbf{1}_{U_0}
\qquad \text{in }\C,
\]
with \(u_0=0\) on \(\C\setminus U_0\), \(\partial_\nu u_0=0\) on \(\partial U_0\), and \(u_0>0\) in \(U_0\).
If the perturbation is sufficiently small, then $U_0$ is a small $C^2$ deformation of $D_0$, and since $\supp\psi_\delta\subset D_0$ we still have
\[
\supp\psi_\delta\subset U_0.
\]
Therefore
\[
\Delta u_0
=
|P(z)|^2e^{-\pi|z|^2}\mathbf{1}_{U_0}
-\frac{\lambda}{c_\delta}\psi_\delta(z)e^{-\pi|z|^2},
\]
so \eqref{eq:inverse-first-new} holds. This proves \eqref{eq:first-domain-new}, and completes Step~1: we have produced a $C^2$ domain $U_0$, lying in a small $C^2$-neighborhood of the Daubechies disk $D_0$, on which the symmetrized moment identity holds with constant right-hand side $\lambda$. We will use $U_0$ in Step~2 as the unperturbed reference for the application of Proposition~\ref{prop:cherednichenko-tailored}.

\medskip
\noindent\textbf{Step 2: construction of a domain satisfying \eqref{eq:almost-reduced-new}.}
With the auxiliary domain $U_0$ at our disposal, we now seek a $C^2$ domain $U$ for which the right-hand side of the moment identity is the polynomial $\lambda Q(w)$ rather than the constant $\lambda$. The crucial observation is that integrating the parameter $w$ along the positive real axis against an exponential weight converts \eqref{eq:almost-reduced-new} into an explicit identity for a weighted Cauchy transform; this is exactly the form to which Proposition~\ref{prop:cherednichenko-tailored} of Cherednichenko \cite{Cherednichenko} applies. The reformulation we are about to perform is closely related to the quadrature-domain framework developed in \cite{Gustafsson,Gustafsson-Shahgholian,Karp,Shahgholian}.

Restrict \eqref{eq:almost-reduced-new} to real values $w=t>0$, multiply by $e^{-\pi \zeta t}$, and integrate in $t>0$. If
\[
\zeta>\sup_{z\in U}|z|,
\]
the integrand of the resulting double integral is absolutely integrable: indeed, for real $\zeta>\sup_{z\in U}|z|$,
\[
\Re(\overline z-\zeta)\le |z|-\zeta<0
\qquad \text{for every } z\in U,
\]
so
\[
\int_U \int_0^\infty \big|e^{\pi(\overline z-\zeta)t}\big|\,dt\,|P(z)|^2 e^{-\pi|z|^2}\,dA(z)
\le
\int_U \frac{|P(z)|^2 e^{-\pi|z|^2}}{\pi(\zeta-|z|)}\,dA(z)<\infty,
\]
because $U$ is bounded. Fubini's theorem therefore applies and yields
\[
\int_U \left(\int_0^\infty e^{\pi(\overline z-\zeta)t}\,dt\right)|P(z)|^2e^{-\pi|z|^2}\,dA(z)
=
\lambda\int_0^\infty e^{-\pi\zeta t}Q(t)\,dt.
\]
Writing
\[
Q(w)=\alpha_0+\alpha_1 w+\cdots+\alpha_n w^n,
\]
we obtain
\[
\frac{1}{\pi}\int_U \frac{|P(z)|^2e^{-\pi|z|^2}}{\zeta-\overline z}\,dA(z)
=
\lambda \sum_{k=0}^n \alpha_k \int_0^\infty e^{-\pi \zeta t} t^k\,dt.
\]
Since
\[
\int_0^\infty e^{-\pi \zeta t} t^k\,dt
=
\frac{\Gamma(k+1)}{(\pi \zeta)^{k+1}},
\]
it follows that
\[
\frac{1}{\pi}\int_U \frac{|P(z)|^2e^{-\pi|z|^2}}{\zeta-\overline z}\,dA(z)
=
\lambda \sum_{k=0}^n \alpha_k \frac{\Gamma(k+1)}{\pi^{k+1}\zeta^{k+1}}.
\]

To put this into the form of a weighted Cauchy transform, write
\[
E^\ast:=\{\overline z:z\in E\}
\]
for reflection across the real axis, and let
\[
V:=U^\ast,
\qquad
\mu(w):=|P(\overline w)|^2 e^{-\pi|w|^2}.
\]
The change of variables $w=\overline z$ yields
\[
\int_U \frac{|P(z)|^2e^{-\pi|z|^2}}{\zeta-\overline z}\,dA(z)
=
\int_V \frac{\mu(w)}{\zeta-w}\,dA(w)
=
g_V(\zeta).
\]
Hence
\begin{equation}\label{eq:analytic-construct-new}
g_V(\zeta)
=
\lambda \sum_{k=0}^n \alpha_k \frac{\Gamma(k+1)}{\pi^k \zeta^{k+1}},
\qquad \zeta\in \R_+, \ \zeta\gg 1.
\end{equation}
Since both sides are holomorphic for large $\zeta$, analytic continuation gives the desired identity on $\C\setminus\overline V$.

Conversely, suppose a domain $V$ satisfies \eqref{eq:analytic-construct-new}. Expanding $1/(\zeta-w)=\sum_{k\ge 0}w^k/\zeta^{k+1}$ for $|\zeta|$ large and matching coefficients yields the moments
\[
\int_V w^k\mu(w)\,dA(w)=\lambda\,\alpha_k\,\frac{\Gamma(k+1)}{\pi^k},
\qquad 0\le k\le n,
\]
together with the vanishing of the corresponding moments for $k>n$ (since $Q$ has degree $n$). These moment identities recover \eqref{eq:almost-reduced-new} for the reflected domain \(U=V^\ast\) via the Laplace-transform calculation read in reverse, in conjunction with the identity theorem for entire functions. Thus the reflection \(V=U^\ast\) is merely a bookkeeping device that places the kernel in the form $1/(\zeta-w)$ required by the weighted Cauchy transform formalism.

Let $g$ denote the rational function on the right-hand side of \eqref{eq:analytic-construct-new}. We now verify the hypotheses of Proposition~\ref{prop:cherednichenko-tailored}. Since
\[
Q(w)=P^\#(\partial_w/\pi)\big(P(w)\big),
\]
and $P(0)=1$, the constant term of $Q$ is
\[
Q(0)=1+\sum_{k=1}^N \frac{k!}{\pi^k}|b_k|^2.
\]
In particular,
\[
Q(0)>0,
\]
so the leading term in the expansion of $g$ at infinity is
\[
\frac{\lambda Q(0)}{z},
\]
and consequently
\[
g(z)\to 0,
\qquad
z\,g(z)\to \lambda Q(0)>0
\qquad \text{as } z\to\infty.
\]
Thus \eqref{eq:appendix-g-hypotheses} holds, with $c_0=\lambda Q(0)$.

We next verify the rational-data smallness condition in Proposition~\ref{prop:cherednichenko-tailored}. The weighted Cauchy transform formulation is written for the reflected domain, so set
\[
V_0:=U_0^\ast.
\]
Since \(U_0\) was obtained in Step~1 as a small perturbation of the Daubechies disk \(D_0\), the reflected domain \(V_0\) is also a small perturbation of \(D_0\). We choose concentric disks \(V_1\) and \(V_2\), centered at the center of \(D_0\), with radii, say, \((1-10^{-3})\) and \((1+10^{-3})\) times the radius of \(D_0\), after shrinking the perturbation so that
\[
\overline{V_1}\subset V_0,\qquad \overline{V_0}\subset V_2 .
\]
By Step~1,
\[
\int_{U_0}|P(z)|^2e^{-\pi|z|^2}e^{\pi \overline zw}\,dA(z)=\lambda,
\qquad w\in \C.
\]
Repeating the Laplace-transform argument above, with the constant $\lambda$ in place of the polynomial $\lambda Q(w)$, gives
\[
g_{V_0}(z)=\frac{\lambda}{z}
\qquad \text{for } z\in \C\setminus \overline{V_0},
\]
where \(g_{V_0}\) denotes the weighted Cauchy transform with weight \(\mu(w)=|P(\overline w)|^2e^{-\pi|w|^2}\). Consequently, on \(\partial V_2\),
\[
g(z)-g_{V_0}(z)
=
\lambda \sum_{k=1}^n \alpha_k \frac{\Gamma(k+1)}{\pi^k z^{k+1}}
+
\lambda\,(Q(0)-1)\,\frac1z .
\]
The right-hand side belongs to the finite-dimensional real rational space
\[
\mathcal R_n:=\operatorname{span}_{\R}
\{z^{-1},z^{-2},iz^{-2},\ldots,z^{-(n+1)},iz^{-(n+1)}\}.
\]
Every element of \(\mathcal R_n\) has real \(z^{-1}\)-coefficient. It is
holomorphic on \(\C\setminus\overline{V_1}\), has only a pole at the origin,
and vanishes at infinity. Since all norms on the finite-dimensional
space \(\mathcal R_n\) are equivalent, there is a constant
\(C(V_1,V_2,\lambda,n)\) such that
\[
\big\|(g-g_{V_0})|_{\partial V_2}\big\|_{C^2(\partial V_2)}
\le
C(V_1,V_2,\lambda,n)
\Big(|Q(0)-1|+\sum_{k=1}^n |\alpha_k|\Big).
\]
Each coefficient \(\alpha_k\) of \(Q-1\) depends polynomially on the
coefficients of \(P-1\) and their conjugates, with no constant term, and
therefore tends to zero with \(\max_{1\le j\le N}|b_j|\). Hence, by choosing
the coefficients sufficiently small, we may ensure that
\[
(g-g_{V_0})\in\mathcal R_n,
\qquad
\big\|(g-g_{V_0})|_{\partial V_2}\big\|_{C^2(\partial V_2)}<\varepsilon_0,
\]
where \(\varepsilon_0\) is the smallness threshold provided by Proposition~\ref{prop:cherednichenko-tailored}. Applying that proposition, we obtain a \(C^2\) domain \(\widetilde U\) such that
\[
g_{\widetilde U}(z)=g(z),
\qquad z\in \C\setminus \overline{\widetilde U}.
\]
Setting \(U:=\widetilde U^\ast\), the converse part of the Laplace-transform argument gives \eqref{eq:almost-reduced-new}.

Finally, since \(\widetilde U\subset V_2\), the domain \(U=\widetilde U^\ast\) is contained in \(V_2^\ast=V_2\), and \(V_2\) is fixed once \(\lambda\) is fixed. By shrinking the admissible size of the coefficients of \(P-1\) once more we may ensure that
\[
P^\#(z)\neq 0
\qquad \text{for all } z\in V_2.
\]
In particular,
\[
P^\#(z)\neq 0
\qquad \text{for all } z\in U.
\]
This non-vanishing of $P^\#$ on $U$ will play no direct role until Step~3, where it is needed to push the zeros of the characteristic polynomial $\widetilde P^\#(\zeta)=P^\#(\zeta/\pi)$ to large modulus.

\medskip
\noindent\textbf{Step 3: recovery of the original eigenvalue equation.}
It remains to upgrade the symmetrized identity \eqref{eq:almost-reduced-new} to the original equation \eqref{eq:eigenfunction-general-new}. Define
\[
G(w):=\int_U P(z)e^{-\pi|z|^2}e^{\pi \overline z w}\,dA(z),
\qquad w\in\C.
\]
Then $G$ is entire (the integrand is entire in $w$, and the dominated-convergence justification of differentiation under the integral sign carried out in Step~0 applies verbatim with $P$ in place of $|P|^2$). By construction, \eqref{eq:almost-reduced-new} holds, so applying $P^\#(\partial_w/\pi)$ to $G$ yields, by the same argument as in Step~0,
\[
P^\#(\partial_w/\pi)G(w)=\int_U|P(z)|^2 e^{-\pi|z|^2}e^{\pi\overline z w}\,dA(z)=\lambda Q(w)=P^\#(\partial_w/\pi)(\lambda P)(w),
\qquad w\in \C.
\]
Hence
\[
H:=G-\lambda P
\]
is an entire function satisfying the constant-coefficient linear differential equation
\[
P^\#(\partial_w/\pi)H=0.
\]

Introduce the characteristic polynomial of this operator,
\[
\widetilde P^\#(\zeta):=P^\#(\zeta/\pi),
\]
and let $\zeta_1,\dots,\zeta_\ell$ be its zeros, with multiplicities $m_1,\dots,m_\ell$. By the structure theorem for entire solutions of constant-coefficient linear differential equations of finite exponential type (see e.g.\ \cite[Ch.\ XII]{Hormander}), every entire solution of
\[
P^\#(\partial_w/\pi)H=0
\]
of at most exponential type is of the form
\begin{equation}\label{eq:exp-decomposition-step3}
H(w)=\sum_{j=1}^\ell R_j(w)\,e^{\zeta_j w},
\end{equation}
where each $R_j$ is a polynomial of degree strictly less than $m_j$.

We will reach a contradiction by showing that all $R_j$ must vanish. First, since $P^\#$ has no zeros on $V_2$, and the zeros of $\widetilde P^\#(\zeta)=P^\#(\zeta/\pi)$ are obtained by scaling those of $P^\#$ by the factor $\pi$, by choosing the coefficients of $P-1$ sufficiently small we may guarantee that all zeros of $\widetilde P^\#$ stay at distance at least $10\sup_{z\in U}|z|$ from the origin:
\[
|\zeta_j|>10\sup_{z\in U}|z|=:10M,
\qquad 1\le j\le \ell.
\]
(Indeed, the roots of $\widetilde P^\#$ depend continuously on its coefficients, and at $b=0$ the polynomial $\widetilde P^\#$ degenerates to the constant $1$, which has no roots; for small $b$ the roots therefore lie in any prescribed neighborhood of infinity.)

Next, we establish an upper bound on $H$ from its definition. Since $P$ is a fixed polynomial of degree $\le N$ and $U$ is bounded with $M=\sup_{z\in U}|z|$,
\[
|G(w)|
\le
\int_U |P(z)|e^{-\pi|z|^2}e^{\pi |z||w|}\,dA(z)
\le
\Big(\sup_{z\in U}|P(z)|e^{-\pi|z|^2}\Big)\,|U|\,e^{\pi M|w|}
\le
C(1+|w|)^N e^{\pi M|w|},
\qquad w\in\C,
\]
where the first inequality uses $|e^{\pi\overline z w}|\le e^{\pi|z||w|}$ and in the last inequality we have absorbed $\sup_{z\in U}|P(z)|\le C_0(1+M)^N$ and the area $|U|$ into the constant $C$. Since every polynomial is dominated by $e^{\varepsilon |w|}$ for arbitrary $\varepsilon>0$, it follows that for every $\varepsilon>0$ there exists $C_\varepsilon>0$ such that
\[
|G(w)|\le C_\varepsilon e^{(\pi+\varepsilon)M|w|},
\qquad w\in \C.
\]
As $\lambda P(w)$ grows only polynomially in $|w|$, the same estimate holds for $H$:
\[
|H(w)|\le C'_\varepsilon e^{(\pi+\varepsilon)M|w|}.
\]
In particular, $H$ has finite exponential type at most $\pi M$, which justifies the use of \eqref{eq:exp-decomposition-step3}.

Suppose, for contradiction, that some $R_j$ is not identically zero. By a small generic rotation of the coordinate $w$ if necessary, we may assume that the zeros $\zeta_j$ are arranged so that there is a unique $\zeta_1$ realizing
\[
|\zeta_1|=\max_j |\zeta_j|
\qquad\text{and}\qquad
R_1\not\equiv 0,
\]
and that, for every other zero $\zeta_j$ with $j\ne 1$,
\[
\Re(\zeta_j\overline{\zeta_1})<|\zeta_1|^2.
\]
(The strict inequality $\Re(\zeta_j\overline{\zeta_1})\le |\zeta_j||\zeta_1|<|\zeta_1|^2$ holds automatically when $|\zeta_j|<|\zeta_1|$; ties at maximal modulus are broken by the rotation.) Then, along the ray
\[
w=t\,\overline{\zeta_1},
\qquad t>0,
\]
the term $R_1(w)e^{\zeta_1 w}$ has modulus $|R_1(t\overline{\zeta_1})|e^{t|\zeta_1|^2}$, while every other term $R_j(w)e^{\zeta_j w}$ has modulus bounded by a polynomial in $t$ times $e^{t\,\Re(\zeta_j\overline{\zeta_1})}$, with $\Re(\zeta_j\overline{\zeta_1})<|\zeta_1|^2$. Hence the dominant exponential term yields
\[
|H(t\overline{\zeta_1})|\ge c\,e^{t|\zeta_1|^2}
\]
for all sufficiently large $t$, with $c>0$.

Evaluating the upper bound on $H$ at $w=t\overline{\zeta_1}$ gives
\[
|H(t\overline{\zeta_1})|
\le
C'_\varepsilon e^{(\pi+\varepsilon)Mt|\zeta_1|}.
\]
Comparing with the lower bound $|H(t\overline{\zeta_1})|\ge c\,e^{t|\zeta_1|^2}$ and letting $t\to\infty$, we obtain
\[
|\zeta_1|^2\le (\pi+\varepsilon)M|\zeta_1|,
\qquad\text{i.e.,}\qquad
|\zeta_1|\le (\pi+\varepsilon)M.
\]
Since $\varepsilon>0$ is arbitrary, this gives
\[
|\zeta_1|\le \pi M,
\]
contradicting our choice $|\zeta_1|>10M$ (recall that $\pi<10$). Hence every $R_j$ vanishes identically, so
\[
H\equiv 0.
\]
Equivalently,
\[
G=\lambda P,
\]
which is exactly \eqref{eq:eigenfunction-general-new}. This proves that $P$, and therefore $f_0$, is an eigenfunction of the localization operator associated with the domain $U$ with eigenvalue $\lambda$.
\end{proof}
\fi

% ==============================================================
% Section 4
% ==============================================================
\section{A Classical Rigidity Corollary}\label{sec:classical-rigidity}

The localization operator $\mathcal{L}_U$ on the Bargmann--Fock space is
diagonalized by Hermite functions whenever $U$ is a disk centered at the
origin: this is an immediate consequence of the rotational symmetry of $U$
together with the radial structure of the Hermite weight. The converse
question, raised by Abreu and D\"orfler, asks whether eigenfunction status of
even a \emph{single} Hermite function already forces the underlying domain to
be a disk. The next theorem records a positive answer in the simply connected
$C^2$ regime.

For convenience we recall the statement of Theorem~\ref{thm:eig-loc-U} from the introduction:
\emph{let $U\subset \C$ be a simply connected bounded domain with $C^2$
boundary, and let $h_k$ be a Hermite function ($k\ge 0$). If $h_k$ is an
eigenfunction of the localization operator $\mathcal{L}_U$, then $U$ is a
disk.} (The original Abreu--D\"orfler theorem~\cite{Abreu-Doerfler} dispenses
with the boundary regularity assumption; we will recover this stronger form
in the second proof below.)

We will give two proofs of Theorem~\ref{thm:eig-loc-U}. The first, presented
below, leverages the Fourier-analytic machinery developed earlier in the
paper, reducing the question to a free boundary problem for a logarithmic
potential. The second, in the next subsection, proceeds along more
classical complex-analytic lines and, as we shall see, dispenses with the
$C^2$ boundary assumption altogether.

Before launching into the arguments, we record, in a form convenient for
later reference, the two analytic inputs underlying both proofs. The first is
the precise variant of the Paley--Wiener theorem we use to interpret the
moment identity \eqref{eq:fourier-delta} below as the vanishing of an entire
function on a complex hypersurface; the second is the
Ehrenpreis--Malgrange-type divisibility lemma that converts this vanishing
into a Poisson equation with prescribed right-hand side. Both statements are
quoted in the exact form in which they will be applied; we have stated them
here, rather than referring back to Section~\ref{sec:fourier-to-fboundary},
in order to keep the regularity hypotheses on $\omega$ visible to the
reader following Section~\ref{sec:classical-rigidity} in isolation.

\begin{theorem}[Paley--Wiener--Schwartz, planar form]\label{thm:PW-section4}
Let $\omega\in\mathcal D'(\R^2)$ be a compactly supported distribution of
order $\le N$ with $\operatorname{supp}\omega\subset \overline{B(0,R)}$.
Then its Fourier transform $\widehat\omega$ extends to an entire function
on $\C^2$ satisfying the Paley--Wiener bound
\begin{equation}\label{eq:PW-bound-section4}
|\widehat\omega(\xi,\eta)|\le C\,(1+|\xi|+|\eta|)^N\,e^{R(|\Im\xi|+|\Im\eta|)},
\qquad (\xi,\eta)\in\C^2,
\end{equation}
for some $C>0$. Conversely, any entire function on $\C^2$ obeying such a
bound is the Fourier--Laplace transform of a unique compactly supported
distribution of order $\le N$ supported in $\overline{B(0,R)}$. If, in
addition, the bound holds with $N=0$ and $\widehat\omega|_{\R^2}\in
L^2(\R^2)$, then $\omega\in L^2(\R^2)$.
\end{theorem}

This is the classical Paley--Wiener--Schwartz theorem; we use the form
stated in \cite[Theorem~7.3.1]{Hormander}, see also \cite[\S 9.3]{Folland}.
The $L^2$ refinement in the last sentence is a direct consequence of
Plancherel applied to the entire function $\widehat\omega|_{\R^2}$, combined
with uniqueness of the Fourier inversion. In our application, $\omega$ is
a finite signed Radon measure of compact support, so the Paley--Wiener
bound \eqref{eq:PW-bound-section4} holds with $N=0$ and $C=\|\omega\|_{\mathrm{TV}}$.

\begin{lemma}[Ehrenpreis--Malgrange-type divisibility along $\xi^2+\eta^2$]\label{lem:EM-section4}
Let $\omega$ be a compactly supported singular-continuous free measure on
$\R^2$, in the sense of Section~\ref{sec:fourier-to-fboundary}: a finite
signed Radon measure of the form
\[
\omega=f\,dx\,dy+\sum_{j=1}^N c_j\,\delta_{p_j},
\]
with $f\in L^\infty(\R^2)$ compactly supported, $c_j\in\R\setminus\{0\}$,
and $p_j\in\R^2$ distinct. Suppose moreover that the Fourier--Laplace
transform $\widehat\omega$, an entire function on $\C^2$ of exponential
type by Theorem~\ref{thm:PW-section4}, vanishes on the complex
characteristic variety of the planar Laplacian,
\begin{equation}\label{eq:complex-char-variety}
\Sigma:=\{(\xi,\eta)\in\C^2:\xi^2+\eta^2=0\}=\{\eta=i\xi\}\cup\{\eta=-i\xi\}.
\end{equation}
Then $\widehat\omega$ is divisible, in the ring of entire functions of
exponential type on $\C^2$, by the polynomial $\xi^2+\eta^2$, with quotient
$F\in L^2(\R^2)\cap L^1(\R^2)$ on the real locus. Equivalently, there
exists a unique compactly supported $u\in L^2(\R^2)$, continuous off a
finite set $P\subset\R^2$, such that
\[
\Delta u=\omega\qquad\text{in }\mathcal D'(\R^2),
\]
namely the logarithmic potential
\[
u(z)=\frac{1}{2\pi}\int_{\R^2}\log|z-z'|\,d\omega(z').
\]
\end{lemma}

This is exactly Lemma~\ref{lemma:fourier-to-fboundary}, restated to
emphasize its character as a planar-Laplacian instance of the
Ehrenpreis--Malgrange division theorem (entire functions of exponential
type can be divided by a polynomial within the same Paley--Wiener class;
see \cite[Ch.~7]{Hormander}). We do \emph{not} invoke the general
Ehrenpreis--Malgrange machinery: a self-contained proof, exploiting the
homogeneity decomposition of $\widehat\omega$ and positivity of
$\xi^2+\eta^2$ on $\R^2$, is given in Section~\ref{sec:fourier-to-fboundary}.
The hypotheses on $\omega$ -- in particular, the absence of a continuous
singular part and the finiteness of the point part -- are essential for
the conclusion that $u$ is continuous off a finite set; without them the
quotient $F$ remains in $L^2$, but $u$ need only be quasi-continuous in
the logarithmic-potential sense.

\subsection{First proof}

\begin{proof}[First proof under the additional assumption that $\partial U$ is $C^2$]
The Bargmann transform of the Hermite function $h_k$ is a nonzero scalar
multiple of the monomial $z^k$. Thus, arguing as in
\cite{Abreu-Doerfler}, the assumption that $h_k$ is an eigenfunction of
the localization operator $\mathcal{L}_U$ is equivalent to the identity
\[
\int_U z^k\, e^{\pi w\overline{z}}\, e^{-\pi |z|^2}\,dA(z)
=\lambda\, w^k,
\qquad w\in \C.
\]
Both sides are entire functions of $w$. Since $U$ is bounded and the
weight $e^{-\pi|z|^2}$ is integrable, the dominated convergence
theorem applied to difference quotients shows that for each
multi-index $\alpha$ in $w$, the partial derivative
$\partial_w^\alpha$ commutes with the integral on the left-hand side;
the absolute integrability of $|z|^{|\alpha|+k} e^{\pi
|w||z|}e^{-\pi|z|^2}\mathbf{1}_U$ for fixed $w$ provides the required
local majorant. Iterating $k$ times, we obtain
\[
\frac{1}{(\pi)^k}\,\partial_w^k
\Big[\int_U z^k\, e^{\pi w\overline{z}}\, e^{-\pi |z|^2}\,dA(z)\Big]
=\int_U z^k\overline{z}^k\, e^{\pi w\overline{z}}\, e^{-\pi |z|^2}\,dA(z)
=\int_U |z|^{2k}\, e^{\pi w\overline{z}}\, e^{-\pi |z|^2}\,dA(z).
\]
On the right-hand side, $\partial_w^k(\lambda w^k)=\lambda\,k!$.
Setting $c(\lambda,k):=\lambda\,k!/\pi^k\in\R$ and using
$\mathbf{1}_U$ to extend the integration to $\C$, we obtain
\begin{equation}\label{eq:fourier-delta}
\int_{\C} \mathbf{1}_U(z)\,|z|^{2k}\,e^{-\pi |z|^2}\, e^{\pi \overline{z}\, w}\,dA(z)
=c(\lambda,k),
\end{equation}
for some real constant $c(\lambda,k)$.

Consider now the compactly supported distribution
\begin{equation}\label{eq:eta-def-section4}
\eta:=|z|^{2k}e^{-\pi |z|^2}\mathbf{1}_U(z)-c(\lambda,k)\,\delta_0,
\end{equation}
which is a finite signed Radon measure with compact support in any closed
ball containing $U\cup\{0\}$. Its absolutely continuous part has bounded
density $|z|^{2k}e^{-\pi|z|^2}\mathbf{1}_U\in L^\infty_c(\R^2)$, and its
pure-point part is the single Dirac mass $-c(\lambda,k)\delta_0$. Hence
$\eta$ is a singular-continuous free measure in the sense of
Section~\ref{sec:fourier-to-fboundary}, and we may apply
Theorem~\ref{thm:PW-section4} and Lemma~\ref{lem:EM-section4} to it.

Writing $z=x+iy$ and $w\in\C$, we rewrite the exponential factor in
\eqref{eq:fourier-delta} as
\[
e^{\pi \overline{z}\, w}=e^{\pi(x-iy)w}=e^{\pi xw}\, e^{-i\pi yw}.
\]
The function class entering the Paley--Wiener step is now unambiguous:
$\eta\in\mathcal E'(\R^2)$ is a compactly supported distribution of order
zero, in fact a finite signed Radon measure. Theorem~\ref{thm:PW-section4}
therefore applies and produces an entire function $\widehat\eta$ on
$\C^2$ obeying the Paley--Wiener bound \eqref{eq:PW-bound-section4} with
$N=0$ and $R$ equal to the radius of any closed disk containing
$U\cup\{0\}$; in particular $\widehat\eta|_{\R^2}$ is bounded by
$\|\eta\|_{\mathrm{TV}}$ and extends to $\C^2$ as an entire function of
exponential type. With the convention
$\widehat g(\xi,\eta):=\int_{\R^2}g(x,y)\,e^{-2\pi i(x\xi+y\eta)}\,dx\,dy$
used in Section~\ref{sec:fourier-to-fboundary}, a comparison of exponents
shows that for $w\in\C$ the kernel $e^{\pi xw}\,e^{-i\pi yw}$ coincides
with $e^{-2\pi i(x\xi+y\eta)}$ on the complex line
$(\xi,\eta)=\bigl(\tfrac{iw}{2},-\tfrac{w}{2}\bigr)$, $w\in\C$, whose
defining equation is $\eta=i\xi$. Reparametrizing this line in the form
$(\xi,\eta)=(w,-iw)$, identity \eqref{eq:fourier-delta}---which says that
the integral on the left-hand side, viewed as an entire function of
$w\in\C$, is identically equal to the constant $c(\lambda,k)$, so that
the Fourier--Laplace transform of the modified distribution $\eta$
vanishes there---becomes
\begin{equation}\label{eq:eta-vanishes-on-line}
\widehat{\eta}(w,-iw)=0,\qquad \forall w\in\C.
\end{equation}
Symmetrically, since $c(\lambda,k)\in\R$ and the absolutely continuous
density of $\eta$ is real-valued, complex conjugation of
\eqref{eq:fourier-delta} after the substitution $w\mapsto\overline w$,
combined with the same Paley--Wiener identification applied to the
conjugate complex line $\eta=-i\xi$, yields
\begin{equation}\label{eq:eta-vanishes-on-conjugate-line}
\widehat{\eta}(w,iw)=0,\qquad \forall w\in\C.
\end{equation}
Identities \eqref{eq:eta-vanishes-on-line} and
\eqref{eq:eta-vanishes-on-conjugate-line} together assert that
$\widehat\eta$ vanishes on the complex characteristic variety
$\Sigma=\{\xi^2+\eta^2=0\}$ of the Laplacian, see
\eqref{eq:complex-char-variety}.

Applying the divisibility lemma (Lemma~\ref{lem:EM-section4}, equivalently
Lemma~\ref{lemma:fourier-to-fboundary}) to the distribution $\eta$, we
obtain a unique compactly supported function $u\in L^2(\R^2)$, continuous
off a finite set, such that
\begin{equation}\label{eq:potato}
\Delta u=f(|z|)\,\mathbf{1}_U-c(\lambda,k)\,\delta_0
\end{equation}
in the sense of distributions on $\R^2$, where
\[
f(|z|)=|z|^{2k}e^{-\pi |z|^2}.
\]
Writing $u$ explicitly as a logarithmic potential of the right-hand side
of \eqref{eq:potato}, we conclude that $U$ satisfies
\begin{equation}\label{eq:potential}
\int_U \log|z-z'|\, f(|z'|)\,dA(z')
=c(\lambda,k)\,\log|z|,
\qquad z\in \C\setminus (U\cup\{0\}).
\end{equation}

\medskip

\noindent\emph{Reduction to $0\notin U$.}\quad
Suppose first that $0\in U$, and pick $\rho>0$ small enough that
$\overline{D_\rho(0)}\subset U$. The radial measure
$f(|z|)\,\mathbf{1}_{D_\rho(0)}\,dA$ generates a logarithmic potential
\[
\Phi_\rho(z):=\int_{D_\rho(0)} \log|z-z'|\, f(|z'|)\,dA(z'),
\]
which, by Newton's shell theorem in the plane, is radial; outside
$D_\rho(0)$ it equals
\[
\Phi_\rho(z)=\Big(\int_{D_\rho(0)}f(|z'|)\,dA(z')\Big)\log|z|+\text{const}.
\]
Subtracting $D_\rho(0)$ from $U$ thus preserves identity
\eqref{eq:potential} on the complement of the new set, at the cost of
replacing the constant $c(\lambda,k)$ by another real constant. Hence we
may assume $0\notin U$. Moreover, $0\notin \partial U$, because the
left-hand side of \eqref{eq:potential} stays bounded as $z\to 0$, whereas
the right-hand side diverges to $-\infty$ if $c(\lambda,k)\neq 0$.
The case $c(\lambda,k)=0$ is incompatible with the integral being
positive at $z=0$, so it is also excluded.

\medskip

Define the logarithmic potential
\[
\Psi_U(z):=\int_U \log|z-z'|\, f(|z'|)\,dA(z').
\]
We claim that $\Psi_U$ is \emph{locally radial}: for each $z_0\in U$,
there exists a ball $B(z_0)\subset U$ such that $\Psi_U$ agrees on
$B(z_0)$ with a radial function of $|z|$.

To prove this, let $B\supset U$ be a ball centered at $0$, and define
\[
\Phi(z):=\int_B \log|z-z'|\, f(|z'|)\,dA(z').
\]
Since both $B$ and the weight $f(|z'|)$ are radial, $\Phi$ is radial; in
particular, $\nabla \Phi(z)$ is parallel to $z$ for every $z\in \R^2$.
Furthermore, on the interior of $B$,
\[
\Delta \Phi(z)=2\pi\, f(|z|).
\]

Now $\Psi_U$ satisfies \eqref{eq:potential} on the unbounded
component of $\C\setminus(U\cup\{0\})$. Differentiating that identity
in $z$ for $z\not\in \overline{U}\cup\{0\}$ gives
\[
\nabla \Psi_U(z)=c(\lambda,k)\,\frac{z}{|z|^2},
\qquad z\in \C\setminus(\overline{U}\cup\{0\}).
\]
The classical jump relations for the single-layer / volume logarithmic
potential of an $L^\infty$ density on a $C^2$ domain (see, e.g.,
\cite{Gustafsson}) imply that $\nabla\Psi_U$ extends continuously
across $\partial U$ from \emph{both} sides, and the boundary values
agree. Hence on $\partial U$,
\[
\nabla \Psi_U(z)=c(\lambda,k)\,\frac{z}{|z|^2},
\qquad z\in \partial U,
\]
which is, in particular, parallel to $z$. Therefore the difference
\[
v:=\Psi_U-\Phi
\]
satisfies $\Delta v=2\pi f(|z|)\mathbf 1_U-2\pi f(|z|)\mathbf 1_B
=-2\pi f(|z|)\mathbf 1_{B\setminus U}$ in the sense of distributions
on $B$, so $v$ is harmonic in $U$ and in $B\setminus\overline{U}$, and
on $\partial U$ satisfies
\[
\nabla v(x)\ \text{is parallel to}\ x.
\]
The next lemma converts this gradient condition into local radiality of $v$.

\begin{lemma}\label{lem:harmonic-locally-radial}
Let $V\subset \R^2$ be a domain, and suppose there exists a harmonic
function $v\in C^1(\overline{V})$ such that $\nabla v(x)$ is parallel to
$x$ for every $x\in \partial V$. Then $\nabla v(x)$ is parallel to $x$
for every $x\in V$. In particular, on every circle
$S=\{|x|=r\}\subset \R^2$, the function $v$ is constant on every
connected component of $S\cap V$.
\end{lemma}

\begin{proof}
Consider the angular derivative
\[
g(x):=x_1\,\partial_2 v(x)-x_2\,\partial_1 v(x)
=\partial_\theta v(x),
\]
which equals the tangential derivative of $v$ along the circle
through $x$ centered at $0$. A direct computation gives
\[
\Delta g
=x_1\,\Delta(\partial_2 v)-x_2\,\Delta(\partial_1 v)
+2\big(\partial_1\partial_2 v-\partial_2\partial_1 v\big)
=0,
\]
since $v$ is harmonic and partial derivatives of harmonic functions
are themselves harmonic. By the boundary hypothesis, $\nabla v(x)$ is
parallel to $x$ on $\partial V$, so $g\equiv 0$ on $\partial V$.
Together with $g\in C(\overline{V})$, the maximum principle for
harmonic functions yields $g\equiv 0$ in $V$. Hence $\nabla v(x)$ is
parallel to $x$ at every interior point as well, and on any circle
$S=\{|x|=r\}$ the tangential derivative $\partial_\theta v$ vanishes
on $S\cap V$. Therefore $v$ is constant on each connected component
of $S\cap V$.
\end{proof}

Applying Lemma~\ref{lem:harmonic-locally-radial} on small subdomains
$V\subset U$ where $v$ is harmonic and $C^1$ up to the boundary, we
conclude that $v$ is locally radial on $U$. Since $\Phi$ is globally
radial, so is $\Psi_U=\Phi+v$, locally on $U$.

\begin{lemma}\label{lem:local-to-global-radial}
Let \(W\subset \R^2\setminus\{0\}\) be a connected open set, and let
\(u\in C(W)\). Suppose that for every \(x\in W\) there exist \(r_x>0\)
and a scalar function \(\phi_x\) such that
\[
u(y)=\phi_x(|y|)
\qquad\text{for all } y\in B(x,r_x)\subset W.
\]
Then there exists a single scalar function \(\phi\), defined on the
interval
\[
I_W:=\{|y|:y\in W\},
\]
such that
\[
u(y)=\phi(|y|)
\qquad\text{for all } y\in W.
\]
\end{lemma}

\begin{proof}
For each \(x\in W\), choose a ball \(B_x:=B(x,r_x)\subset W\) on which
\(u\) is a radial function of \(|y|\), with profile $\phi_x$. We
first note that whenever $B_x\cap B_{x'}\ne\emptyset$,
\[
\phi_x(|y|)=u(y)=\phi_{x'}(|y|)
\qquad\text{for every }y\in B_x\cap B_{x'}.
\]
Since $B_x\cap B_{x'}$ is open and avoids the origin, the map $y\mapsto
|y|$ is a submersion on it, so the set of radii it attains contains a
nontrivial open interval on which $\phi_x$ and $\phi_{x'}$ agree.

Fix $x_*\in W$ and define a candidate function $\phi$ on $I_W$ as
follows. Given $y\in W$, connect $x_*$ to $y$ by a polygonal path in
$W$, which exists because $W$ is open and connected. Its image is
compact, hence covered by finitely many balls
$B_{x_1},\dots,B_{x_m}$ with successive nonempty overlaps and
$x_1=x_*$, $y\in B_{x_m}$. The previous paragraph implies $\phi_{x_j}$
and $\phi_{x_{j+1}}$ agree on a nontrivial interval of radii, so by
analytic-continuation along the chain we may unambiguously set
$\phi(|y|):=\phi_{x_m}(|y|)$. Independence of the chain follows from
the same overlap argument: any two chains from $x_*$ to $y$ can be
interleaved into a single chain, and successive overlaps force the
final value to coincide. Hence $\phi$ is well-defined on $I_W$ and
satisfies $u(y)=\phi(|y|)$ for every $y\in W$.
\end{proof}

We now conclude the proof. Since \(U\) is simply connected and bounded,
it suffices to rule out the possibility that \(U\) fails to be a disk
centered at the origin. Suppose therefore, for contradiction, that
\(U\) is not such a disk. Then \(\partial U\) is not a single circle
centered at $0$, so there exist radii \(0<r_1<r_2\) such that the
annulus
\[
A:=\{z\in \R^2:\ r_1<|z|<r_2\}
\]
meets both $U$ and $\partial U$ in a non-trivial way; in particular,
we may pick a point
\[
x_0\in \partial U\cap A.
\]
Let
\[
\widetilde{U}_0:=U\cap B_r(x_0),
\]
for $r>0$ sufficiently small that $\widetilde U_0$ is connected and
$0\notin \overline{\widetilde U_0}$. Since \(\Psi_U\) is locally radial
in \(U\), Lemma~\ref{lem:local-to-global-radial} applied to
$W=\widetilde U_0$ yields a scalar function \(\psi_U\) on
$I_{\widetilde U_0}$ such that
\[
\Psi_U(x)=\psi_U(|x|),\qquad x\in \widetilde{U}_0.
\]

On the other hand, by \eqref{eq:potential} and standard regularity for
logarithmic potentials with $C^2$ boundary, the gradient $\nabla\Psi_U$
extends continuously across $\partial U$, and for $z\in\partial U$ we
have
\[
\nabla \Psi_U(z)=c\,\frac{z}{|z|^2}
\]
for the same constant $c=c(\lambda,k)\in \R$ as in
\eqref{eq:potential}. Combining this with $\Psi_U(x)=\psi_U(|x|)$ on
$\widetilde U_0$ and passing to $z\in \partial \widetilde U_0\cap
\partial U$ from the inside,
\[
\psi_U'(|z|)\,\frac{z}{|z|}=c\,\frac{z}{|z|^2},
\qquad z\in \partial U\cap \overline{B_r(x_0)},
\]
hence
\[
\psi_U'(r)=\frac{c}{r}
\]
for every $r$ in the set of radii attained by points of
$\partial U\cap \overline{B_r(x_0)}$. Since $\partial U$ is a $C^2$
curve, we may, after shrinking $r$ once more, parametrize
$\partial U\cap B_r(x_0)$ by a $C^2$ embedding
$\gamma:(-\varepsilon,\varepsilon)\to \R^2$ with $\gamma(0)=x_0$.
Because $\partial U$ is not a circle centered at $0$, the function
$t\mapsto |\gamma(t)|$ is not constant on any subinterval; by choosing
$x_0\in\partial U\cap A$ at which $|\gamma|$ is non-constant on every
neighbourhood (such points form an open dense subset of $\partial
U\cap A$ when $\partial U$ is not a circle centered at $0$), we may
further assume that the derivative
$\frac{d}{dt}|\gamma(t)|\big|_{t=0}\ne 0$. Hence by the open-mapping
property for $C^1$ functions of one variable with nonzero derivative,
the image of $|\gamma|$ contains a nontrivial open interval
\(I\subset (r_1,r_2)\).
Therefore
\[
\psi_U'(r)=\frac{c}{r}
\qquad\text{for every }r\in I,
\]
whence
\[
\psi_U(r)=c\log r+C_0,
\qquad r\in I,
\]
for some real constant $C_0$.

Let
\[
W_I:=\{x\in \widetilde U_0:\ |x|\in I\}.
\]
Then \(W_I\) is a nonempty open subset of \(U\), and on \(W_I\),
\[
\Psi_U(x)=c\log|x|+C_0.
\]
Hence
\[
\Delta \Psi_U\equiv 0 \qquad \text{in } W_I,
\]
since the origin is excluded from $\widetilde U_0$. On the other hand,
distribution-theoretically,
\[
\Delta \Psi_U=2\pi\, f(|z|)\,\mathbf{1}_U-c\,\delta_0
\]
in $\C$, so on the open set $W_I\subset U\setminus\{0\}$ we have
\[
\Delta \Psi_U=2\pi\, f(|z|)=2\pi\, |z|^{2k}e^{-\pi |z|^2}.
\]
The right-hand side vanishes only at the origin, which lies outside
$W_I$, while the left-hand side just shown to be $0$ on $W_I$ — a
contradiction. This rules out the case in which $U$ is not a disk
centered at $0$, completing the proof.
\end{proof}

\begin{remark}
The structure of the argument above is in the spirit of
\cite{Aharonov-Schiffer-Zalcman}, where logarithmic-potential
rigidity statements were derived from Schiffer-style overdetermined
boundary conditions, and of the quadrature-domain literature
\cite{Gustafsson,Gustafsson-Shahgholian,Shahgholian,Karp}, in which
the regularity of free boundaries and the matching of an interior
potential with the Newtonian potential of a radially symmetric
density play a central role. Concretely, identity
\eqref{eq:potential} can be read as saying that $U$ is a
\emph{generalized quadrature domain} for the radial weight $f(|z|)$,
with quadrature data concentrated at the origin. From this perspective,
the conclusion that $U$ is a disk is a (sharp) instance of the
classical principle that radial quadrature data forces a radial
domain. The role of the $C^2$ boundary regularity is concentrated
in the boundary extension of $\nabla\Psi_U$ and in the
open-interval-of-radii argument that ends the proof; both can be
relaxed at the cost of additional technicalities, which we will
revisit in the second proof.
\end{remark}

\subsection{Second proof}\label{subsec:second-proof}

The first proof established Theorem~\ref{thm:eig-loc-U} under the auxiliary
hypothesis that $\partial U$ is of class $C^2$. The argument we now give does
not use any boundary regularity at all; it therefore strengthens
Theorem~\ref{thm:eig-loc-U} to a form which agrees with the original
Abreu--D\"orfler statement \cite{Abreu-Doerfler}.

\begin{theorem}[Strengthened rigidity, no boundary regularity]\label{thm:eig-loc-U-strong}
Let $U\subset\C$ be an open, bounded, simply connected set, with no
hypothesis on the regularity of $\partial U$ beyond what is implicit in
``simply connected''. Let $h_k$ ($k\ge 0$) be a Hermite function. If $h_k$
is an eigenfunction of the localization operator $\mathcal L_U$, then $U$
is a disk.
\end{theorem}

The hypotheses on $U$ in Theorem~\ref{thm:eig-loc-U-strong} reduce to: $U$
is open, bounded, connected, and its complement in $\overline\C$ is
connected (i.e.\ $U$ is simply connected). In particular, $\partial U$ is
allowed to be merely a closed set with $|U|<\infty$; no assumption of
positive Lebesgue density, finite perimeter, or differentiability is
imposed. The proof below shows that the boundary identity $\nabla u\equiv
0$ on $\partial U$, which in the first proof was extracted via boundary
regularity of the logarithmic potential and jump relations on a $C^2$
curve, now follows directly from the explicit formula for $u$ as a
convolution with $\log|\cdot|$, together with the local integrability of
the kernel $z/|z|^2$. Since $u$ vanishes identically on the open exterior
$\R^2\setminus\overline U$, continuity of $u$ and of $\nabla u$ on
$\R^2\setminus\{0\}$ forces both to vanish on $\partial U$, regardless of
how rough the boundary is. The remainder of the proof proceeds via a
holomorphic auxiliary function on $U$ whose boundary values are forced to
be real, and the maximum principle on $U$ (which is well-defined even for
rough boundaries) then yields the conclusion. We summarize this in
Remark~\ref{rem:bookkeeping-no-regularity} below; the precise place where
the $C^2$ assumption is bypassed is at the derivation
of~\eqref{eq:dz-zero-second-proof}.

\begin{proof}[Proof of Theorem~\ref{thm:eig-loc-U-strong} (second proof of Theorem~\ref{thm:eig-loc-U})]
We keep the notation
\[
f(|z|):=|z|^{2k}e^{-\pi|z|^2},
\qquad
c:=c(\lambda,k),
\]
from the first proof. The argument below is in the spirit of the holomorphic
moving-plane/quadrature-domain methods of \cite{Aharonov-Schiffer-Zalcman,Gustafsson,Karp},
in that we exploit the existence of a holomorphic function on \(U\) with prescribed real
boundary values to recover full radial symmetry. Starting again from \eqref{eq:fourier-delta} and
applying Lemma~\ref{lemma:fourier-to-fboundary}, we obtain a compactly supported distributional
solution of
\begin{equation}\label{eq:potato-second-proof}
\Delta u=f(|z|)\mathbf{1}_U-c\,\delta_0
\qquad\text{in }\mathcal D'(\R^2).
\end{equation}
Using the planar fundamental solution \((2\pi)^{-1}\log|z|\), we may write
\begin{equation}\label{eq:explicit-u-second-proof}
u(z):=\frac{1}{2\pi}\int_U \log|z-z'|\,f(|z'|)\,dA(z')-\frac{c}{2\pi}\log|z|.
\end{equation}
By \eqref{eq:potential}, this function satisfies
\[
u\equiv 0
\qquad\text{on }\R^2\setminus (U\cup\{0\}).
\]

We first show that \(0\in U\). Since the singular term in
\eqref{eq:potato-second-proof} is \(-c\delta_0\), one must have \(0\in\overline U\). If
\(0\in\partial U\), then there exist exterior points \(z_j\to 0\) with
\[
u(z_j)=0.
\]
On the other hand, the first term in \eqref{eq:explicit-u-second-proof} remains finite as
\(z\to 0\), because \(f\in L^\infty(U)\) and \(\log|z-z'|\) is locally integrable, whereas the
second term equals \(-\frac{c}{2\pi}\log|z|\to +\infty\). This contradicts the exterior identity
\(u(z_j)=0\). Hence \(0\notin\partial U\), and therefore
\[
0\in U.
\]

We next show directly from \eqref{eq:explicit-u-second-proof} that
\[
u\in C^1(\R^2\setminus\{0\}).
\]
Indeed, if we write
\[
g(z):=f(|z|)\mathbf{1}_U(z)\in L^\infty_c(\R^2),
\]
then
\[
u(z)=\frac{1}{2\pi}\int_{\R^2}\log|z-z'|\,g(z')\,dA(z')-\frac{c}{2\pi}\log|z|.
\]
Since the kernel \(K(\zeta)=\zeta/|\zeta|^2\) belongs to \(L^1_{\mathrm{loc}}(\R^2)\) (note that
\(|K(\zeta)|=|\zeta|^{-1}\), which is integrable on bounded subsets of \(\R^2\)), the translation
map \(z\mapsto K(z-\cdot)\) is continuous in \(L^1_{\mathrm{loc}}(\R^2)\) by the standard density
argument: continuous compactly supported functions are dense in \(L^1_{\mathrm{loc}}\), and for
each such function continuity of translations is immediate. Together with \(g\in L^\infty_c(\R^2)\),
this implies that the convolution
\[
(K\ast g)(z):=\int_{\R^2}\frac{z-z'}{|z-z'|^2}\,g(z')\,dA(z')
\]
is a continuous function of \(z\in\R^2\). On the other hand, a direct distributional computation
shows
\[
\nabla\!\left(\int_{\R^2}\log|z-z'|\,g(z')\,dA(z')\right)=K\ast g
\qquad\text{in }\mathcal D'(\R^2),
\]
so the gradient of the convolution \(\log|\cdot|\ast g\) coincides almost everywhere with the
continuous function \(K\ast g\). It follows that \(\log|\cdot|\ast g\) is of class \(C^1\) on
\(\R^2\), and hence \(u\in C^1(\R^2\setminus\{0\})\) (the only obstruction to global \(C^1\)
regularity being the explicit logarithmic singularity at the origin coming from the
\(-\frac{c}{2\pi}\log|z|\) term).

Now let \(z_*\in\partial U\). Since \(0\in U\), the boundary point \(z_*\) is bounded away from
\(0\); fix a small ball \(B_\rho(z_*)\Subset \R^2\setminus\{0\}\). On the exterior open set
\(\R^2\setminus \overline U\) we have \(u\equiv 0\), so \(\nabla u\equiv 0\) there. As
\(\partial U\) is contained in the boundary of the exterior, every \(z_*\in\partial U\) is a limit
of exterior points \(z_j\to z_*\) along which \(\nabla u(z_j)=0\). By the \(C^1\) regularity of
\(u\) on \(B_\rho(z_*)\), the gradient \(\nabla u\) is continuous at \(z_*\), and we conclude
\(\nabla u(z_*)=0\). In complex notation \(\partial_z=\tfrac{1}{2}(\partial_x-i\partial_y)\), this
yields
\begin{equation}\label{eq:dz-zero-second-proof}
\partial_z u=0
\qquad\text{on }\partial U.
\end{equation}
We emphasize that \emph{no} boundary regularity of \(\partial U\) has been used at this point, in
contrast with the first proof: the vanishing of \(\nabla u\) on \(\partial U\) is forced solely by
exterior continuity, which in turn is supplied by the explicit
formula~\eqref{eq:explicit-u-second-proof}.

Define
\[
\Phi(t):=\int_0^t s^k e^{-\pi s}\,ds
\qquad\text{and}\qquad
\Psi(z):=\frac{\Phi(|z|^2)}{z},
\quad z\neq 0.
\]
A direct computation, using \(|z|^2=z\bar z\) and \(\partial_{\bar z}(z\bar z)=z\), gives, for
\(z\neq 0\),
\[
\partial_{\bar z}\Psi(z)
=
\frac{\Phi'(|z|^2)\,\partial_{\bar z}(|z|^2)}{z}
=
\frac{|z|^{2k}e^{-\pi|z|^2}\cdot z}{z}
=
|z|^{2k}e^{-\pi|z|^2}
=
f(|z|).
\]
At the origin, the apparent singularity \(1/z\) is integrable against the vanishing factor
\(\Phi(|z|^2)\): since \(k\ge 0\) and \(\Phi(t)=\int_0^t s^k e^{-\pi s}\,ds=O(t^{k+1})\) as
\(t\to 0\), we have \(\Psi(z)=O(|z|^{2k+1})\) near \(z=0\), so \(\Psi\in L^\infty_{\mathrm{loc}}\)
(in fact continuous, with \(\Psi(0)=0\)). Consequently, the identity
\(\partial_{\bar z}\Psi=f(|z|)\) persists in the distributional sense on all of \(U\); no extra
\(\delta_0\) contribution arises at the origin.

Recall also the standard identity
\[
\partial_{\bar z}\!\left(\frac{1}{z}\right)=\pi\delta_0
\qquad\text{in }\mathcal D'(\R^2),
\]
which is equivalent to the fact that \(\frac{1}{\pi z}\) is a fundamental solution of
\(\partial_{\bar z}\); equivalently, \(\Delta\log|z|=4\partial_z\partial_{\bar z}\log|z|=2\pi\delta_0\)
combined with \(\partial_z\log|z|=\frac{1}{2z}\). Now set
\[
F(z):=4\partial_z u(z)-\Psi(z)+\frac{c}{\pi z},
\qquad z\in U\setminus\{0\}.
\]
Using \(\Delta=4\partial_z\partial_{\bar z}\), the relation \(\partial_{\bar z}(1/z)=\pi\delta_0\),
the computation of \(\partial_{\bar z}\Psi\) above, and \eqref{eq:potato-second-proof}, we
compute, in \(\mathcal D'(U)\),
\[
\partial_{\bar z}F
=
4\partial_z\partial_{\bar z}u-\partial_{\bar z}\Psi+\frac{c}{\pi}\partial_{\bar z}\!\left(\frac{1}{z}\right)
=
\Delta u-f(|z|)+c\,\delta_0
=
\bigl(f(|z|)-c\delta_0\bigr)-f(|z|)+c\delta_0
=0.
\]
Hence \(F\) is a distributional solution of the Cauchy--Riemann equation on \(U\), and by
Weyl's lemma (or, equivalently, hypoellipticity of \(\partial_{\bar z}\)) it is represented by a
genuine holomorphic function on \(U\); the apparent pole of \(\Psi\) and of \(c/(\pi z)\) at the
origin is exactly cancelled by the contribution of the singular term \(-c\delta_0\) in
\eqref{eq:potato-second-proof} arriving through \(4\partial_z\partial_{\bar z}u\). Thus \(F\) is
holomorphic on all of \(U\).

Define
\[
W(z):=\pi z\,F(z).
\]
Then \(W\) is holomorphic on \(U\) and \(W(0)=0\). Since \(\partial_z u\) is continuous up to
\(\partial U\), so is \(W\), and by \eqref{eq:dz-zero-second-proof},
\[
W|_{\partial U}
=
\pi z\left(4\partial_z u-\frac{\Phi(|z|^2)}{z}+\frac{c}{\pi z}\right)\Big|_{\partial U}
=
c-\pi\Phi(|z|^2).
\]
In particular, \(W|_{\partial U}\) is real-valued. The imaginary part of \(W\) is therefore
harmonic in \(U\), continuous up to \(\partial U\), and vanishes on \(\partial U\). By the
maximum principle,
\[
\Im W\equiv 0
\qquad\text{in }U.
\]
Since a holomorphic function with zero imaginary part is constant, there exists \(\kappa\in\R\)
such that
\[
W\equiv \kappa
\qquad\text{in }U.
\]

Restricting to the boundary, and using that \(W\) is continuous up to \(\partial U\) with
boundary value \(c-\pi\Phi(|z|^2)\), we obtain
\[
c-\pi\Phi(|z|^2)=\kappa
\qquad\text{for all }z\in\partial U.
\]
The function \(t\mapsto \Phi(t)=\int_0^t s^k e^{-\pi s}\,ds\) has derivative
\(\Phi'(t)=t^k e^{-\pi t}>0\) for \(t>0\), hence is strictly increasing on \([0,\infty)\). The
identity above therefore forces \(|z|^2\) to take a single value on \(\partial U\), i.e.\ \(|z|\)
is constant on \(\partial U\). Consequently \(\partial U\) is contained in a circle centered at
the origin. Let \(R>0\) be such that
\(\partial U\subset \{|z|=R\}\). Since \(U\) is open and connected, it cannot meet both
\(\{|z|<R\}\) and \(\{|z|>R\}\), because these are disjoint open sets and \(U\cap \{|z|=R\}=\varnothing\).
As \(U\) is bounded, we must therefore have \(U\subset \{|z|<R\}\). If this inclusion were strict,
then \(\{|z|<R\}\setminus U\) would contain a point, and hence, being open, would contribute a
boundary point of \(U\) lying strictly inside the disk \(\{|z|<R\}\), contradicting
\(\partial U\subset \{|z|=R\}\). Therefore
\[
U=\{|z|<R\}
\]
that is, \(U\) is a disk.
\end{proof}

\begin{remark}\label{rem:bookkeeping-no-regularity}
The two proofs use the boundary assumptions in rather different ways, and it is instructive to
keep careful track of the regularity bookkeeping.

In the first proof, once the logarithmic potential \(\Psi_U\) is written explicitly, no elliptic
boundary regularity is needed to conclude that \(\Psi_U\in C^1(\R^2\setminus\{0\})\): this follows
directly from convolution with the locally integrable kernel \(\nabla\log|z|\), exactly as in the
\(L^1_{\mathrm{loc}}\)-translation argument used here for \eqref{eq:explicit-u-second-proof}. Thus
the local radiality argument itself only uses continuity of the gradient up to the boundary. The
genuine extra input appears only in the final geometric step, where one needs that, near a suitable
point \(x_0\in \partial U\), the set \(\partial U\cap B_r(x_0)\) contains a nontrivial boundary
arc, so that the radii attained there form an interval. Accordingly, the first proof does not
really need \(C^2\) regularity; it works under any hypothesis guaranteeing such a local arc
structure, for instance when \(U\) is a Jordan domain. In particular, any Lipschitz, \(C^1\), or
\(C^2\) boundary is more than sufficient.

By contrast, the second proof does not use any boundary regularity at all. The explicit formula
\eqref{eq:explicit-u-second-proof} already yields \(u\in C^1(\R^2\setminus\{0\})\), and since
\(u\equiv 0\) on the exterior open set \(\R^2\setminus \overline U\), continuity implies that both
\(u\) and \(\nabla u\) vanish on \(\partial U\) -- precisely the boundary information packaged in
\eqref{eq:dz-zero-second-proof}. This is exactly the input needed to construct the holomorphic
function \(W=\pi z F\) and conclude that \(|z|\) is constant on \(\partial U\). Hence the second
proof recovers the theorem under the original assumptions of Abreu and D\"orfler
\cite{Abreu-Doerfler}, namely for arbitrary bounded simply connected domains, with no boundary
smoothness hypothesis -- this is the content of
Theorem~\ref{thm:eig-loc-U-strong}. The mechanism is also conceptually closer to the classical
quadrature-domain rigidity results of \cite{Aharonov-Schiffer-Zalcman,Gustafsson,Karp}, where the
existence of a holomorphic function on \(U\) with real boundary values forces the boundary to lie
on a level set of a real-analytic function, here the elementary radial function
\(z\mapsto |z|^2\).

To make the comparison explicit: the only step at which the first proof
genuinely uses \(\partial U\in C^2\) is in the open-mapping argument that
extracts a nontrivial open interval of radii from a boundary parametrization
\(\gamma:(-\varepsilon,\varepsilon)\to\partial U\). In the second proof, this
geometric step is replaced by the harmonic-function-with-real-boundary-values
mechanism, which is purely topological (continuity of \(W=\pi zF\) up to
\(\partial U\) plus the maximum principle on the open set \(U\)) and is
therefore insensitive to the regularity of \(\partial U\). It is in this
precise sense that Theorem~\ref{thm:eig-loc-U-strong} drops the \(C^2\)
hypothesis present in Theorem~\ref{thm:eig-loc-U} while delivering the same
conclusion.
\end{remark}

% ==============================================================
% Section 5
% ==============================================================
\section{Proof of Theorem \ref{thm:GGRT-sharp}}

We prove that the quantitative stability estimate in Corollary~1.4 of
\cite{Gomez-Guerra-Ramos-Tilli}, which builds on the rigidity result of
\cite{Nicola-Tilli} characterizing the disk as the unique optimizer of the
first localization eigenvalue, is sharp. The proof proceeds by constructing a
family of domains which are perturbations of the disk, whose first localization
eigenvalue has the same first-order behaviour as the disk, but whose distance
to the class of disks is of order $\varepsilon$.

\begin{proof}[Proof of Theorem~\ref{thm:GGRT-sharp}]
Fix
\[
\lambda_0:=1-e^{-\pi},
\]
which is the eigenvalue of the localization operator of the unit disk $D_1$
associated with the ground state $h_0$. For $\varepsilon$ small, define
\[
P_\varepsilon(w):=1+\varepsilon w^2.
\]
By Theorem~\ref{thm:perturb}, applied with \(\lambda=\lambda_0\) and
\(P_\varepsilon(w)=1+\varepsilon w^2\), for all sufficiently small
\(|\varepsilon|\) there exists a domain \(U_\varepsilon\) such that
\begin{equation}\label{eq:functiona-P-eps-new}
\int_{U_\varepsilon} P_\varepsilon(z)e^{\pi \overline zw}e^{-\pi|z|^2}\,dA(z)
=
\lambda_0 P_\varepsilon(w),
\qquad w\in\C.
\end{equation}
Moreover, by construction, $U_\varepsilon$ is a $C^2$ perturbation of the unit
disk. Hence, after possibly shrinking the range of admissible $\varepsilon$,
its boundary can be written as
\begin{equation}\label{eq:def-u-eps-new}
\partial U_\varepsilon
=
\{(1+\varphi_\varepsilon(y))y:\ y\in \mathbb S^1\},
\end{equation}
where
\[
\varphi_\varepsilon\in C^2(\mathbb S^1),
\qquad
\|\varphi_\varepsilon\|_{C^2(\mathbb S^1)}\lesssim |\varepsilon|.
\]
The map
\[
\varepsilon\mapsto \varphi_\varepsilon
\]
may in fact be chosen \(C^1\) as a map from a neighborhood of \(0\) into
\(C^1(\mathbb S^1)\). To see this, we briefly retrace the proof of
Theorem~\ref{thm:perturb}. The construction is the disk moment-map
construction of Proposition~\ref{prop:disk-inverse-tailored}. Thus
\(\varphi_\varepsilon\) is obtained from an equation of the form
\(F(\varepsilon,\varphi)=0\), where
\[
F:(-\varepsilon_0,\varepsilon_0)\times \mathcal U\to \mathcal V
\]
is the \(C^1\) moment map between the analytic boundary spaces used there,
with \(F(0,0)=0\) and \(\partial_\varphi F(0,0)\) invertible by the explicit
Fourier diagonalization at the unit disk. Since the polynomial
\(P_\varepsilon(z)=1+\varepsilon z^2\)
depends polynomially on \(\varepsilon\), the dependence of \(F\) on
\(\varepsilon\) is \(C^\infty\). The \(C^1\) version of the implicit function
theorem in Banach spaces then yields a \(C^1\) branch
\(\varepsilon\mapsto \varphi_\varepsilon\) into \(C^1(\mathbb S^1)\). We
therefore write
\[
\dot\varphi_0(y):=
\left.\partial_\varepsilon\right|_{\varepsilon=0}\varphi_\varepsilon(y).
\]

\medskip
\noindent
{\bf Step 1: first variation of the eigenvalue identity.} We differentiate \eqref{eq:functiona-P-eps-new} at $\varepsilon=0$.
Since \(P_\varepsilon\to 1\) in \(C^\infty\) and
\(\varepsilon\mapsto \varphi_\varepsilon\) is \(C^1\) into \(C^1(\mathbb S^1)\),
all terms in \eqref{eq:functiona-P-eps-new} are differentiable at
\(\varepsilon=0\). Since $U_\varepsilon$ is given by a normal graph over
$\mathbb S^1$, we apply Reynolds' transport formula. Concretely, in polar
coordinates \(z=re^{it}\), for any smooth test function \(H\) we have
\[
\int_{U_\varepsilon} H(z)\,dA(z)
=
\int_0^{2\pi}\int_0^{1+\varphi_\varepsilon(e^{it})}
H(re^{it})\,r\,dr\,dt.
\]
Differentiating in \(\varepsilon\) at \(\varepsilon=0\) and using the
fundamental theorem of calculus,
\[
\left.\frac{d}{d\varepsilon}\right|_{\varepsilon=0}
\int_{U_\varepsilon} H(z)\,dA(z)
=
\int_0^{2\pi}\dot\varphi_0(e^{it})H(e^{it})\,dt
=
\int_{\mathbb S^1}\dot\varphi_0(z)H(z)\,d\mathcal H^1(z),
\]
where the factor \(r=1\) on the unit circle has been absorbed into the
arc-length measure. This is precisely Reynolds' transport identity:
\[
\left.\frac{d}{d\varepsilon}\right|_{\varepsilon=0}
\int_{U_\varepsilon} H(z)\,dA(z)
=
\int_{\mathbb S^1} \dot\varphi_0(z)H(z)\,d\mathcal H^1(z).
\]
Applying this to
\[
H(z)=e^{\pi \overline zw}e^{-\pi|z|^2},
\]
and differentiating the factor $P_\varepsilon(z)=1+\varepsilon z^2$, we obtain
\begin{equation}\label{eq:first-variation-identity-new}
\int_{\mathbb S^1}
\dot\varphi_0(z)e^{\pi \overline zw}e^{-\pi|z|^2}\,d\mathcal H^1(z)
+
\int_{D_1} z^2 e^{\pi \overline zw}e^{-\pi|z|^2}\,dA(z)
=
\lambda_0 w^2.
\end{equation}
Since $D_1$ is radial, the monomial $z^2$ is an eigenfunction of the
localization operator of $D_1$; therefore there exists $\lambda_2>0$ such that
\[
\int_{D_1} z^2 e^{\pi \overline zw}e^{-\pi|z|^2}\,dA(z)=\lambda_2 w^2.
\]
Hence \eqref{eq:first-variation-identity-new} becomes
\begin{equation}\label{eq:boundary-identity-new}
\int_{\mathbb S^1}
\dot\varphi_0(z)e^{\pi \overline zw}e^{-\pi|z|^2}\,d\mathcal H^1(z)
=
(\lambda_0-\lambda_2)w^2.
\end{equation}
In fact, \(\lambda_2\) can be computed explicitly. Expanding
\(e^{\pi \overline zw}=\sum_{n\geq 0}\frac{\pi^n}{n!}\overline z^{\,n}w^n\)
and using orthogonality of the monomials \(\overline z^{\,n}\) over the disk,
the only term contributing to \(w^2\) is the one with \(n=2\):
\[
\int_{D_1} z^2 e^{\pi \overline zw}e^{-\pi|z|^2}\,dA(z)
=
\frac{\pi^2}{2}w^2\int_{D_1}z^2\overline z^{\,2}e^{-\pi|z|^2}\,dA(z)
=
\frac{\pi^2}{2}w^2\int_{D_1}|z|^4 e^{-\pi|z|^2}\,dA(z).
\]
Passing to polar coordinates and integrating by parts twice,
\[
\int_{D_1}|z|^4 e^{-\pi|z|^2}\,dA(z)
=
2\pi\int_0^1 r^4 e^{-\pi r^2}\,r\,dr
=
\pi\int_0^1 s^2 e^{-\pi s}\,ds
\]
where we substituted \(s=r^2\). A direct computation, integrating by parts
twice, yields
\[
\pi\int_0^1 s^2 e^{-\pi s}\,ds
=
\frac{2}{\pi^2}\Bigl(1-\Bigl(1+\pi+\frac{\pi^2}{2}\Bigr)e^{-\pi}\Bigr).
\]
Multiplying by \(\frac{\pi^2}{2}\), we obtain
\[
\lambda_2
=
\frac{\pi^2}{2}\int_{D_1}|z|^4e^{-\pi|z|^2}\,dA(z)
=
1-\Bigl(1+\pi+\frac{\pi^2}{2}\Bigr)e^{-\pi}.
\]
Therefore
\[
\lambda_0-\lambda_2
=
(1-e^{-\pi})-\Bigl(1-\Bigl(1+\pi+\frac{\pi^2}{2}\Bigr)e^{-\pi}\Bigr)
=
\pi\Bigl(1+\frac{\pi}{2}\Bigr)e^{-\pi}>0,
\]
so the coefficient of \(w^2\) on the right-hand side is indeed nonzero.
In particular, since $|z|=1$ on $\mathbb S^1$, we may rewrite this as
\begin{equation}\label{eq:boundary-identity-circle-new}
e^{-\pi}
\int_{\mathbb S^1}
\dot\varphi_0(z)e^{\pi \overline zw}\,d\mathcal H^1(z)
=
(\lambda_0-\lambda_2)w^2.
\end{equation}

Setting $w=0$ gives
\begin{equation}\label{eq:cancellation-f-0-new}
\int_{\mathbb S^1}\dot\varphi_0\,d\mathcal H^1=0.
\end{equation}

\medskip
\noindent
{\bf Step 2: identification of the leading mode.} Expanding
\[
e^{\pi \overline zw}
=
\sum_{n\geq 0}\frac{\pi^n}{n!}\overline z^{\,n}w^n,
\]
and comparing coefficients in \eqref{eq:boundary-identity-circle-new}, we find
\[
\int_{\mathbb S^1}\dot\varphi_0(z)\overline z^{\,n}\,d\mathcal H^1(z)=0
\qquad \text{for all } n\neq 2,
\]
whereas the coefficient corresponding to $n=2$ is nonzero. Since
$\dot\varphi_0$ is real-valued, it follows that
\begin{equation}\label{eq:mode-decomposition-new}
\dot\varphi_0(e^{it})
=
a\cos(2t)+b\sin(2t)
\end{equation}
for some pair $(a,b)\neq (0,0)$. In particular,
\begin{equation}\label{eq:l1-nontrivial-new}
\int_{\mathbb S^1} |\dot\varphi_0|\,d\mathcal H^1>0.
\end{equation}

\medskip
\noindent
{\bf Step 3: normalization of the area.} The area of $U_\varepsilon$ is given by
\begin{equation}\label{eq:area-u-eps-new}
|U_\varepsilon|
=
\frac12\int_{\mathbb S^1} (1+\varphi_\varepsilon(y))^2\,d\mathcal H^1(y).
\end{equation}
Since
\[
\varphi_\varepsilon
=
\varepsilon \dot\varphi_0 + O(\varepsilon^2)
\quad \text{in } C^1(\mathbb S^1),
\]
and \eqref{eq:cancellation-f-0-new} gives the vanishing of the first-order
term, we obtain from \eqref{eq:area-u-eps-new}
\begin{equation}\label{eq:measure-close-u-eps-new}
|U_\varepsilon|=\pi+O(\varepsilon^2).
\end{equation}

Let
\[
r_\varepsilon:=\sqrt{\frac{\pi}{|U_\varepsilon|}},
\qquad
\widetilde U_\varepsilon:=r_\varepsilon U_\varepsilon.
\]
Then $|\widetilde U_\varepsilon|=\pi$, and by
\eqref{eq:measure-close-u-eps-new},
\[
r_\varepsilon=1+O(\varepsilon^2).
\]
Hence $\widetilde U_\varepsilon$ still admits a normal graph representation
\[
\partial \widetilde U_\varepsilon
=
\{(1+\widetilde\varphi_\varepsilon(y))y:\ y\in \mathbb S^1\},
\]
with
\[
\widetilde\varphi_\varepsilon
=
\varepsilon \dot\varphi_0 + O(\varepsilon^2)
\quad \text{in } C^1(\mathbb S^1).
\]
Thus the area normalization does not alter the first-order deformation.

\medskip
\noindent
{\bf Step 4: distance from the centered disk.} Since both $\widetilde U_\varepsilon$ and $D_1$ are star-shaped with respect to
the origin, the symmetric difference may be computed in polar coordinates:
\[
|\widetilde U_\varepsilon \triangle D_1|
=
\frac12\int_{\mathbb S^1}
\left|(1+\widetilde\varphi_\varepsilon(y))^2-1\right|
\,d\mathcal H^1(y).
\]
Since \(\widetilde\varphi_\varepsilon=\varepsilon\dot\varphi_0+O(\varepsilon^2)\)
in \(C^1(\mathbb S^1)\), we have in particular
\(\|\widetilde\varphi_\varepsilon\|_{L^\infty(\mathbb S^1)}\lesssim |\varepsilon|\),
whence \(\widetilde\varphi_\varepsilon^2=O(\varepsilon^2)\) in
\(L^\infty(\mathbb S^1)\). Hence
\[
(1+\widetilde\varphi_\varepsilon)^2-1
=
2\widetilde\varphi_\varepsilon+\widetilde\varphi_\varepsilon^2
=
2\varepsilon \dot\varphi_0 + O(\varepsilon^2)
\quad \text{in } L^\infty(\mathbb S^1),
\]
where the implied constant depends only on the \(C^1\)-norm bound on
\(\widetilde\varphi_\varepsilon\). Therefore
\[
|\widetilde U_\varepsilon \triangle D_1|
=
\frac12\int_{\mathbb S^1}
\left|2\varepsilon \dot\varphi_0 + O(\varepsilon^2)\right|
\,d\mathcal H^1.
\]
Set \(\alpha:=\int_{\mathbb S^1}|\dot\varphi_0|\,d\mathcal H^1\). By
\eqref{eq:l1-nontrivial-new}, \(\alpha>0\). The reverse triangle inequality in
\(L^1\) and the \(O(\varepsilon^2)\) remainder give
\[
\int_{\mathbb S^1}\bigl|2\varepsilon\dot\varphi_0+O(\varepsilon^2)\bigr|\,d\mathcal H^1
\geq
2|\varepsilon|\alpha-O(\varepsilon^2)
\]
and similarly the upper bound
\[
\int_{\mathbb S^1}\bigl|2\varepsilon\dot\varphi_0+O(\varepsilon^2)\bigr|\,d\mathcal H^1
\leq
2|\varepsilon|\alpha+O(\varepsilon^2).
\]
Consequently there exist constants $c,C>0$ such that
\begin{equation}\label{eq:distance-centered-disk-new}
c|\varepsilon|
\leq
|\widetilde U_\varepsilon \triangle D_1|
\leq
C|\varepsilon|
\end{equation}
for all sufficiently small $\varepsilon$.

\medskip
\noindent
{\bf Step 5: comparison with arbitrary disks.} Let $D$ be any disk of area $\pi$. Then $D=D_1+a$ for some $a\in \C$. We first consider disks whose centers are not too close to the origin.
Since
\[
|D_1\triangle (D_1+a)|
\gtrsim \min\{|a|,1\},
\]
the triangle inequality and \eqref{eq:distance-centered-disk-new} imply that
there exists $K\gg 1$ such that, if $|a|\geq K|\varepsilon|$, then
\begin{equation}\label{eq:far-disks-new}
|\widetilde U_\varepsilon \triangle D|
\geq
|D_1\triangle D|-|\widetilde U_\varepsilon\triangle D_1|
\gtrsim
|a|-C|\varepsilon|
\gtrsim
|\varepsilon|.
\end{equation}

It remains to consider disks whose centers satisfy $|a|\leq K|\varepsilon|$.
Write
\[
a=\eta \varepsilon e^{i\theta},
\qquad
|\eta|\leq K,
\quad
\theta\in [0,2\pi).
\]
We claim that, for $\varepsilon$ sufficiently small, the translated domain
\[
\widetilde U_\varepsilon-a
\]
is still a normal graph over $\mathbb S^1$:
\[
\partial(\widetilde U_\varepsilon-a)
=
\{(1+\psi_\varepsilon^{\theta,\eta}(y))y:\ y\in \mathbb S^1\},
\]
with
\[
\psi_\varepsilon^{\theta,\eta}\in C^1(\mathbb S^1),
\qquad
\|\psi_\varepsilon^{\theta,\eta}\|_{C^1}\lesssim |\varepsilon|,
\]
uniformly in $(\theta,\eta)$ with $|\eta|\leq K$. This follows because
$\widetilde U_\varepsilon$ is already an $O(\varepsilon)$-$C^1$ perturbation of
$D_1$, and the translation vector $a$ is also of order $\varepsilon$. We now compute the first-order term of $\psi_\varepsilon^{\theta,\eta}$.
A translation by $a$ contributes, to first order in $\varepsilon$, the normal
displacement
\[
-\eta \Re(e^{-i\theta}y).
\]
Indeed, for \(y\in \mathbb S^1\),
\[
|y-a|
=
\bigl(1-2\Re(\overline y a)+|a|^2\bigr)^{1/2}
=
1-\Re(\overline y a)+O(|a|^2),
\]
uniformly in \(y\), and here \(a=\eta\varepsilon e^{i\theta}\).
Therefore
\begin{equation}\label{eq:psi-expansion-new}
\psi_\varepsilon^{\theta,\eta}(y)
=
\varepsilon\Big(\dot\varphi_0(y)-\eta \Re(e^{-i\theta}y)\Big)
+
O(\varepsilon^2)
\quad \text{in } C^1(\mathbb S^1),
\end{equation}
uniformly for $|\eta|\leq K$.

Define
\[
g^{\theta,\eta}(y)
:=
\dot\varphi_0(y)-\eta \Re(e^{-i\theta}y).
\]
Using \eqref{eq:mode-decomposition-new}, we see that $g^{\theta,\eta}$ belongs
to the finite-dimensional space spanned by the first and second harmonics.
More precisely,
\[
g^{\theta,\eta}(e^{it})
=
a\cos(2t)+b\sin(2t)-\eta \cos(t-\theta).
\]
The second-harmonic part is independent of $(\theta,\eta)$ and nonzero, because
$(a,b)\neq (0,0)$. To extract a uniform \(L^1\) lower bound, note that
\(g^{\theta,\eta}\) lies in the \(4\)-dimensional space
\[
V:=\mathrm{span}\bigl\{\cos t,\sin t,\cos 2t,\sin 2t\bigr\}\subset L^1(\mathbb S^1).
\]
On this finite-dimensional space, all norms are equivalent; in particular, the
\(L^1(\mathbb S^1)\) and \(L^\infty(\mathbb S^1)\) norms are equivalent, and
convergence in coefficients is equivalent to convergence in any of these norms. The map
\[
(\theta,\eta)\in [0,2\pi)\times [-K,K]\longmapsto g^{\theta,\eta}\in V
\]
is continuous, and its image
\[
\mathcal G_K:=\{g^{\theta,\eta}:\ \theta\in[0,2\pi),\ |\eta|\leq K\}
\]
is the continuous image of a compact set, hence compact in \(V\). Moreover, the
\((\cos 2t,\sin 2t)\)-coefficients of every \(g^{\theta,\eta}\) are exactly
\((a,b)\neq (0,0)\), independent of \((\theta,\eta)\); thus \(0\notin
\mathcal G_K\). The function \(g\mapsto\|g\|_{L^1(\mathbb S^1)}\) is continuous
and strictly positive on \(V\setminus\{0\}\), so it attains a positive minimum
on the compact set \(\mathcal G_K\). It follows that there exists
$\delta_K>0$ such that
\begin{equation}\label{eq:uniform-l1-lower-new}
\int_{\mathbb S^1}|g^{\theta,\eta}(y)|\,d\mathcal H^1(y)\geq \delta_K
\qquad
\text{for all }\theta\in[0,2\pi),\ |\eta|\leq K.
\end{equation}

Now,
\[
|\widetilde U_\varepsilon \triangle D|
=
|(\widetilde U_\varepsilon-a)\triangle D_1|.
\]
Since $\widetilde U_\varepsilon-a$ is star-shaped with radial graph
$1+\psi_\varepsilon^{\theta,\eta}$, the same polar-coordinate computation as in
Step~4 gives
\[
|(\widetilde U_\varepsilon-a)\triangle D_1|
=
\frac12\int_{\mathbb S^1}
\left|(1+\psi_\varepsilon^{\theta,\eta}(y))^2-1\right|
\,d\mathcal H^1(y).
\]
Using \eqref{eq:psi-expansion-new},
\[
(1+\psi_\varepsilon^{\theta,\eta})^2-1
=
2\varepsilon g^{\theta,\eta}+O(\varepsilon^2)
\quad \text{in } L^\infty(\mathbb S^1),
\]
uniformly for $|\eta|\leq K$. Thus, by \eqref{eq:uniform-l1-lower-new},
there exists $c_K>0$ such that
\begin{equation}\label{eq:near-disks-new}
|\widetilde U_\varepsilon \triangle D|
=
|(\widetilde U_\varepsilon-a)\triangle D_1|
\geq
c_K |\varepsilon|
\end{equation}
whenever $|a|\leq K|\varepsilon|$. Combining \eqref{eq:far-disks-new} and \eqref{eq:near-disks-new}, we conclude
that
\[
\inf_{\substack{D \text{ disk}\\ |D|=\pi}}
|\widetilde U_\varepsilon \triangle D|
\gtrsim |\varepsilon|.
\]
Since \(r_\varepsilon=1+O(\varepsilon^2)\), the dilation from \(U_\varepsilon\)
to \(\widetilde U_\varepsilon\) changes the set only by
\[
|\widetilde U_\varepsilon\triangle U_\varepsilon|=O(\varepsilon^2),
\]
and therefore does not affect the first-order distance estimate. Thus the
area-normalized family \(\widetilde U_\varepsilon\) remains at
symmetric-difference distance \(\gtrsim |\varepsilon|\) from the class of
disks.

\medskip
\noindent
{\bf Step 6: square-root sharpness of the deficit estimate.} It remains to translate the distance estimate into the corresponding deficit
estimate. Recall that, by construction via Theorem~\ref{thm:perturb}, the
first localization eigenvalue of \(U_\varepsilon\) equals \(\lambda_0\)
\emph{exactly}, since \eqref{eq:functiona-P-eps-new} states that
\(P_\varepsilon\) is an eigenfunction with eigenvalue \(\lambda_0\) for the
localization operator on \(U_\varepsilon\), and a standard variational
argument shows that this is the principal eigenvalue. By
\eqref{eq:measure-close-u-eps-new},
\(|U_\varepsilon|=\pi+O(\varepsilon^2)\), and rescaling by
\(r_\varepsilon=1+O(\varepsilon^2)\) to area \(\pi\) changes the eigenvalue
only at order \(O(\varepsilon^2)\) (since the localization operator scales
analytically in the dilation parameter, and the disk is a critical point for
this scaling). It follows that the eigenvalue deficit of
\(\widetilde U_\varepsilon\) with respect to the disk satisfies
\[
\lambda_0-\lambda(\widetilde U_\varepsilon)=O(\varepsilon^2),
\]
where \(\lambda(\widetilde U_\varepsilon)\) denotes the first localization
eigenvalue of \(\widetilde U_\varepsilon\). On the other hand, by Step~5,
the symmetric-difference asymmetry satisfies
\[
\inf_{\substack{D\text{ disk}\\ |D|=\pi}}|\widetilde U_\varepsilon\triangle D|
\gtrsim |\varepsilon|.
\]
Comparing the two,
\[
\inf_{\substack{D\text{ disk}\\ |D|=\pi}}|\widetilde U_\varepsilon\triangle D|
\gtrsim |\varepsilon|
\simeq
\sqrt{\lambda_0-\lambda(\widetilde U_\varepsilon)}.
\]
This shows that the inequality
\[
\inf_{\substack{D\text{ disk}\\ |D|=\pi}}
|\Omega\triangle D|
\lesssim
\sqrt{\lambda_0-\lambda(\Omega)}
\]
of Corollary~1.4 in \cite{Gomez-Guerra-Ramos-Tilli} is square-root-tight along
the family \(\widetilde U_\varepsilon\): the deficit is of order
\(\varepsilon^2\), the asymmetry is of order \(\varepsilon\), and so the
square root of the deficit and the asymmetry are of comparable size. This
exhibits the optimality of the exponent \(1/2\) and concludes the proof of the
sharpness of the estimate in Corollary~1.4 of
\cite{Gomez-Guerra-Ramos-Tilli}, complementing the rigidity result of
\cite{Nicola-Tilli}.

\end{proof}

% ==============================================================
% Section 6
% ==============================================================
\section{Local maximizers for the Faber--Krahn problem in the Gaussian STFT setting}

The purpose of this section is to establish a local rigidity statement for the
Faber--Krahn problem associated with the Gaussian STFT. The global problem,
solved in \cite{Nicola-Tilli}, asserts that, among measurable subsets of \(\C\)
of fixed measure, the centered disk uniquely maximizes the first eigenvalue of
the Gaussian STFT localization operator. Here we show that this rigidity is
already visible at the \emph{local} level: any set whose first eigenvalue cannot
be increased by small symmetric-difference perturbations must, up to a
translation, be a disk. This complements the quantitative stability result of
\cite{Gomez-Guerra-Ramos-Tilli} and provides the variational counterpart of the
inverse-spectral analysis carried out in the previous sections.

We work throughout in the Bargmann--Fock representation; we refer the reader to
\cite{Bargmann,Folland,Grochenig} for the relevant background. Thus, if
\(F\in \mathcal F^2(\C)\) is the Bargmann transform of \(f\in L^2(\R)\), then
\[
|\mathcal V_\varphi f(z)|^2 = |F(z)|^2 e^{-\pi |z|^2},
\]
up to the normalization already fixed in the previous sections. Accordingly, for
\(F\in \mathcal F^2(\C)\) we set
\[
u_F(z):=|F(z)|^2 e^{-\pi |z|^2}.
\]

For a measurable set \(\Omega\subset \C\), let
\begin{equation}\label{eq:lambda1-def-sec6}
\lambda_{1,\varphi}(\Omega)
:=
\sup_{\|f\|_2=1}\int_\Omega |\mathcal V_\varphi f(z)|^2\,dz.
\end{equation}

\subsection{The weighted symmetric-difference topology}

The locality of a maximizer must be measured in a topology compatible with the
Lebesgue-measure constraint and with the natural symmetries of the localization
operator. The right choice in our setting is the weighted symmetric-difference
metric obtained by integrating against the Gaussian density
\(e^{-\pi|z|^2}\,dA(z)\), which is the same density that appears both in the
Fock-space norm and in the eigenfunction \(u_F\).

\begin{definition}[Weighted symmetric-difference metric]\label{def:weighted-sym-diff}
For measurable sets \(\Omega_1,\Omega_2\subset \C\) of finite Lebesgue measure,
set
\begin{equation}\label{eq:weighted-sym-diff}
d_w(\Omega_1,\Omega_2)
:=
\int_{\Omega_1\triangle\Omega_2} w(z)\,dA(z),
\qquad
w(z):=e^{-\pi |z|^2}.
\end{equation}
We call \(d_w\) the \emph{weighted symmetric-difference distance}, and we denote
by \(\mathcal M_a\) the Borel-measurable subsets of \(\C\) of Lebesgue measure
\(a>0\), modulo Lebesgue-null modifications. The map
\(d_w\colon \mathcal M_a\times \mathcal M_a\to[0,\infty)\) is a genuine metric
on \(\mathcal M_a\): symmetry, the triangle inequality, and vanishing only on
null modifications all follow from \(w(z)=e^{-\pi|z|^2}>0\) and standard
properties of the symmetric difference.
\end{definition}

\begin{remark}[Naturality of \(d_w\) for the Gaussian Faber--Krahn problem]\label{rem:topology-natural}
Three observations clarify why \(d_w\) is the natural topology for our problem.

\begin{enumerate}
\item[\textnormal{(i)}] (\emph{Compatibility with the measure constraint.})
Since \(w\) is bounded above by \(1\) and bounded below by \(e^{-\pi R^2}>0\)
on every disk \(\{|z|\le R\}\), the metric \(d_w\) is, on every fixed bounded
region of \(\C\), equivalent to the unweighted symmetric-difference distance
\(d_0(\Omega_1,\Omega_2):=|\Omega_1\triangle\Omega_2|\). Moreover, every
\(d_w\)-perturbation \(E\) of \(\Omega\) used in this section is constructed
to satisfy the equimeasurability constraint \(|E|=|\Omega|\); hence the
\(d_w\)-balls intersected with \(\mathcal M_a\) form a basis of neighborhoods
on the measure-constraint manifold.

\item[\textnormal{(ii)}] (\emph{Translation- and rotation-invariance via Weyl.})
The Gaussian density \(w(z)=e^{-\pi|z|^2}\) is the symbol density associated
with the metaplectic representation: under the Weyl translations
\(T_a\colon F\mapsto e^{\pi az-\frac{\pi}{2}|a|^2}F(z-a)\) (cf.\
\eqref{eq:Ta}) and rotations \(F(z)\mapsto F(e^{-i\theta}z)\), one has the
covariance
\[
u_{T_aF}(z)=u_F(z-a),\qquad u_{F(\,e^{-i\theta}\cdot\,)}(z)=u_F(e^{-i\theta}z),
\]
so that the spectral data \((\lambda_{1,\varphi},u_F)\) of the localization
operator transform unitarily under these symmetries. Any choice of weight
other than a Gaussian would break this covariance with the operator algebra of
the Fock space; see, e.g., \cite{Folland,Grochenig}.

\item[\textnormal{(iii)}] (\emph{Control of Fock-space tails.}) For any
\(F\in \mathcal F^2(\C)\) one has \(u_F(z)\le \|F\|_{\mathcal F^2}^2\) for
every \(z\in\C\), as a consequence of the reproducing property
\eqref{eq:reproducing-kernel} below and Cauchy--Schwarz. Hence, for first
eigenfunctions \(F\) (which are bounded; see Lemma~\ref{lem:simplicity-local-max}),
\[
\Bigl|\int_{\Omega_1}u_F\,dA-\int_{\Omega_2}u_F\,dA\Bigr|
\le
\int_{\Omega_1\triangle\Omega_2}u_F\,dA
\le
\|F\|_{\mathcal F^2}^2\,e^{\pi R^2}\, d_w(\Omega_1,\Omega_2)
\]
on any bounded region \(\{|z|\le R\}\), so that the localization functional
\(\Omega\mapsto \int_\Omega u_F\,dA\) is Lipschitz in \(d_w\) on bounded
sets. The Gaussian weight thus tames the lack of compactness at infinity by
ensuring that \(d_w\)-small perturbations cannot transport mass to large
\(|z|\) without paying a proportional cost.
\end{enumerate}

In summary, \(d_w\) is essentially the unique Gaussian-weighted \(L^1\)-metric
on indicator functions that is compatible with the equimeasurable constraint
\(|\Omega|=a\), respects the Weyl-translation and rotation symmetries of the
Fock space, and renders \(\lambda_{1,\varphi}\) and the localization energy
continuous on bounded sets.
\end{remark}

\begin{definition}[Local maximizer]\label{def:local-maximizer}
A set \(\Omega\in\mathcal M_a\) is a \emph{local maximizer} for
\(\lambda_{1,\varphi}\) under the Lebesgue-measure constraint
\(|\,\cdot\,|=a\) if there exists \(\delta_0>0\) such that, whenever
\(E\in \mathcal M_a\) satisfies \(d_w(\Omega,E)<\delta_0\), one has
\begin{equation}\label{eq:local-max-ineq}
\lambda_{1,\varphi}(\Omega)\ge \lambda_{1,\varphi}(E).
\end{equation}
\end{definition}

We can now state the main theorem of this section.

\begin{theorem}[Local maximizers are disks]\label{thm:local-max}
Let \(\Sigma\in\mathcal M_a\) be a local maximizer for \(\lambda_{1,\varphi}\)
under the Lebesgue-measure constraint, in the sense of
Definition~\ref{def:local-maximizer}. Then there exists \(a_0\in\C\) such that
\(\Sigma\) coincides, up to a Lebesgue-null modification, with the centered
disk
\[
\Sigma=D_R+a_0,\qquad R=\sqrt{a/\pi}.
\]
\end{theorem}

The remainder of this section is devoted to the proof of
Theorem~\ref{thm:local-max}. The
proof proceeds in six steps. Steps 1--3 are essentially the analogues, in our
setting, of standard arguments in the Faber--Krahn theory: \(\Sigma\) is forced
to be a level set of \(u_F\), the corresponding eigenvalue is simple, and a
first-variation identity ties \(F\) to a divergence-free vector field built from
holomorphic data. The crucial new ingredient lies in Step 4, where we show that
the \emph{derivative} \(F'\) is again a first eigenfunction. Combined with the
simplicity from Step 2, this forces \(F\) to satisfy a first-order linear ODE,
and hence to be a Gaussian. Once this is in place, Steps 5--6 use the Weyl
invariance of the Bargmann--Fock setting and a one-dimensional rearrangement
argument to identify \(\Sigma\) with a centered disk.

\begin{proof}[Proof of the local-maximizer theorem]
Let \(\Sigma\subset \C\) be a local maximizer for \(\lambda_{1,\varphi}\) under
the measure constraint, and let \(F\in \mathcal F^2(\C)\) be a normalized first
eigenfunction for the localization operator associated with \(\Sigma\).
Equivalently,
\begin{equation}\label{eq:eigenvalue-equation}
\int_\Sigma H(z)\overline{F(z)}\,d\mu(z)
=
\lambda_{1,\varphi}(\Sigma)\int_\C H(z)\overline{F(z)}\,d\mu(z)
\end{equation}
for every \(H\in \mathcal F^2(\C)\), where
\[
d\mu(z):=e^{-\pi |z|^2}\,dA(z).
\]

The proof of the superlevel-set characterization in Step~1 below uses the
following classical rearrangement inequality, in the formulation of
Lieb--Loss \cite[Theorem~1.14]{Lieb-Loss} (the so-called
\emph{bathtub principle}). We record it here for ease of reference and to fix
notation.

\begin{lemma}[Bathtub principle]\label{lem:bathtub}
Let \((X,\Sigma_X,\mu_X)\) be a measure space and let
\(g\colon X\to[0,\infty)\) be a measurable function such that
\(\mu_X(\{g>t\})<\infty\) for every \(t>0\). Fix \(0<a<\mu_X(X)\). Then
\begin{equation}\label{eq:bathtub-equation}
\sup\Bigl\{\int_X g\,\mathbf 1_E\,d\mu_X\ \Big|\ E\in\Sigma_X,\ \mu_X(E)=a\Bigr\}
\end{equation}
is attained, and any maximizer \(E^\ast\) coincides, up to a \(\mu_X\)-null
set, with a superlevel set \(\{g>t_\ast\}\) of \(g\) (possibly with a
measure-zero correction at the level set \(\{g=t_\ast\}\) needed to fit the
constraint \(\mu_X(E^\ast)=a\)). The threshold \(t_\ast\ge 0\) is uniquely
determined by the condition
\[
\mu_X(\{g>t_\ast\})\le a\le \mu_X(\{g\ge t_\ast\}),
\]
and the maximum is strict in the sense that, if \(g\) has no plateaus of
positive measure on the level \(\{g=t_\ast\}\), then \(E^\ast=\{g>t_\ast\}\)
\(\mu_X\)-a.e.
\end{lemma}

In our application, \((X,\mu_X)=(\C,\mathrm{Leb}_\C)\) and \(g=u_F\). Since
\(u_F\) is real-analytic on the (open) complement of the zero set of \(F\),
its level sets \(\{u_F=t\}\) are real-analytic varieties of strictly lower
dimension and hence Lebesgue-null for every \(t>0\); the no-plateau
hypothesis of Lemma~\ref{lem:bathtub} is therefore automatic, and the
maximizer is unique up to null sets.

\medskip
\noindent
{\bf Step 1: \(\Sigma\) is a superlevel set of \(u_F\).}\label{step:superlevel}

We claim that there exists \(t_0\ge 0\) such that, up to sets of measure zero,
\begin{equation}\label{eq:superlevel-claim}
\Sigma=\{u_F>t_0\}.
\end{equation}

The strategy is a bathtub-style perturbation argument: if \(\Sigma\) were not a
superlevel set of \(u_F\), one could move a small piece of \(\Sigma\) where
\(u_F\) is small to a region where \(u_F\) is larger, increasing the
localization integral while preserving Lebesgue measure and remaining within
the local-maximizer neighborhood of \(\Sigma\) (cf.\ Lemma~\ref{lem:bathtub}).
The contradiction we extract is local, but as we shall see in
Lemma~\ref{lem:bathtub-superlevel-from-localmax}, the same argument actually
proves the slightly stronger statement that \(\Sigma\) is the unique
\(\mu\)-essential maximizer of the functional \(E\mapsto \int_E u_F\,dA\)
under the equimeasurable constraint, in agreement with the bathtub principle.

Recall that \(F\) is admissible in the variational problem defining
\(\lambda_{1,\varphi}(\Sigma')\) for any measurable set \(\Sigma'\) of finite
measure. In particular,
\[
\lambda_{1,\varphi}(\Sigma')\ge \int_{\Sigma'} u_F\,dA,
\qquad
\lambda_{1,\varphi}(\Sigma)= \int_\Sigma u_F\,dA,
\]
the latter equality holding because \(F\) is the actual eigenfunction for
\(\Sigma\). It therefore suffices to find a competitor \(\Sigma'\) with
\(|\Sigma'|=|\Sigma|\), close to \(\Sigma\) in the sense of weighted symmetric
difference, and with \(\int_{\Sigma'} u_F\,dA>\int_\Sigma u_F\,dA\); that would
violate local maximality.

Suppose, for the sake of contradiction, that \eqref{eq:superlevel-claim} fails.
Choose, by the layer-cake representation, the unique \(t_0\ge 0\) such that
\[
A:=\{u_F>t_0\}\quad\text{satisfies}\quad |A|=|\Sigma|;
\]
then \(A\) and \(\Sigma\) do not coincide a.e., so both \(A\setminus\Sigma\) and
\(\Sigma\setminus A\) have positive Lebesgue measure (and in fact positive
weighted measure, since \(e^{-\pi|z|^2}>0\) everywhere).

We now perform the bathtub construction explicitly. By inner regularity of the
Lebesgue measure, choose an open subset
\(S'\subset A\setminus\Sigma\) and an open subset
\(T'\subset \Sigma\setminus A\) with \(|S'|,|T'|>0\). Replacing each by a
sufficiently small open subset, we may further assume
\(\int_{S'}d\mu,\,\int_{T'}d\mu<\delta_0/4\). Since the Lebesgue measure
restricted to a fixed open set has the intermediate-value property, we may
choose subsets \(S\subset S'\) and \(T\subset T'\) with the common Lebesgue
measure
\[
|S|=|T|=\tfrac12\min(|S'|,|T'|)>0.
\]
Then
\[
\int_\C (\mathbf 1_S+\mathbf 1_T)\,e^{-\pi |z|^2}\,dz<\tfrac{\delta_0}{2}<\delta_0.
\]
Define
\[
\Sigma':=(\Sigma\setminus T)\cup S.
\]
Then \(|\Sigma'|=|\Sigma|\) and
\[
\int_\C |\mathbf 1_\Sigma-\mathbf 1_{\Sigma'}|e^{-\pi |z|^2}\,dz
=\int_\C(\mathbf 1_S+\mathbf 1_T)e^{-\pi|z|^2}\,dz<\delta_0.
\]
On the other hand, by construction,
\[
u_F(z)>t_0 \quad \text{on } S\subset A,
\qquad
u_F(z)\le t_0 \quad \text{on } T\subset \Sigma\setminus A.
\]
Moreover, since \(u_F\) is continuous and \(S\subset A=\{u_F>t_0\}\) is open
with positive measure, we may by shrinking \(S\) further assume that there
exists \(\eta>0\) with \(u_F\ge t_0+\eta\) on \(S\). It then follows that
\[
\int_S u_F\,dA \ge (t_0+\eta)\,|S|,
\qquad
\int_T u_F\,dA \le t_0\,|T|=t_0\,|S|,
\]
and hence
\[
\int_{\Sigma'} u_F\,dA
=
\int_\Sigma u_F\,dA + \int_S u_F\,dA-\int_T u_F\,dA
\ge \int_\Sigma u_F\,dA + \eta\,|S|
>
\int_\Sigma u_F\,dA.
\]
Since \(F\) is admissible for the variational problem defining
\(\lambda_{1,\varphi}(\Sigma')\),
\[
\lambda_{1,\varphi}(\Sigma')
\ge \int_{\Sigma'} u_F\,dA
>
\int_\Sigma u_F\,dA
=
\lambda_{1,\varphi}(\Sigma),
\]
contradicting the local maximality of \(\Sigma\). This proves
\eqref{eq:superlevel-claim}.

In particular, since \(u_F\) is real-analytic outside the (discrete) zero set of
\(F\), \(u_F\) is constant on \(\partial \Sigma\) in the sense that there exists
\(c_\Sigma\ge 0\) with
\begin{equation}\label{eq:boundary-level-set-new}
u_F(z)=c_\Sigma
\qquad \text{for } z\in \partial \Sigma.
\end{equation}
We will see in Step 2 that \(F\) does not vanish on \(\overline\Sigma\), and in
particular \(c_\Sigma>0\).

We collect the conclusions of this step in a self-contained statement that we
will reuse below.

\begin{lemma}[Superlevel-set characterization]\label{lem:bathtub-superlevel-from-localmax}
Let \(\Sigma\in\mathcal M_a\) be a local maximizer for \(\lambda_{1,\varphi}\)
in the sense of Definition~\ref{def:local-maximizer}, and let \(F\in
\mathcal F^2(\C)\) be a normalized first eigenfunction associated with
\(\Sigma\). Then there exists \(t_0>0\) such that
\begin{equation}\label{eq:superlevel-strong}
\Sigma=\{z\in\C\,:\, u_F(z)>t_0\}\quad
\text{up to a Lebesgue-null set,}
\end{equation}
and \(\Sigma\) coincides up to a null set with the unique (up to null sets)
maximizer of the bathtub problem
\(\sup\{\int_E u_F\,dA\,:\,|E|=a\}\) provided by Lemma~\ref{lem:bathtub}
applied to \(g=u_F\).
\end{lemma}

\begin{proof}
The first part is exactly \eqref{eq:superlevel-claim}; the strict positivity
\(t_0>0\) follows from the fact, derived in Step~2 below, that \(F\) cannot
vanish on \(\overline\Sigma\), so that \(c_\Sigma=t_0>0\). Uniqueness within
bathtub maximizers follows from the no-plateau remark below
Lemma~\ref{lem:bathtub}, since \(u_F\) is real-analytic on the open set
\(\{F\ne 0\}\supset\overline\Sigma\) and hence has no positive-measure level
sets there.
\end{proof}

\begin{lemma}[Boundary regularity at a local maximizer]\label{lem:boundary-reg}
Under the hypotheses of Lemma~\ref{lem:bathtub-superlevel-from-localmax},
choosing the open superlevel-set representative
\(\Sigma=\{u_F>t_0\}\), we have:
\begin{enumerate}
\item[\textnormal{(i)}] \(\Sigma\) is open and bounded (because
\(u_F(z)\to 0\) as \(|z|\to\infty\) and \(t_0>0\));
\item[\textnormal{(ii)}] \(F\) does not vanish on \(\overline\Sigma\), and
\(u_F\) is real-analytic on a neighborhood of \(\overline\Sigma\);
\item[\textnormal{(iii)}] \(\partial\Sigma\) is a real-analytic
\(1\)-manifold, except for a finite set
\(\Sigma_{\mathrm{sing}}\subset \partial\Sigma\) consisting of points where
\(\nabla u_F\) vanishes; in particular, \(\partial\Sigma\) is locally
Lipschitz, has finite \(\mathcal H^1\)-measure, and
\(\mathcal H^1(\Sigma_{\mathrm{sing}})=0\);
\item[\textnormal{(iv)}] the divergence theorem
\(\int_\Sigma \mathrm{div}\,X\,dA=\int_{\partial\Sigma}X\cdot\nu\,
d\mathcal H^1\) holds for every \(C^1\)-vector field \(X\) defined on a
neighborhood of \(\overline\Sigma\).
\end{enumerate}
\end{lemma}

\begin{proof}
Statement (i) is immediate, since \(u_F(z)=|F(z)|^2 e^{-\pi|z|^2}\) tends to
\(0\) at infinity (as \(F\) is entire and the Gaussian factor decays
super-exponentially against any polynomial growth of \(|F|^2\) along any
sequence \(|z|\to\infty\), since \(F\in\mathcal F^2\) is in particular bounded
along any horizontal line by \cite[Ch.~1]{Folland}). The non-vanishing in
(ii) is established in Step~2 (more precisely, in
Lemma~\ref{lem:simplicity-local-max} below; the argument in Step~2 only uses
(i) of the present lemma plus \eqref{eq:superlevel-strong}). Once \(F\) is
non-vanishing on \(\overline\Sigma\), \(u_F\) is real-analytic on a
neighborhood of \(\overline\Sigma\). Statement (iii) is the standard local
structure theorem for level sets of non-constant real-analytic functions
% TODO: bib key Krantz-Parks-Primer, Lojasiewicz
(see, e.g., \cite{Krantz-Parks-Primer,Lojasiewicz}), which gives that
\(\{u_F=t_0\}\) decomposes locally into finitely many real-analytic arcs
meeting at finitely many singular points (the zeros of \(\nabla u_F\) on
\(\{u_F=t_0\}\), of which there are finitely many on the compact set
\(\overline\Sigma\) by analyticity). Lipschitz regularity of the boundary
then follows by parametrizing each arc by arc length, and the
divergence-theorem statement (iv) follows from the Gauss--Green theorem on
Lipschitz domains
% TODO: bib key Evans-Gariepy
(cf.\ \cite{Evans-Gariepy}), using that the Lipschitz regular set
\(\partial\Sigma\setminus \Sigma_{\mathrm{sing}}\) carries
\(\mathcal H^1\)-almost every point of \(\partial\Sigma\).
\end{proof}

\begin{remark}\label{rem:reg-vs-hypothesis}
No boundary regularity is hypothesized on the local maximizer: the only input
is the measurability of \(\Sigma\), and the real-analyticity of
\(\partial\Sigma\) (away from finitely many points) is \emph{derived} from
the superlevel-set characterization combined with the entirety of \(F\) and
the Gaussian weight. Thus the divergence-theorem identity used in Step~3 is
fully justified by Lemma~\ref{lem:boundary-reg}(iv), even though \(\Sigma\)
is a priori only a measurable set.
\end{remark}

\medskip
\noindent
{\bf Step 2: simplicity of the first eigenvalue.}\label{step:simplicity}

In this step we prove the following lemma, which lies at the heart of the
remainder of the argument: it is what allows us, in Step~4, to conclude from
the fact that \(F'\) is a first eigenfunction that \(F'\) is a scalar multiple
of \(F\) and hence that \(F\) satisfies a linear ODE.

\begin{lemma}[Simplicity at a local maximizer]\label{lem:simplicity-local-max}
Let \(\Sigma\) be a local maximizer for \(\lambda_{1,\varphi}\) in the sense
of Definition~\ref{def:local-maximizer} and let \(F\in\mathcal F^2(\C)\) be a
normalized first eigenfunction associated with \(\Sigma\). Then the
eigenspace \(\ker(\mathcal L_\Sigma-\lambda_{1,\varphi}(\Sigma)\,\mathrm I)\)
is one-dimensional; that is, \(\lambda_{1,\varphi}(\Sigma)\) is simple, and
\(F\) is uniquely determined up to a unimodular constant. In particular,
\(F\) does not vanish on \(\overline\Sigma\), and \(F\) is bounded on \(\C\).
\end{lemma}

The proof uses a modular-deformation principle: any two first eigenfunctions
have the same modulus on \(\partial\Sigma\) up to a constant, and a holomorphic
function with constant modulus on a connected open set must itself be constant.
This argument plays a role also in the stability analysis of
\cite{Gomez-Guerra-Ramos-Tilli}.

\begin{proof}[Proof of Lemma~\ref{lem:simplicity-local-max}]
Suppose, for contradiction, that \(G\in \mathcal F^2(\C)\) is another first
eigenfunction, linearly independent of \(F\). Step 1 applied to \(G\) shows that
\(\Sigma\) is also a superlevel set of \(u_G\); hence there exists
\(c'_\Sigma\ge 0\) such that
\[
u_G(z)=c'_\Sigma
\qquad\text{for } z\in \partial \Sigma.
\]
Equivalently,
\[
|F(z)|^2 e^{-\pi |z|^2}=c_\Sigma,
\qquad
|G(z)|^2 e^{-\pi |z|^2}=c'_\Sigma,
\qquad z\in \partial \Sigma.
\]

We next observe that first eigenfunctions are zero-free in the interior of
\(\Sigma\). Indeed, by Step 1 we have \(u_F>c_\Sigma\) on the (open) interior of
\(\Sigma\). Since
\[
u_F=|F|^2 e^{-\pi |z|^2}
\]
and the Gaussian factor \(e^{-\pi|z|^2}\) is strictly positive everywhere, this
yields \(|F|>0\) throughout the interior of \(\Sigma\); the same applies to
\(|G|\). In particular, \(c_\Sigma>0\) and \(c'_\Sigma>0\). Moreover, since
\(F,G\) are entire and continuous, \(|F|>0\) on a neighborhood of any compact
subset of \(\Sigma^\circ\); by Step 1, we may take a slight enlargement of
\(\Sigma\) (still a superlevel set of \(u_F\)) on which \(|F|,|G|>0\), so we may
work on a connected open set \(U\Supset \Sigma\) where both \(F\) and \(G\) are
zero-free.

The ratio \(F/G\) is holomorphic and nowhere zero on \(U\), and the function
\[
\log \left|\frac{F}{G}\right|
=
\frac12 \log \frac{|F|^2}{|G|^2}
=
\frac12 \log \frac{u_F}{u_G}
\]
is harmonic on \(U\): it is the real part of any local branch of
\(\log(F/G)\), which is well defined on simply connected subsets, and harmonic
in any case as the logarithm of the modulus of a zero-free holomorphic
function. Note that the Gaussian factors cancel inside the logarithm, so on
\(\partial \Sigma\) by \eqref{eq:boundary-level-set-new} this function equals
the constant \(\frac12\log(c_\Sigma/c'_\Sigma)\). By the maximum principle
applied separately on each connected component of \(\Sigma\), the function
\(\log|F/G|\) is constant on each component. Hence
\[
\left|\frac{F}{G}\right| \equiv \text{constant}
\qquad \text{on each connected component of } \Sigma.
\]
A nonconstant holomorphic function is open, so a holomorphic function with
constant modulus on a connected open set must itself be constant. Therefore
\[
F=\gamma_j G
\qquad \text{on each connected component } \Sigma_j\text{ of } \Sigma.
\]
The relation \(F=\gamma_j G\) is an identity between two entire functions on a
set with an accumulation point; the analytic continuation principle gives
\(F=\gamma_j G\) on all of \(\C\) for any single \(j\). In particular all
\(\gamma_j\) coincide, and \(F\) and \(G\) are linearly dependent, contradicting
our assumption. This proves simplicity.

\medskip
\noindent
{\bf Step 3: a first-variation identity.}

By Step 2, the eigenfunction \(F\) is determined up to a unimodular constant.
We now derive a first-variation identity satisfied by \(F\), based on the fact
that \(\Sigma\) is a level set of \(u_F\) and that holomorphic vector fields are
divergence-free.

The function \(u_F=|F|^2e^{-\pi|z|^2}\) is real-analytic on
\(\C\setminus\{F=0\}\), since \(F\) is entire and the exponential factor is
real-analytic. By Step 2, \(F\) does not vanish in the interior of \(\Sigma\),
and by the argument of Step 4 below it does not vanish on \(\overline\Sigma\)
either; so \(u_F\) is in fact real-analytic on a neighborhood of
\(\overline\Sigma\). The boundary \(\partial\Sigma\) is therefore a level set
of a non-constant real-analytic function and so is the union of finitely many
real-analytic arcs, plus a possible finite set of singular points where
\(\nabla u_F\) vanishes (see, e.g., the standard local structure of
real-analytic varieties). In particular, \(\partial\Sigma\) is Lipschitz, and
\(C^1\) away from the critical set of \(u_F\); the divergence theorem applies
to \(\Sigma\) for any \(C^1\)-vector field on a neighborhood of
\(\overline\Sigma\), the singular boundary points forming a set of
one-dimensional Hausdorff measure zero.

Let \(\psi\) be holomorphic on an open neighborhood \(V\) of
\(\overline \Sigma\), and write
\[
\psi=u+iv
\]
with \(u,v\) real-valued. Consider the real vector field
\[
X_\psi:=(u,-v).
\]
Since \(\psi\) is holomorphic, the Cauchy--Riemann equations
\[
u_x=v_y,
\qquad
u_y=-v_x,
\]
imply
\[
\operatorname{div}X_\psi=u_x+(-v)_y=u_x-v_y=0
\qquad \text{in }V.
\]
Because \(u_F\equiv c_\Sigma\) on \(\partial\Sigma\), the divergence theorem on
\(\Sigma\) gives
\[
\int_\Sigma \langle \nabla u_F, X_\psi\rangle\,dA
=
\int_\Sigma \operatorname{div}(u_F X_\psi)\,dA
=
\int_{\partial \Sigma} u_F\, X_\psi\cdot \nu\, d\mathcal H^1
=
c_\Sigma\int_{\partial \Sigma} X_\psi\cdot \nu\, d\mathcal H^1
=
0,
\]
where the last equality uses
\(\int_{\partial\Sigma}X_\psi\cdot\nu\,d\mathcal H^1
=\int_\Sigma\operatorname{div}X_\psi\,dA=0\) by Cauchy--Riemann. Since
\(u_F=|F|^2e^{-\pi |z|^2}\) and
\(\nabla(e^{-\pi|z|^2})=-2\pi e^{-\pi|z|^2}(x,y)\), this becomes
\begin{equation}\label{eq:vector-field-identity-real-new}
0=
\int_\Sigma
\left(
\langle \nabla |F|^2,X_\psi\rangle
-
2\pi |F|^2 \langle z,X_\psi\rangle
\right)
\,d\mu.
\end{equation}

We now express \eqref{eq:vector-field-identity-real-new} in complex notation.
A direct computation, using \(\partial_z|F|^2=F'\overline F\) and
\(\partial_{\bar z}|F|^2=F\overline{F'}\), gives
\[
\partial_x |F|^2 = F'\overline F+\overline{F'}F = 2\Re(F'\overline F),
\qquad
\partial_y |F|^2 = i(F'\overline F-\overline{F'}F) = -2\Im(F'\overline F).
\]
Writing \(\psi=u+iv\) and recalling \(X_\psi=(u,-v)\),
\[
\langle \nabla |F|^2,X_\psi\rangle
=
2u\,\Re(F'\overline F)+2v\,\Im(F'\overline F)
=
2\,\Re\!\big(\overline{\psi}\,F'\overline F\big)
=
2\,\Re\!\big(\psi\,\overline{F'}F\big),
\]
where in the last equality we used \(\Re w=\Re\overline w\). Likewise, writing
\(z=x+iy\),
\[
\langle z,X_\psi\rangle
=
xu+y(-v)
=
xu-yv
=
\Re\big((x+iy)(u+iv)\big)
=
\Re(z\psi).
\]
Substituting these expressions into
\eqref{eq:vector-field-identity-real-new}, we obtain
\[
0=
\Re\!\int_\Sigma
\psi(z)\Big(\overline{F'(z)}\,F(z)-\pi z\,|F(z)|^2\Big)\,d\mu(z).
\]
Replacing \(\psi\) by \(i\psi\) (which is also holomorphic on \(V\)) gives the
analogous identity for the imaginary part:
\[
0=
\Im\!\int_\Sigma
\psi(z)\Big(\overline{F'(z)}\,F(z)-\pi z\,|F(z)|^2\Big)\,d\mu(z).
\]
Combining the two,
\begin{equation}\label{eq:orthogonality-main-new}
0=
\int_\Sigma
\psi(z)\,F(z)\Big(\overline{F'(z)}-\pi z\,\overline{F(z)}\Big)\,d\mu(z),
\end{equation}
which holds for every holomorphic \(\psi\) on a neighborhood of
\(\overline \Sigma\). Equivalently, after factoring out \(F(z)\) inside the
integrand,
\begin{equation*}
0=
\int_\Sigma
\psi(z)\Big(\overline{F'(z)}\,F(z)-\pi z\,|F(z)|^2\Big)\,d\mu(z).
\end{equation*}
The form \eqref{eq:orthogonality-main-new} is the one used in Step 4.

\medskip
\noindent
{\bf Step 4: \(F'\) is again a first eigenfunction.}

This step is the crux of the argument and is the main new ingredient compared
to classical Faber--Krahn rigidity proofs. Rather than working with general
test functions in the eigenvalue equation, we exploit the divergence identity
\eqref{eq:orthogonality-main-new} together with a clever choice of test
function involving the reproducing kernel of the Fock space. The end result is
that \emph{the derivative} \(F'\) of any first eigenfunction is itself a first
eigenfunction; since the eigenspace is one-dimensional by Step 2, this forces
a first-order ODE on \(F\), and hence pins down \(F\) up to translation as a
Gaussian.

For every \(w\in \C\), set
\[
K_w(z):=e^{\pi \overline w z};
\]
this is the reproducing kernel of \(\mathcal F^2(\C)\) at the point \(w\), in
the sense that
\begin{equation}\label{eq:reproducing-kernel}
\int_\C H(z)\,e^{\pi w \overline z}\,d\mu(z)=H(w),
\qquad H\in\mathcal F^2(\C);
\end{equation}
see, e.g., \cite[Ch.~1]{Folland} or \cite{Bargmann}. Equivalently,
\(\langle H,K_w\rangle_{\mathcal F^2}=H(w)\).

By Step 1, after replacing \(\Sigma\) by the corresponding superlevel-set
representative if necessary, we may suppose that
\[
\Sigma=\{u_F>c_\Sigma\}.
\]
In particular, \(u_F\ge c_\Sigma>0\) on \(\overline \Sigma\). Since
\(u_F(z)=|F(z)|^2e^{-\pi |z|^2}\), it follows that \(F\) has no zeros on
\(\overline \Sigma\). By continuity of \(F\) and the compactness of
\(\overline \Sigma\), there exists an open neighborhood
\(V\supset \overline \Sigma\) on which \(F\) has no zeros. Hence, for every
\(w\in \C\), the function
\[
\psi_w(z):=\frac{K_w(z)}{F(z)}=\frac{e^{\pi\overline wz}}{F(z)}
\]
is holomorphic on \(V\). Applying \eqref{eq:orthogonality-main-new} with
\(\psi=\psi_w\), so that
\(\psi_w(z)F(z)=e^{\pi\overline w z}\), we obtain
\[
0=
\int_\Sigma
e^{\pi \overline wz}
\Big(\overline{F'(z)}-\pi z\,\overline{F(z)}\Big)\,d\mu(z),
\]
or, separating the two terms,
\begin{equation}\label{eq:key-kernel-identity-new}
\int_\Sigma
e^{\pi \overline wz}\overline{F'(z)}\,d\mu(z)
=
\int_\Sigma
\pi z\,e^{\pi \overline wz}\overline{F(z)}\,d\mu(z).
\end{equation}

We now relate the right-hand side of \eqref{eq:key-kernel-identity-new} to
\(\overline{F'(w)}\) using the eigenvalue equation
\eqref{eq:eigenvalue-equation} together with the reproducing identity
\eqref{eq:reproducing-kernel}. Choosing
\[
H(z):=z\,e^{\pi\overline w z}\in \mathcal F^2(\C)
\]
(the product of the kernel \(K_w\) by the polynomial \(z\); both factors
preserve membership in \(\mathcal F^2\)) in \eqref{eq:eigenvalue-equation}, we
obtain
\begin{equation}\label{eq:eigvalue-applied}
\int_\Sigma z\, e^{\pi \overline wz}\,\overline{F(z)}\,d\mu(z)
=
\lambda_{1,\varphi}(\Sigma)\int_\C z\, e^{\pi \overline wz}\,\overline{F(z)}\,d\mu(z).
\end{equation}
On the other hand, taking complex conjugates in \eqref{eq:reproducing-kernel}
applied to \(F\) at the point \(w\), one finds
\[
\int_\C e^{\pi \overline wz}\,\overline{F(z)}\,d\mu(z)=\overline{F(w)}.
\]
Both sides depend on \(w\) through the antiholomorphic variable
\(\overline w\); we differentiate with respect to \(\overline w\), exchanging
differentiation and integration on the left. The interchange is legitimate by
the dominated convergence theorem: locally in \(w\), the integrand
\(\partial_{\overline w}(e^{\pi\overline wz}\overline{F(z)})
=\pi z\,e^{\pi\overline wz}\overline{F(z)}\) is dominated by an
\(L^1(d\mu)\) function on \(\C\), as a consequence of the Cauchy--Schwarz
inequality and the fact that \(F\in \mathcal F^2(\C)\) and the Fock-space norm
of \(zK_w\) is locally bounded. We conclude
\[
\pi\int_\C z\, e^{\pi \overline wz}\,\overline{F(z)}\,d\mu(z)
=
\overline{F'(w)}.
\]
Multiplying \eqref{eq:eigvalue-applied} by \(\pi\) and substituting,
\[
\pi\int_\Sigma z\, e^{\pi \overline wz}\,\overline{F(z)}\,d\mu(z)
=
\lambda_{1,\varphi}(\Sigma)\,\overline{F'(w)}.
\]
Combining this with \eqref{eq:key-kernel-identity-new},
\[
\int_\Sigma e^{\pi \overline wz}\,\overline{F'(z)}\,d\mu(z)
=
\lambda_{1,\varphi}(\Sigma)\,\overline{F'(w)},
\qquad w\in \C.
\]
Taking complex conjugates, this is exactly the eigenvalue equation
\eqref{eq:eigenvalue-equation} applied with \(F\) replaced by \(F'\) (and an
arbitrary \(K_w\) playing the role of the test function). Since the linear
span of \(\{K_w\}_{w\in\C}\) is dense in \(\mathcal F^2(\C)\), the identity
extends to every \(H\in\mathcal F^2(\C)\), and we conclude that \(F'\) is a
first eigenfunction of the localization operator on \(\Sigma\), with eigenvalue
\(\lambda_{1,\varphi}(\Sigma)\).

By the simplicity established in Step 2,
\[
F'=\theta F
\]
for some constant \(\theta\in \C\). Solving this ODE gives
\[
F(z)=C e^{\theta z},
\]
for some nonzero constant \(C\).

\medskip
\noindent
{\bf Step 5: reduction to the constant eigenfunction.}

We now use the Weyl invariance of the Bargmann--Fock setting to reduce the
analysis to the case \(F\equiv 1\). Recall from \eqref{eq:Ta} that, for
\(a\in\C\), the Weyl translation
\[
(T_a F)(z)= e^{\pi a z-\frac{\pi}{2}|a|^2}\,F(z-a)
\]
is unitary on \(\mathcal F^2(\C)\) and satisfies
\(u_{T_aF}(z)=u_F(z-a)\). In particular, \(T_a\) maps eigenfunctions of the
localization operator on \(\Sigma\) to eigenfunctions of the localization
operator on the translated set \(\Sigma+a\), with the same eigenvalue, and the
quantity \(\int_\Sigma u_F\,dA\) is invariant under the simultaneous
substitutions \(F\mapsto T_aF\) and \(\Sigma\mapsto \Sigma+a\).

We claim that we may choose \(a\in\C\) so that \(T_{-a}F\) is a constant
function. Indeed, by Step 4, \(F(z)=Ce^{\theta z}\) for some \(C,\theta\in\C\)
with \(C\ne 0\). Compute, using the identity
\(u_{T_{-a}F}(z)=u_F(z+a)\):
\[
u_{T_{-a}F}(z)=|F(z+a)|^2 e^{-\pi|z+a|^2}
=|C|^2 e^{2\Re(\theta(z+a))}\,e^{-\pi|z+a|^2}.
\]
Expanding,
\[
2\Re(\theta(z+a))-\pi|z+a|^2
=
2\Re(\theta z)+2\Re(\theta a)-\pi|z|^2-2\pi\Re(\overline a z)-\pi|a|^2.
\]
The linear terms in \(z\) on the right cancel if and only if
\(2\Re(\theta z)=2\pi\Re(\overline a z)\) for every \(z\in\C\), which is the
condition \(\overline a=\theta/\pi\), i.e.
\[
a:=\frac{\overline\theta}{\pi}.
\]
With this choice,
\[
u_{T_{-a}F}(z)=|C'|^2 e^{-\pi|z|^2},
\qquad
|C'|^2:=|C|^2 e^{2\Re(\theta a)-\pi|a|^2}=|C|^2 e^{|\theta|^2/\pi}.
\]
Since \(T_{-a}F\) is entire and \(|T_{-a}F(z)|^2 e^{-\pi|z|^2}=|C'|^2
e^{-\pi|z|^2}\), it follows that \(|T_{-a}F|\equiv|C'|\) is constant; a
holomorphic function with constant modulus is constant, hence
\(T_{-a}F\equiv C''\) for some \(C''\in\C\) with \(|C''|=|C'|\). Normalizing,
we may absorb the unimodular constant into the eigenfunction and assume
\[
F\equiv 1
\]
after the change \(\Sigma\rightsquigarrow \Sigma-a\) (and the corresponding
unitary substitution \(F\rightsquigarrow T_{-a}F\)).

In this normalization the ground state \(h_0\) (corresponding to
\(F\equiv 1\) in Bargmann--Fock coordinates) is an eigenfunction of the
localization operator of \(\Sigma\). Since \(\Sigma\) is a level set of
\(u_F(z) = e^{-\pi |z|^2}\), and the latter is a strictly radial function of
\(|z|\), the set \(\Sigma\) is itself radial; in particular, it is either a
centered disk or a finite union of centered annuli. The next step rules out
the annular configurations and identifies \(\Sigma\) with the disk.

\medskip
\noindent
{\bf Step 6: exclusion of annuli and conclusion.}

It remains to verify that, among radial sets of prescribed measure, the only
local maximizer is the centered disk. Although this is essentially automatic
after Step 5, we record the standard rearrangement argument since it makes
explicit \emph{why} no annular local maximizer can exist; this provides the
endpoint matching the global Faber--Krahn theorem of \cite{Nicola-Tilli}.

If \(\Sigma\) is radial, write
\[
\widetilde \Sigma:=\{r>0:\ re^{it}\in \Sigma \text{ for some } t\in [0,2\pi)\}.
\]
By Step 1 and the Lipschitz regularity established in Step 3,
\(\widetilde \Sigma\) is open and is a finite union of disjoint open intervals:
\[
\widetilde \Sigma=\bigcup_{j=1}^N (a_j,b_j),
\qquad
0\le a_1<b_1\le a_2<b_2\le \cdots\le a_N<b_N.
\]
Then, using \(\int re^{-\pi r^2}\,dr=-\frac{1}{2\pi}e^{-\pi r^2}\),
\[
\int_\Sigma e^{-\pi |z|^2}\,dA(z)
=
2\pi \int_{\widetilde \Sigma} r e^{-\pi r^2}\,dr
=
\sum_{j=1}^N \big(e^{-\pi a_j^2}-e^{-\pi b_j^2}\big),
\]
under the constraint
\[
|\Sigma|
=
2\pi \int_{\widetilde \Sigma} r\,dr
=
\pi \sum_{j=1}^N (b_j^2-a_j^2).
\]
We are therefore led to the optimization problem
\begin{equation}\label{eq:radial-opt}
\text{maximize}\qquad
\sum_{j=1}^N \big(e^{-\pi a_j^2}-e^{-\pi b_j^2}\big)
\qquad \text{subject to}\qquad
\sum_{j=1}^N (b_j^2-a_j^2)=c,
\end{equation}
where \(c=|\Sigma|/\pi\) is fixed. Substituting \(s_j:=a_j^2\),
\(t_j:=b_j^2\) (with \(0\le s_1<t_1\le s_2<\cdots\)), and writing
\(\Phi(x):=e^{-\pi x}\), \eqref{eq:radial-opt} becomes
\[
\text{maximize}\qquad
\sum_{j=1}^N \big(\Phi(s_j)-\Phi(t_j)\big)
\qquad \text{subject to}\qquad
\sum_{j=1}^N (t_j-s_j)=c.
\]
The function \(\Phi\) is strictly convex and strictly decreasing on
\([0,\infty)\); its derivative \(-\Phi'(\sigma)=\pi e^{-\pi\sigma}\) is strictly
decreasing in \(\sigma\). For any \(t>s\ge 0\),
\[
\Phi(s)-\Phi(t)
=
\int_s^t (-\Phi'(\sigma))\,d\sigma,
\]
so
\[
\sum_{j=1}^N\big(\Phi(s_j)-\Phi(t_j)\big)
=
\int_{\widetilde\Sigma_2}(-\Phi'(\sigma))\,d\sigma,
\qquad
\widetilde\Sigma_2:=\bigcup_{j=1}^N (s_j,t_j).
\]
The set \(\widetilde\Sigma_2\subset[0,\infty)\) has Lebesgue measure equal to
\(c\), and the integrand \(-\Phi'\) is strictly decreasing in \(\sigma\). The
bathtub principle (or, equivalently, a direct rearrangement argument) shows
that, among all measurable subsets of \([0,\infty)\) of total Lebesgue measure
\(c\), the integral \(\int_{\widetilde\Sigma_2}(-\Phi'(\sigma))\,d\sigma\) is
strictly maximized by the leftmost interval \((0,c)\). This corresponds to
\(N=1\), \(s_1=0\), \(t_1=c\), i.e.
\[
\widetilde \Sigma=(0,R),
\qquad R:=\sqrt c,
\]
and equivalently \(\Sigma\) is the centered disk \(D_R\) with
\(|D_R|=|\Sigma|\). Any annular configuration (\(N\ge 2\), or \(N=1\) with
\(s_1>0\)) produces a strictly smaller value of \(\int_\Sigma u_F\,dA\), and is
therefore not a local maximizer.

We conclude that every local maximizer of \(\lambda_{1,\varphi}\) under the
measure constraint is, up to translation, a disk.
\end{proof}


% ==============================================================
% Section 7
% ==============================================================
\section{Global maximizers for the Faber--Krahn problem in the Gaussian STFT setting}

In this section we prove existence of global maximizers for the Gaussian
time-frequency concentration problem under a measure constraint, and then
combine that existence result with the local rigidity theorem from the previous
section to recover the disk as the unique maximizer, up to translation. The
strategy is by now classical in time-frequency analysis: a profile decomposition
in the spirit of P.-L.~Lions' concentration-compactness, transplanted to the
Fock space via the Bargmann transform, dichotomizes the loss of compactness of a
maximizing sequence into the action of the Weyl translation group. A vanishing
remainder is then ruled out by an iterative extraction whose maximality forces
all the energy to remain captured by the extracted profiles. Closely related
arguments have appeared in the time-frequency setting in
\cite{Nicola-Romero-Trapasso, Marceca-Romero-Speckbacher, Fulsche-Hagger,
Samuelsen}, and the disk-rigidity result we recover is the celebrated theorem
of Nicola--Tilli \cite{Nicola-Tilli, Nicola-Tilli-2}.

Since the Bargmann transform is unitary, the variational problem may be written
on the Fock space side as
\[
M(s):=\sup_{|\Omega|=s}\lambda_{1,\varphi}(\Omega)
=
\sup_{\|F\|_{\mathcal F^2}=1} I_F(s).
\]
Here, and throughout the section, we exploit the translation invariance
\eqref{eq:I-translation-global} together with the pointwise identity
\eqref{eq:u-translation-global} (cf.\ also \eqref{eq:uF-translate}), the
basic regularity properties recorded in Lemma~\ref{lem:basic-properties-global},
and the elementary properties of the Fock space gathered in
Lemma~\ref{lem:Fock-basic}.

We begin with a Weyl-orthogonality lemma that captures the asymptotic
decoupling of two functions whose translation parameters drift apart. This
plays the role of \emph{Pythagorean orthogonality at infinity} on the Fock
space.

\begin{lemma}\label{lem:weyl-orthogonality-global}
Let \(F,G\in \mathcal F^2(\C)\), and let \(\{a_k\}\subset \C\) satisfy
\(|a_k|\to \infty\). Then
\[
T_{a_k}G \rightharpoonup 0
\qquad\text{weakly in }\mathcal F^2(\C),
\]
and in particular
\[
\langle F,T_{a_k}G\rangle_{\mathcal F^2}\to 0.
\]
Consequently,
\[
\|F+T_{a_k}G\|_{\mathcal F^2}^2
=
\|F\|_{\mathcal F^2}^2+\|G\|_{\mathcal F^2}^2+o(1).
\]
More generally, if \(|a_k-b_k|\to\infty\), then
\[
\langle T_{a_k}F,T_{b_k}G\rangle_{\mathcal F^2}\to 0.
\]
\end{lemma}

\begin{proof}
Using unitarity of the Weyl operators \(T_a\) on \(\mathcal F^2(\C)\) (cf.\
\eqref{eq:Ta}),
\[
\langle F,T_{a_k}G\rangle_{\mathcal F^2}
=
\langle T_{-a_k}F,G\rangle_{\mathcal F^2}.
\]
Now \(T_{-a_k}F\) is bounded in \(\mathcal F^2\) by unitarity, and the
identity \eqref{eq:u-translation-global} (equivalently
\eqref{eq:uF-translate}) yields
\[
u_{T_{-a_k}F}(z)=u_F(z+a_k)\qquad (z\in\C),
\]
so that the densities \(u_{T_{-a_k}F}\) drift uniformly to infinity.
Equivalently, since \(u_F\in C_0(\C)\) by Lemma~\ref{lem:Fock-basic}(iii), for
every compact \(K\subset \C\),
\[
\sup_{z\in K}|T_{-a_k}F(z)|e^{-\pi |z|^2/2}
=
\sup_{z\in K}u_{T_{-a_k}F}(z)^{1/2}
=
\sup_{z\in K}u_F(z+a_k)^{1/2}\to 0.
\]
Hence \(T_{-a_k}F\to 0\) locally uniformly. Combining this with the
boundedness of \(\{T_{-a_k}F\}\) in \(\mathcal F^2(\C)\), and using the fact
that the linear span of the reproducing kernels \(\{K_w\}_{w\in\C}\) is dense
in \(\mathcal F^2\), we conclude that \(T_{-a_k}F\rightharpoonup 0\) weakly in
\(\mathcal F^2\); in particular \(\langle T_{-a_k}F,G\rangle\to 0\). This
proves the first claim.

For the norm decoupling, expand the square:
\[
\|F+T_{a_k}G\|_{\mathcal F^2}^2
=
\|F\|_{\mathcal F^2}^2+\|T_{a_k}G\|_{\mathcal F^2}^2
+2\Re\langle F,T_{a_k}G\rangle_{\mathcal F^2}
=
\|F\|_{\mathcal F^2}^2+\|G\|_{\mathcal F^2}^2+o(1),
\]
where we used unitarity of \(T_{a_k}\) and the cross-term decay just proved.

Finally, the general case reduces to the first by writing
\[
\langle T_{a_k}F,T_{b_k}G\rangle_{\mathcal F^2}
=
\langle F,T_{-a_k}T_{b_k}G\rangle_{\mathcal F^2}
=
e^{i\theta_k}\langle F,T_{b_k-a_k}G\rangle_{\mathcal F^2},
\]
for some real phase \(\theta_k\) coming from the cocycle of the Weyl
representation; since this phase factor has modulus one, the conclusion
follows from \(|b_k-a_k|\to\infty\).
\end{proof}

The next lemma compares the density of a finite sum of widely separated
translates with the density of each individual translate. The precise statement
is what makes the profile decomposition compatible with the variational
quantities \(I_F(s)\), since these depend only on the densities and not on
the phases of the underlying functions.

\begin{lemma}\label{lem:pointwise-decoupling-global}
Fix \(J\ge 1\), let \(G^{(1)},\dots,G^{(J)}\in \mathcal F^2(\C)\), and let
\(\{a_k^{(j)}\}_{k\ge 1}\subset \C\) for \(1\le j\le J\) satisfy
\[
|a_k^{(j)}-a_k^{(\ell)}|\to\infty
\qquad\text{whenever }j\neq \ell.
\]
For each \(k\), set
\[
S_k:=\sum_{j=1}^J T_{a_k^{(j)}}G^{(j)}.
\]
Then the following hold.
\begin{enumerate}
\item[(i)] For every \(R>0\) and every \(1\le j\le J\),
\[
\sup_{z\in D_R(a_k^{(j)})}
\big|u_{S_k}(z)-u_{T_{a_k^{(j)}}G^{(j)}}(z)\big|
\to 0
\qquad (k\to\infty).
\]

\item[(ii)] For every \(R>0\), if
\[
E_{k,R}:=\C\setminus \bigcup_{j=1}^J D_R(a_k^{(j)}),
\]
then
\[
\limsup_{k\to\infty}\sup_{z\in E_{k,R}}u_{S_k}(z)
\le
\left(
\sum_{j=1}^J
\sup_{\zeta\in \C\setminus D_R} u_{G^{(j)}}(\zeta)^{1/2}
\right)^2.
\]
In particular, given \(\varepsilon>0\), there exists \(R_\varepsilon>0\) such that
for every \(R\ge R_\varepsilon\),
\[
\sup_{z\in E_{k,R}}u_{S_k}(z)\le \varepsilon
\]
for all sufficiently large \(k\).
\end{enumerate}
\end{lemma}

\begin{proof}
Throughout the proof we use repeatedly the identity
\[
u_{T_a F}(z)=u_F(z-a),
\qquad z,a\in \C,
\]
which is \eqref{eq:u-translation-global} (equivalently
\eqref{eq:uF-translate}); together with the elementary subadditivity
\[
u_{F+G}(z)^{1/2}\le u_F(z)^{1/2}+u_G(z)^{1/2},
\qquad z\in\C,
\]
which follows from the definition \(u_F=|F|^2 e^{-\pi |z|^2}\) and the
triangle inequality applied to \(|F+G|\le |F|+|G|\).

\medskip
\noindent\emph{Proof of (i).}
Fix \(R>0\) and \(j\in\{1,\dots,J\}\). Write
\[
S_k=T_{a_k^{(j)}}G^{(j)}+H_{k,j},
\qquad
H_{k,j}:=\sum_{\ell\neq j} T_{a_k^{(\ell)}}G^{(\ell)}.
\]
For every \(z\in \C\), the elementary inequality
\(\big||a+b|^2-|a|^2\big|\le 2|a|\,|b|+|b|^2\) yields
\begin{align*}
\big|u_{S_k}(z)-u_{T_{a_k^{(j)}}G^{(j)}}(z)\big|
&=
e^{-\pi |z|^2}
\big||T_{a_k^{(j)}}G^{(j)}(z)+H_{k,j}(z)|^2
-|T_{a_k^{(j)}}G^{(j)}(z)|^2\big| \\
&\le
2\,u_{T_{a_k^{(j)}}G^{(j)}}(z)^{1/2}u_{H_{k,j}}(z)^{1/2}
+u_{H_{k,j}}(z).
\end{align*}
By Lemma~\ref{lem:Fock-basic}(ii),
\[
u_{T_{a_k^{(j)}}G^{(j)}}(z)\le \|G^{(j)}\|_{\mathcal F^2}^2
\qquad \text{for every }z\in\C\text{ and every }k.
\]
Thus, to prove (i) it is enough to show that
\[
\sup_{z\in D_R(a_k^{(j)})} u_{H_{k,j}}(z)\to 0.
\]

For \(z\in D_R(a_k^{(j)})\) and \(\ell\neq j\),
\[
z-a_k^{(\ell)}\in D_R(a_k^{(j)}-a_k^{(\ell)}).
\]
Since \(|a_k^{(j)}-a_k^{(\ell)}|\to\infty\),
\[
\operatorname{dist}\big(D_R(a_k^{(j)}-a_k^{(\ell)}),0\big)
\ge |a_k^{(j)}-a_k^{(\ell)}|-R\to\infty.
\]
Because \(u_{G^{(\ell)}}\in C_0(\C)\) by Lemma~\ref{lem:Fock-basic}(iii),
\[
\sup_{z\in D_R(a_k^{(j)})} u_{T_{a_k^{(\ell)}}G^{(\ell)}}(z)
=
\sup_{\zeta\in D_R(a_k^{(j)}-a_k^{(\ell)})} u_{G^{(\ell)}}(\zeta)
\xrightarrow[k\to\infty]{} 0.
\]
The subadditivity of \(u_F^{1/2}\) yields
\[
u_{H_{k,j}}(z)^{1/2}
\le
\sum_{\ell\neq j} u_{T_{a_k^{(\ell)}}G^{(\ell)}}(z)^{1/2},
\]
so the preceding convergence implies
\[
\sup_{z\in D_R(a_k^{(j)})} u_{H_{k,j}}(z)^{1/2}\to 0.
\]
This proves (i).

\medskip
\noindent\emph{Proof of (ii).}
Fix \(R>0\) and \(z\in E_{k,R}\). By definition, \(z\notin D_R(a_k^{(j)})\)
for every \(j\), and hence \(|z-a_k^{(j)}|\ge R\). Therefore
\[
u_{T_{a_k^{(j)}}G^{(j)}}(z)
=
u_{G^{(j)}}(z-a_k^{(j)})
\le
\sup_{\zeta\in \C\setminus D_R}u_{G^{(j)}}(\zeta).
\]
Using again the subadditivity of \(u_F^{1/2}\),
\[
u_{S_k}(z)^{1/2}
\le
\sum_{j=1}^J u_{T_{a_k^{(j)}}G^{(j)}}(z)^{1/2}
\le
\sum_{j=1}^J \sup_{\zeta\in \C\setminus D_R}u_{G^{(j)}}(\zeta)^{1/2}.
\]
Squaring yields
\[
u_{S_k}(z)
\le
\left(
\sum_{j=1}^J \sup_{\zeta\in \C\setminus D_R}u_{G^{(j)}}(\zeta)^{1/2}
\right)^2.
\]
Taking the supremum over \(z\in E_{k,R}\) proves the displayed estimate in
(ii). Since each \(u_{G^{(j)}}\in C_0(\C)\), the right-hand side tends to
\(0\) as \(R\to\infty\); the final claim follows by choosing
\(R_\varepsilon\) so that the displayed bracket is at most
\(\sqrt{\varepsilon}\).
\end{proof}

The next lemma is a robust uniform-density perturbation estimate. Concretely,
it says that adding to a bounded sequence in \(\mathcal F^2(\C)\) a residue
whose density vanishes uniformly leaves the density essentially unchanged in
\(L^\infty\). This will be applied with \(S_k\) the truncated profile sum and
\(R_k\) the residual of the profile decomposition.

\begin{lemma}\label{lem:density-perturbation-global}
Let \(\{S_k\}\subset \mathcal F^2(\C)\) be bounded and let
\(\{R_k\}\subset \mathcal F^2(\C)\) satisfy
\[
\|u_{R_k}\|_{L^\infty(\C)}\to 0.
\]
Then
\[
\|u_{S_k+R_k}-u_{S_k}\|_{L^\infty(\C)}\to 0.
\]
More precisely, if \(\sup_k \|S_k\|_{\mathcal F^2}\le M\), then
\[
\|u_{S_k+R_k}-u_{S_k}\|_{L^\infty(\C)}
\le
2M\,\|u_{R_k}\|_{L^\infty(\C)}^{1/2}
+\|u_{R_k}\|_{L^\infty(\C)}.
\]
\end{lemma}

\begin{proof}
For every \(z\in \C\),
\begin{align*}
\big|u_{S_k+R_k}(z)-u_{S_k}(z)\big|
&=
e^{-\pi |z|^2}\big||S_k(z)+R_k(z)|^2-|S_k(z)|^2\big| \\
&\le
2\,u_{S_k}(z)^{1/2}u_{R_k}(z)^{1/2}
+u_{R_k}(z),
\end{align*}
again by the elementary inequality
\(\big||a+b|^2-|a|^2\big|\le 2|a|\,|b|+|b|^2\). By
Lemma~\ref{lem:Fock-basic}(ii),
\[
u_{S_k}(z)\le \|S_k\|_{\mathcal F^2}^2\le M^2
\qquad (z\in\C).
\]
Hence
\[
\big|u_{S_k+R_k}(z)-u_{S_k}(z)\big|
\le
2M\,\|u_{R_k}\|_{L^\infty}^{1/2}
+\|u_{R_k}\|_{L^\infty}.
\]
Taking the supremum over \(z\) yields the claim.
\end{proof}

The next lemma supplies the soft input on the function \(s\mapsto I_F(s)\)
that will be used at the end of the variational argument. It is essentially a
consequence of the bathtub principle and the basic regularity of \(u_F\),
already noted in Lemma~\ref{lem:basic-properties-global}.

\begin{lemma}\label{lem:I-continuity-global}
Let \(F\in \mathcal F^2(\C)\). Then the function
\[
s\mapsto I_F(s)
\]
is finite, nondecreasing, and continuous on \([0,\infty)\).
\end{lemma}

\begin{proof}
By Lemma~\ref{lem:basic-properties-global},
\(u_F\in C_0(\C)\cap L^1(\C)\), so \(I_F(s)\) is finite for every \(s\ge 0\)
and bounded above by \(\|F\|_{\mathcal F^2}^2\). Monotonicity is immediate
from the definition of \(I_F(s)\) as a supremum over \(|\Omega|=s\): if
\(s_1\le s_2\), every admissible set for \(s_1\) can be enlarged to a set of
measure \(s_2\) without decreasing the integral, since \(u_F\ge 0\).

To prove continuity, write \(u:=u_F\) and let \(u^*\) denote the decreasing
rearrangement of \(u\) on \([0,\infty)\). Since \(u\in C_0(\C)\cap L^1(\C)\),
\(u^*\in L^1([0,\infty))\) and is equimeasurable with \(u\). The bathtub
principle yields
\[
I_F(s)=\int_0^s u^*(\tau)\,d\tau,\qquad s\ge 0.
\]
The right-hand side is the integral of an
\(L^1_{\mathrm{loc}}([0,\infty))\) function over \([0,s]\) and is therefore
absolutely continuous, hence continuous, in \(s\).
\end{proof}

The next result is the profile decomposition adapted to the present setting,
in the spirit of P.-L.~Lions' concentration-compactness method. Although the
overall scheme is by now standard, we spell out the points used in the
sequel; the version stated here is a Fock-space variant of profile
decompositions used in time-frequency analysis, see e.g.\
\cite{Nicola-Romero-Trapasso, Marceca-Romero-Speckbacher, Fulsche-Hagger,
Samuelsen}.

\begin{remark}[Translation non-compactness on \(\mathcal F^2\)]\label{rmk:translation-noncompact}
Although the Weyl translations \(T_a\) act unitarily on \(\mathcal F^2(\C)\),
the family \(\{T_a\}_{a\in\C}\) is not relatively compact in the strong
operator topology: by Lemma~\ref{lem:weyl-orthogonality-global}, for every
nonzero \(F\in \mathcal F^2\) and every sequence \(\{a_k\}\subset \C\) with
\(|a_k|\to\infty\) one has \(T_{a_k}F\rightharpoonup 0\) weakly, while the
norm \(\|T_{a_k}F\|_{\mathcal F^2}=\|F\|_{\mathcal F^2}\) is preserved.
Consequently, a bounded maximizing sequence in \(\mathcal F^2\) need not
admit any strongly convergent subsequence, even after factoring out the
Weyl translation group. The profile decomposition stated below is the
device that quantifies this non-compactness: it identifies all directions
along which mass can escape, by extracting a (possibly infinite) family of
nontrivial profiles \(G^{(j)}\) translated by mutually divergent centers
\(a_k^{(j)}\). It is essential to stress that, since the
\(\mathcal F^2\)-norm is invariant under \(T_a\), the remainder
\(R_k^{(J)}\) does \emph{not} converge to zero in \(\mathcal F^2\): the
Pythagorean identity~(iii) below merely redistributes mass between the
profiles and the remainder, and pure \(\mathcal F^2\)-vanishing of the
remainder is incompatible with the unitarity of \(T_a\) when several
nontrivial profiles are present. The natural and sharp vanishing statement
for the remainder is therefore the dispersive estimate~(iv) below,
expressed through the density \(u_{R_k^{(J)}}\); this is exactly the
quantity that controls all the localization functionals \(I_F(s)\) used in
the sequel, which depend only on the density \(u_F\) and not on the phase
of \(F\). The exposition below extracts what is needed in the sequel; an
analogous decomposition with dispersive-style remainders in the Fock
setting is discussed in
\cite{Nicola-Romero-Trapasso,Marceca-Romero-Speckbacher,Fulsche-Hagger,Samuelsen}.
\end{remark}

\begin{proposition}[Profile decomposition in the Fock space]\label{prop:profile-global}
Let \(\{F_k\}_{k\ge 1}\subset \mathcal F^2(\C)\) be a bounded sequence. Then,
after passing to a subsequence, there exist profiles
\[
G^{(j)}\in \mathcal F^2(\C),
\qquad j\ge 1,
\]
and sequences of points
\[
a_k^{(j)}\in \C,
\qquad j\ge 1,\ k\ge 1,
\]
such that:
\begin{enumerate}
    \item[(i)] for each \(j\neq \ell\),
    \[
    |a_k^{(j)}-a_k^{(\ell)}|\to \infty
    \qquad (k\to \infty);
    \]
    \item[(ii)] for every \(J\ge 1\), if
    \[
    R_k^{(J)}
    :=
    F_k-\sum_{j=1}^J T_{a_k^{(j)}}G^{(j)},
    \]
    then
    \[
    T_{-a_k^{(j)}}R_k^{(J)}\to 0
    \qquad\text{locally uniformly on }\C
    \]
    for every \(1\le j\le J\);
    \item[(iii)] for every \(J\ge 1\),
    \[
    \|F_k\|_{\mathcal F^2}^2
    =
    \sum_{j=1}^J \|G^{(j)}\|_{\mathcal F^2}^2
    +\|R_k^{(J)}\|_{\mathcal F^2}^2
    +o_k(1);
    \]
    \item[(iv)]
    \[
    \lim_{J\to\infty}\left(\limsup_{k\to\infty}
    \|u_{R_k^{(J)}}\|_{L^\infty(\C)}\right)=0.
    \]
\end{enumerate}
\end{proposition}

\begin{proof}
We argue by the standard iterative extraction scheme, but spell out the points
used in the rest of the section. Throughout, we say that a function in
\(\mathcal F^2(\C)\) is \emph{nontrivial} if it is not identically zero.

\medskip
\noindent\emph{Extraction of the first profile.}
If
\[
\limsup_{k\to\infty}\|u_{F_k}\|_{L^\infty(\C)}=0,
\]
there is nothing to prove: take all profiles equal to zero, all centers
arbitrary (say zero), and observe that (i)--(iii) hold trivially while (iv)
holds by hypothesis. Otherwise, after passing to a subsequence, choose
\(a_k^{(1)}\in \C\) such that
\[
u_{F_k}(a_k^{(1)})\ge \tfrac12\|u_{F_k}\|_{L^\infty(\C)}.
\]
By \eqref{eq:I-translation-global} and the density-translation identity
\eqref{eq:u-translation-global}, the sequence \(T_{-a_k^{(1)}}F_k\) has the
same Fock norm and the corresponding density translated; relabelling, we may
assume \(a_k^{(1)}=0\). Since \(\{F_k\}\) is bounded in \(\mathcal F^2(\C)\),
the reproducing-kernel bound \(|F_k(z)|\le \|F_k\|_{\mathcal F^2}\,
e^{\pi |z|^2/2}\) (Lemma~\ref{lem:growth}(ii)) makes \(\{F_k\}\) a normal
family of entire functions on every compact subset of \(\C\). Passing to a
subsequence and applying Montel's theorem, \(F_k\to G^{(1)}\) locally
uniformly on \(\C\), with \(G^{(1)}\in\mathcal F^2(\C)\) by Fatou's lemma. By
the choice of \(a_k^{(1)}=0\),
\[
u_{G^{(1)}}(0)=\lim_{k\to\infty}u_{F_k}(0)
\ge \tfrac12 \limsup_{k\to\infty}\|u_{F_k}\|_{L^\infty(\C)}>0,
\]
so \(G^{(1)}\) is nontrivial. Set
\[
R_k^{(1)}:=F_k-G^{(1)}.
\]

\medskip
\noindent\emph{Iteration.}
Suppose profiles \(G^{(1)},\dots,G^{(J-1)}\) and centers
\(a_k^{(1)},\dots,a_k^{(J-1)}\) have been extracted, with remainder
\(R_k^{(J-1)}\). If
\[
\limsup_{k\to\infty}\|u_{R_k^{(J-1)}}\|_{L^\infty(\C)}=0,
\]
we stop, set \(G^{(j)}=0\) and \(a_k^{(j)}=0\) for all \(j\ge J\), and
properties (i)--(iv) are easily seen to hold. Otherwise choose
\(a_k^{(J)}\in\C\) so that
\[
u_{R_k^{(J-1)}}(a_k^{(J)})
\ge \tfrac12\|u_{R_k^{(J-1)}}\|_{L^\infty(\C)}.
\]
Property (i) at this stage (proved below) forces
\(|a_k^{(J)}-a_k^{(m)}|\to\infty\) for every \(m<J\). Granted this, the
sequence \(\{T_{-a_k^{(J)}}R_k^{(J-1)}\}\) is bounded in \(\mathcal F^2\) (by
unitarity of \(T_{-a_k^{(J)}}\) together with (iii) at stage \(J-1\)), so
after passing to a further subsequence and applying Montel's theorem, it
converges locally uniformly to some \(G^{(J)}\in\mathcal F^2(\C)\). By the
choice of \(a_k^{(J)}\),
\[
u_{G^{(J)}}(0)
\ge \tfrac12 \limsup_{k\to\infty}\|u_{R_k^{(J-1)}}\|_{L^\infty(\C)}>0,
\]
so \(G^{(J)}\) is nontrivial. Set
\[
R_k^{(J)}:=R_k^{(J-1)}-T_{a_k^{(J)}}G^{(J)}.
\]
Iterating either terminates at some finite step (in which case (iv) is
immediate) or produces an infinite sequence of nontrivial profiles. By a
standard diagonal extraction we obtain a single subsequence (still denoted
\(\{F_k\}\)) for which all the partial decompositions are valid.

\medskip
\noindent\emph{Proof of (i).}
Fix \(j<\ell\), and suppose for contradiction that
\(|a_k^{(\ell)}-a_k^{(j)}|\) does not tend to infinity. After passing to a
subsequence, \(|a_k^{(\ell)}-a_k^{(j)}|\le C\) for some constant \(C\). At
the \(\ell\)-th step, by construction
\[
T_{-a_k^{(\ell)}}R_k^{(\ell-1)}\to G^{(\ell)}
\qquad\text{locally uniformly on }\C,
\]
with \(G^{(\ell)}\) nontrivial. On the other hand, property (ii) verified up
to stage \(\ell-1\) (proved below) gives
\[
T_{-a_k^{(j)}}R_k^{(\ell-1)}\to 0
\qquad\text{locally uniformly on }\C.
\]
The Weyl identity
\(T_{-a_k^{(\ell)}}=e^{i\theta_k}T_{-a_k^{(j)}}T_{a_k^{(j)}-a_k^{(\ell)}}\)
(for an appropriate real phase \(\theta_k\) coming from the Weyl co-cycle, of
modulus one) shows that
\[
T_{-a_k^{(\ell)}}R_k^{(\ell-1)}(z)
=
e^{i\theta_k}\big(T_{-a_k^{(j)}}R_k^{(\ell-1)}\big)\!
\big(z-(a_k^{(\ell)}-a_k^{(j)})\big).
\]
Since \(|a_k^{(\ell)}-a_k^{(j)}|\le C\), the points
\(z-(a_k^{(\ell)}-a_k^{(j)})\) remain in a fixed compact set as \(z\)
varies over a compact set; so the local uniform convergence of
\(T_{-a_k^{(j)}}R_k^{(\ell-1)}\) to \(0\) forces
\(T_{-a_k^{(\ell)}}R_k^{(\ell-1)}\to 0\) locally uniformly as well.
This contradicts the nontriviality of \(G^{(\ell)}\). Hence
\[
|a_k^{(j)}-a_k^{(\ell)}|\to\infty
\qquad\text{whenever }j\neq \ell.
\]

\medskip
\noindent\emph{Proof of (ii).}
Fix \(J\ge 1\) and \(1\le j\le J\). By construction, at the \(j\)-th
extraction step,
\[
T_{-a_k^{(j)}}R_k^{(j-1)}\to G^{(j)}
\qquad\text{locally uniformly on }\C.
\]
Subtracting the \(j\)-th extracted profile,
\[
T_{-a_k^{(j)}}R_k^{(j)}
=
T_{-a_k^{(j)}}R_k^{(j-1)}-G^{(j)}
\to 0
\qquad\text{locally uniformly on }\C.
\]
For \(J>j\), since
\[
R_k^{(J)}
=
R_k^{(j)}-\sum_{\ell=j+1}^J T_{a_k^{(\ell)}}G^{(\ell)},
\]
the linearity of the Weyl operators yields
\[
T_{-a_k^{(j)}}R_k^{(J)}
=
T_{-a_k^{(j)}}R_k^{(j)}
-\sum_{\ell=j+1}^J e^{i\theta_{k,\ell}}\,T_{a_k^{(\ell)}-a_k^{(j)}}G^{(\ell)}
\]
for some real phases \(\theta_{k,\ell}\) coming from the Weyl co-cycle. The
first term tends to \(0\) locally uniformly by the previous paragraph. For
each fixed \(\ell>j\), part (i) gives \(|a_k^{(\ell)}-a_k^{(j)}|\to\infty\),
and the identity \(u_{T_b G}(z)=u_G(z-b)\) (i.e.\
\eqref{eq:u-translation-global}) implies that
\(T_{a_k^{(\ell)}-a_k^{(j)}}G^{(\ell)}\to 0\) locally uniformly in modulus
(its density is \(u_{G^{(\ell)}}(z-(a_k^{(\ell)}-a_k^{(j)}))\to 0\) for
\(z\) in a compact set, since \(u_{G^{(\ell)}}\in C_0(\C)\)). Thus
\[
T_{-a_k^{(j)}}R_k^{(J)}\to 0
\qquad\text{locally uniformly on }\C,
\]
which is exactly (ii).

\medskip
\noindent\emph{Proof of (iii).}
For fixed \(J\ge 1\),
\[
F_k=\sum_{j=1}^J T_{a_k^{(j)}}G^{(j)}+R_k^{(J)}.
\]
Expanding the squared \(\mathcal F^2\)-norm and using
Lemma~\ref{lem:weyl-orthogonality-global} together with the pairwise
divergence of the translation parameters from (i),
\[
\Big\|\sum_{j=1}^J T_{a_k^{(j)}}G^{(j)}\Big\|_{\mathcal F^2}^2
=
\sum_{j=1}^J \|G^{(j)}\|_{\mathcal F^2}^2+o_k(1).
\]
Moreover, since
\[
T_{-a_k^{(m)}}R_k^{(J)}\to 0
\qquad\text{locally uniformly on }\C
\]
for every \(1\le m\le J\) by (ii), and the sequence
\(\{T_{-a_k^{(m)}}R_k^{(J)}\}_k\) is bounded in \(\mathcal F^2\), the same
weak-convergence argument as in
Lemma~\ref{lem:weyl-orthogonality-global} yields
\[
\langle T_{a_k^{(m)}}G^{(m)},R_k^{(J)}\rangle_{\mathcal F^2}
=
\langle G^{(m)},T_{-a_k^{(m)}}R_k^{(J)}\rangle_{\mathcal F^2}\to 0,
\]
for every fixed \(m\). Expanding \(\|F_k\|_{\mathcal F^2}^2\) and collecting
terms,
\[
\|F_k\|_{\mathcal F^2}^2
=
\sum_{j=1}^J \|G^{(j)}\|_{\mathcal F^2}^2
+\|R_k^{(J)}\|_{\mathcal F^2}^2
+o_k(1).
\]
In particular, since \(\|F_k\|_{\mathcal F^2}\) is bounded and the terms on
the right are nonnegative, summing the profile-norm contributions over \(j\)
yields the energy bound
\begin{equation}\label{eq:energy-bound-profiles}
\sum_{j\ge 1}\|G^{(j)}\|_{\mathcal F^2}^2
\le \limsup_{k\to\infty}\|F_k\|_{\mathcal F^2}^2.
\end{equation}
This is the key compactness ingredient that controls the iteration and
ensures (iv).

\medskip
\noindent\emph{Proof of (iv).}
Set
\[
\eta_J:=\limsup_{k\to\infty}\|u_{R_k^{(J)}}\|_{L^\infty(\C)},
\]
which is monotone nonincreasing in \(J\) by construction (since
\(R_k^{(J+1)}\) is obtained from \(R_k^{(J)}\) by extracting one further
nontrivial profile). Let \(\eta_\infty:=\lim_{J\to\infty}\eta_J\); we must
show \(\eta_\infty=0\). Suppose, for contradiction, that
\(\eta_\infty\ge \delta>0\).

By the iteration, at each stage \(J\) such that \(\eta_J\ge \delta\), the
extracted profile \(G^{(J+1)}\) satisfies \(u_{G^{(J+1)}}(0)\ge \delta/2\) by
construction. By Lemma~\ref{lem:Fock-basic}(ii),
\[
\|G^{(J+1)}\|_{\mathcal F^2}^2 \ge u_{G^{(J+1)}}(0)\ge \delta/2.
\]
If the iteration produced infinitely many such stages (which is the case
under the assumption \(\eta_\infty\ge \delta\)), then
\[
\sum_{j\ge 1}\|G^{(j)}\|_{\mathcal F^2}^2=\infty,
\]
contradicting \eqref{eq:energy-bound-profiles}. Hence \(\eta_\infty=0\),
proving (iv).
\end{proof}

We now turn to the global problem. The proof of existence consists of
combining the profile decomposition with the elementary continuity properties
of \(I_F(s)\), and using the homogeneity \(I_{cF}(s)=|c|^2 I_F(s)\) recorded
in Lemma~\ref{lem:basic-properties-global}(ii) to redistribute mass to a
single nontrivial profile.

\begin{theorem}\label{thm:existence-global-maximizer}
For every \(s>0\), there exists \(F_*\in \mathcal F^2(\C)\) with
\(\|F_*\|_{\mathcal F^2}=1\) such that
\[
I_{F_*}(s)=M(s).
\]
Equivalently, there exists a measurable set \(\Omega_*\subset \C\) with
\(|\Omega_*|=s\) such that
\[
\lambda_{1,\varphi}(\Omega_*)=M(s).
\]
\end{theorem}

\begin{proof}
Let \(\{F_k\}\subset \mathcal F^2(\C)\) be a maximizing sequence:
\[
\|F_k\|_{\mathcal F^2}=1,
\qquad
I_{F_k}(s)\to M(s).
\]
By replacing \(F_k\) with a Weyl translate if necessary, and using the
translation invariance \eqref{eq:I-translation-global} together with
\eqref{eq:u-translation-global}, we may suppose that
\[
u_{F_k}(0)=\|u_{F_k}\|_{L^\infty(\C)}.
\]
Note that \(M(s)>0\): indeed, taking \(F\equiv 1\), one has
\(F\in \mathcal F^2(\C)\) with \(\|F\|_{\mathcal F^2}^2=1\) and density
\(u_F(z)=e^{-\pi |z|^2}>0\) everywhere, so for any \(s>0\), \(I_F(s)>0\).
Consequently, \(\|u_{F_k}\|_{L^\infty(\C)}\) is bounded below by a positive
constant uniformly in \(k\); otherwise, for any \(\eta>0\) we would have
\(\|u_{F_k}\|_{L^\infty}<\eta\) eventually, hence \(I_{F_k}(s)\le \eta s\),
contradicting \(I_{F_k}(s)\to M(s)>0\).

Applying Proposition~\ref{prop:profile-global}, after passing to a
subsequence we may write
\[
F_k=\sum_{j=1}^J T_{a_k^{(j)}}G^{(j)}+R_k^{(J)}
\]
for every fixed \(J\ge 1\), with the conclusions of that proposition.

\medskip
\noindent\emph{Step 1: an upper bound for the limit.}
Fix \(J\ge 1\) and let \(\Omega_k^{(J)}\) be a measurable set of measure
\(s\) such that
\[
I_{F_k}(s)=\int_{\Omega_k^{(J)}} u_{F_k}\,dA;
\]
by Lemma~\ref{lem:basic-properties-global}(i) we may take
\(\Omega_k^{(J)}\) to be a superlevel set of \(u_{F_k}\). Let
\(\varepsilon>0\). Since \(u_{G^{(j)}}\in L^1(\C)\) for every \(j\) by
Lemma~\ref{lem:basic-properties-global}, choose \(R>0\) so large that
\[
\int_{\C\setminus D_R} u_{G^{(j)}}\,dA<\varepsilon
\qquad\text{for all } 1\le j\le J,
\]
and large enough that the conclusion of
Lemma~\ref{lem:pointwise-decoupling-global}(ii) holds with this choice of
\(R\). Since the translation parameters are pairwise divergent by
Proposition~\ref{prop:profile-global}(i), for \(k\) large enough the disks
\[
D_R(a_k^{(1)}),\dots,D_R(a_k^{(J)})
\]
are pairwise disjoint. Write
\[
B_{k,j}:=D_R(a_k^{(j)}),
\qquad
E_{k,j}:=\Omega_k^{(J)}\cap B_{k,j},
\qquad
E_{k,0}:=\Omega_k^{(J)}\setminus \bigcup_{j=1}^J B_{k,j}.
\]
Set
\[
s_{k,j}:=|E_{k,j}|,
\qquad 0\le j\le J.
\]
Then, since the \(\{B_{k,j}\}_{j=1}^J\) are disjoint and
\(\Omega_k^{(J)}\) has measure \(s\),
\[
\sum_{j=0}^J s_{k,j}=s.
\]

Set
\[
S_k^{(J)}:=\sum_{j=1}^J T_{a_k^{(j)}}G^{(j)}.
\]
By Proposition~\ref{prop:profile-global}(iii), the sequence
\(\{S_k^{(J)}\}_k\) is bounded in \(\mathcal F^2(\C)\), with bound
\(\le 1+o_k(1)\). By Proposition~\ref{prop:profile-global}(iv), if \(J\) is
taken sufficiently large then
\[
\limsup_{k\to\infty}\|u_{R_k^{(J)}}\|_{L^\infty(\C)}<\varepsilon^2.
\]
Lemma~\ref{lem:density-perturbation-global} therefore gives
\[
\|u_{F_k}-u_{S_k^{(J)}}\|_{L^\infty(\C)}\to 0
\qquad (k\to\infty),
\]
and in particular, after taking \(k\) large enough,
\begin{equation}\label{eq:density-approximation-global}
\|u_{F_k}-u_{S_k^{(J)}}\|_{L^\infty(\C)}<\varepsilon.
\end{equation}

Now
\[
E_{k,0}
=
\Omega_k^{(J)}\setminus \bigcup_{j=1}^J B_{k,j}
\subset
\C\setminus \bigcup_{j=1}^J D_R(a_k^{(j)}).
\]
Lemma~\ref{lem:pointwise-decoupling-global}(i) gives, for each fixed
\(j\in\{1,\dots,J\}\),
\[
\sup_{z\in B_{k,j}}
\big|u_{S_k^{(J)}}(z)-u_{T_{a_k^{(j)}}G^{(j)}}(z)\big|\to 0
\qquad (k\to\infty),
\]
while Lemma~\ref{lem:pointwise-decoupling-global}(ii) yields, for our
choice of \(R\) and for \(k\) sufficiently large,
\[
\sup_{z\in E_{k,0}}u_{S_k^{(J)}}(z)\le \varepsilon.
\]
Combining these bounds with \eqref{eq:density-approximation-global}, we
conclude that for all sufficiently large \(k\),
\begin{enumerate}
\item[(a)] on each \(B_{k,j}\),
\[
u_{F_k}(z)\le u_{T_{a_k^{(j)}}G^{(j)}}(z)+2\varepsilon;
\]
\item[(b)] on \(E_{k,0}\),
\[
u_{F_k}(z)\le 2\varepsilon.
\]
\end{enumerate}

Therefore, decomposing the integral defining \(I_{F_k}(s)\) according to the
partition \(\Omega_k^{(J)}=E_{k,0}\sqcup\bigsqcup_{j=1}^J E_{k,j}\),
\begin{align*}
I_{F_k}(s)
&=
\int_{\Omega_k^{(J)}} u_{F_k}\,dA \\
&\le
\sum_{j=1}^J \int_{E_{k,j}} u_{T_{a_k^{(j)}}G^{(j)}}\,dA
+2\varepsilon \sum_{j=0}^J |E_{k,j}| \\
&=
\sum_{j=1}^J \int_{E_{k,j}-a_k^{(j)}} u_{G^{(j)}}\,dA
+2\varepsilon s \\
&\le
\sum_{j=1}^J I_{G^{(j)}}(s_{k,j})
+2\varepsilon s,
\end{align*}
where in the next-to-last step we used the change of variables
\(z\mapsto z-a_k^{(j)}\) together with \eqref{eq:u-translation-global}, and
in the last step the definition of \(I_{G^{(j)}}(s_{k,j})\) as the supremum
of \(\int_E u_{G^{(j)}}\,dA\) over sets of measure
\(s_{k,j}=|E_{k,j}|=|E_{k,j}-a_k^{(j)}|\).

Passing to a further subsequence if necessary, we may suppose
\[
s_{k,j}\to \sigma_j\in [0,s]
\qquad (k\to\infty)
\]
for each \(1\le j\le J\). By the continuity of \(I_{G^{(j)}}\) in the
measure parameter, given by Lemma~\ref{lem:I-continuity-global},
\[
M(s)=\lim_{k\to\infty}I_{F_k}(s)
\le
\sum_{j=1}^J I_{G^{(j)}}(\sigma_j)+2\varepsilon s.
\]
Since \(\varepsilon>0\) was arbitrary,
\begin{equation}\label{eq:M-upper-profiles-global}
M(s)\le \sum_{j=1}^J I_{G^{(j)}}(\sigma_j).
\end{equation}

\medskip
\noindent\emph{Step 2: redistribution of mass and saturation of the bound.}
We now use Lemma~\ref{lem:basic-properties-global}(ii). For each \(j\) with
\(G^{(j)}\neq 0\), set
\[
\widehat G^{(j)}:=\frac{G^{(j)}}{\|G^{(j)}\|_{\mathcal F^2}},
\qquad
\|\widehat G^{(j)}\|_{\mathcal F^2}=1,
\]
and use the homogeneity \(I_{cF}(\sigma)=|c|^2 I_F(\sigma)\) (with
\(c=\|G^{(j)}\|_{\mathcal F^2}\)) to obtain
\[
I_{G^{(j)}}(\sigma_j)
=
\|G^{(j)}\|_{\mathcal F^2}^2
\,I_{\widehat G^{(j)}}(\sigma_j).
\]
Since \(\widehat G^{(j)}\) has unit \(\mathcal F^2\)-norm, by definition of
\(M\) and the monotonicity of \(I_{\widehat G^{(j)}}\) given by
Lemma~\ref{lem:I-continuity-global} (using \(\sigma_j\le s\)),
\[
I_{\widehat G^{(j)}}(\sigma_j)\le I_{\widehat G^{(j)}}(s)\le M(s).
\]
Hence
\[
I_{G^{(j)}}(\sigma_j)\le \|G^{(j)}\|_{\mathcal F^2}^2\, M(s).
\]
For \(j\) with \(G^{(j)}=0\) the inequality is trivial. Summing over
\(1\le j\le J\) and using Proposition~\ref{prop:profile-global}(iii)
together with the normalization \(\|F_k\|_{\mathcal F^2}^2=1\),
\[
\sum_{j=1}^J I_{G^{(j)}}(\sigma_j)
\le
M(s)\sum_{j=1}^J \|G^{(j)}\|_{\mathcal F^2}^2
\le M(s).
\]
Combined with \eqref{eq:M-upper-profiles-global}, this yields
\[
M(s)\le \sum_{j=1}^J I_{G^{(j)}}(\sigma_j)\le M(s)
\]
for every \(J\). Therefore equality holds throughout for every \(J\), which
forces
\[
\sum_{j\ge 1}\|G^{(j)}\|_{\mathcal F^2}^2=1
\]
(so the residual \(R_k^{(J)}\) carries no \(\mathcal F^2\)-mass in the
limit), and moreover, for every \(j\) with \(G^{(j)}\neq 0\),
\[
I_{\widehat G^{(j)}}(\sigma_j)=M(s).
\]
By monotonicity of \(I_{\widehat G^{(j)}}\),
\(I_{\widehat G^{(j)}}(s)\ge I_{\widehat G^{(j)}}(\sigma_j)=M(s)\); since
\(M(s)\) is the supremum over unit-norm functions, we conclude
\[
I_{\widehat G^{(j)}}(s)=M(s).
\]
Thus every nonzero normalized profile is a global maximizer.

\medskip
\noindent\emph{Step 3: producing a maximizing set.}
Choose any \(j_0\) with \(G^{(j_0)}\neq 0\) (such \(j_0\) exists since
\(\sum_{j\ge 1}\|G^{(j)}\|_{\mathcal F^2}^2=1\)) and set
\(F_*:=\widehat G^{(j_0)}\), so that
\[
\|F_*\|_{\mathcal F^2}=1,
\qquad
I_{F_*}(s)=M(s).
\]
By Lemma~\ref{lem:basic-properties-global}(i), there exists a superlevel
set \(\Omega_*\) of \(u_{F_*}\) with \(|\Omega_*|=s\) and
\[
\int_{\Omega_*}u_{F_*}\,dA=I_{F_*}(s)=M(s).
\]
Then
\[
\lambda_{1,\varphi}(\Omega_*)\ge \int_{\Omega_*}u_{F_*}\,dA=M(s),
\]
while by definition of \(M(s)\) we have
\(\lambda_{1,\varphi}(\Omega_*)\le M(s)\). Hence
\[
\lambda_{1,\varphi}(\Omega_*)=M(s),
\]
and \(\Omega_*\) is a global maximizer of measure \(s\).
\end{proof}

We may now identify the global maximizers, by combining the existence theorem
above with the local-rigidity theorem of the previous section.

\begin{theorem}\label{thm:global-maximizers-are-disks}
For every \(s>0\), every global maximizer for
\[
\sup_{|\Omega|=s}\lambda_{1,\varphi}(\Omega)
\]
is, up to translation, a disk.
\end{theorem}

\begin{proof}
Let \(\Omega_*\) be a global maximizer of measure \(s\) (which exists by
Theorem~\ref{thm:existence-global-maximizer}). Then \(\Omega_*\) is, in
particular, a local maximizer of \(\lambda_{1,\varphi}\) under the same
measure constraint: every sufficiently small smooth deformation of
\(\Omega_*\) preserving the measure \(|\Omega|=s\) yields a competitor
whose value of \(\lambda_{1,\varphi}\) is at most
\(M(s)=\lambda_{1,\varphi}(\Omega_*)\). The local-rigidity theorem proved in
the previous section therefore applies: every such local maximizer is, up
to translation, a disk. Since the problem is invariant under Weyl
translations, this characterization is sharp.
\end{proof}

As an immediate consequence, we recover the classical Nicola--Tilli theorem
\cite{Nicola-Tilli, Nicola-Tilli-2}.

\begin{corollary}
Among all measurable sets of prescribed measure, the disk maximizes Gaussian
time-frequency concentration:
\[
\lambda_{1,\varphi}(\Omega)\le \lambda_{1,\varphi}(D)
\]
whenever \(|\Omega|=|D|\) and \(D\) is a disk.
\end{corollary}

\begin{proof}
Existence of a global maximizer is given by
Theorem~\ref{thm:existence-global-maximizer}, and
Theorem~\ref{thm:global-maximizers-are-disks} shows that every such
maximizer is a disk, up to translation. Since translation by an element of
\(\C\) preserves both Lebesgue measure and the value of
\(\lambda_{1,\varphi}\), the claim follows.
\end{proof}


\bibliography{biblio}
\bibliographystyle{amsplain}

\end{document}
